% ============================================================
%  SY02-20《几何 Satake、范畴迹与志村簇的算术》中文译注本·主文件
%  底本：Xinwen Zhu, arXiv:1810.07375（50 页）
% ============================================================
\documentclass[UTF8, a4paper, 12pt, fontset=fandol, scheme=chinese]{ctexart}
% ============================================================
%  SY02-20《几何 Satake、范畴迹与志村簇的算术》中文译注本·导言区
%  底本：Xinwen Zhu, Geometric Satake, categorical traces, and
%        arithmetic of Shimura varieties, arXiv:1810.07375（50 页）
%  体例：定理类挂 [thm] 计数器、由 subsection 重置；公式 \numberwithin{equation}{subsection}
%        复现原文「Example 2.1.1.」与「(2.1.1)」——两套计数器互相独立
% ============================================================

% ---------- 页面 ----------
\usepackage[a4paper, left=2.6cm, right=2.6cm, top=2.8cm, bottom=2.6cm]{geometry}
\linespread{1.22}
\setlength{\parskip}{0.25em}
\setlength{\emergencystretch}{1.5em}
\sloppy

% ---------- 数学 ----------
\usepackage{amsmath, amsthm, amssymb, amsfonts}
\usepackage{mathtools}
\usepackage{mathrsfs}
\usepackage{stmaryrd}
\usepackage{bm}
\usepackage{amscd}
\allowdisplaybreaks[2]

% ---------- 图 ----------
\usepackage{graphicx}
\usepackage{xcolor}
\usepackage{tikz}
\usetikzlibrary{cd, arrows.meta, positioning, calc}
\usepackage[most]{tcolorbox}
\usepackage{enumitem}
\setlist{itemsep=2pt, topsep=3pt, parsep=0pt}
\setlist[enumerate,1]{label=(\arabic*)}

% ---------- 超链接 ----------
\usepackage[hyperfootnotes=false, colorlinks=true, linkcolor=blue!45!black,
            citecolor=green!40!black, urlcolor=teal!60!black, bookmarksnumbered=true]{hyperref}
\usepackage{bookmark}

% ---------- 中文微调 ----------
\xeCJKsetup{CJKecglue={\hskip 0.24em plus 0.05em minus 0.05em}}
\ctexset{
  section    = {name={,}, aftername={.\enspace}, format=\Large\bfseries\centering},
  subsection = {aftername={.\enspace}, format=\large\bfseries},
  subsubsection = {aftername={.\enspace}, format=\normalsize\bfseries}
}

% ---------- 定理环境 ----------
% 原文：\newtheorem{thm}{Theorem}[subsection]，其余定理类挂 [thm]；
% 故定理印「定理 2.1.1.」（名在前、三段号），公式印 (2.1.1)，两者独立计数。
\theoremstyle{plain}
\newtheorem{thm}{定理}[subsection]
\newtheorem{prop}[thm]{命题}
\newtheorem{lem}[thm]{引理}
\newtheorem{sublem}[thm]{子引理}
\newtheorem{lem-def}[thm]{引理-定义}
\newtheorem{cor}[thm]{推论}
\newtheorem{theorem}[thm]{定理}
\newtheorem{proposition}[thm]{命题}
\newtheorem{lemma}[thm]{引理}
\newtheorem{sublemma}[thm]{子引理}
\newtheorem{conjecture}[thm]{猜想}
\newtheorem{corollary}[thm]{推论}
\newtheorem{conj}[thm]{猜想}
\theoremstyle{definition}
\newtheorem{example}[thm]{例}
\newtheorem{remark}[thm]{注}
\newtheorem{hypothesis}[thm]{假设}
\newtheorem{definition}[thm]{定义}
\newtheorem{question}[thm]{问题}
\newtheorem{ex}[thm]{例}
\newtheorem{exs}[thm]{练习}
\newtheorem{notation}[thm]{记号}
\newtheorem{construction}[thm]{构造}
\newtheorem{caution}[thm]{注意}
\newtheorem{convention}[thm]{约定}
\newtheorem{rmk}[thm]{注}
\newtheorem{dfn}[thm]{定义}
\numberwithin{equation}{subsection}

\renewcommand{\proofname}{证明}
\renewcommand{\qedsymbol}{$\square$}
\renewcommand{\refname}{参考文献}
\renewcommand{\contentsname}{目\quad 录}
\renewcommand{\abstractname}{摘\quad 要}

% ---------- 译注（独立脚注序列，不占作者脚注编号） ----------
\usepackage{manyfoot}
\DeclareNewFootnote[plain]{trans}
\renewcommand{\thefootnotetrans}{译注\arabic{footnotetrans}}
\newcommand{\zy}[1]{\footnotetrans{\hspace{0.25em}#1}}
\newtcolorbox{yizhu}[1][]{
  enhanced, breakable, colback=orange!4!white, colframe=orange!55!black,
  boxrule=0.5pt, left=6pt, right=6pt, top=4pt, bottom=4pt,
  fonttitle=\bfseries\small, title=#1,
  before skip=6pt, after skip=6pt
}

% ---------- 原文宏 ----------
\newcommand{\quash}[1]{}   % 原文第 1185–1190 行的注即被此宏吞掉（PDF 中不出现）
\newcommand{\nc}{\newcommand}   % 原文用 \nc 作 \newcommand 的缩写
% 由 工程文件/scripts/port_macros2.py 从原文导言区自动移植，勿手改。
\newcommand{\ca}{\mathfrak{a}}
\newcommand{\cb}{\mathfrak{b}}
\newcommand{\cl}{\mathrm{cl}}
\DeclareMathOperator{\Int}{Int}
\newcommand{\disc}{\mathrm{disc}}
\newcommand{\lead}{\mathrm{lead}}
\DeclareMathOperator{\HC}{HC}
\DeclareMathOperator{\Ker}{Ker}
\newcommand{\bbb}{\mathbf{b}}
\newcommand{\bba}{\mathbf{a}}
\newcommand{\AAA}{\mathbb{A}}
\newcommand{\BB}{\mathbb{B}}
\newcommand{\CC}{\mathbb{C}}
\newcommand{\DD}{\mathbb{D}}
\newcommand{\EE}{\mathbb{E}}
\newcommand{\FF}{\mathbb{F}}
\newcommand{\GG}{\mathbb{G}}
\newcommand{\HH}{\mathbb{H}}
\newcommand{\II}{\mathbb{I}}
\newcommand{\JJ}{\mathbb{J}}
\newcommand{\KK}{\mathbb{K}}
\newcommand{\LL}{\mathbb{L}}
\newcommand{\MM}{\mathbb{M}}
\newcommand{\NN}{\mathbb{N}}
\newcommand{\OO}{\mathbb{O}}
\newcommand{\PP}{\mathbb{P}}
\newcommand{\QQ}{\mathbb{Q}}
\newcommand{\RR}{\mathbb{R}}
\newcommand{\SSS}{\mathbb{S}}
\newcommand{\TT}{\mathbb{T}}
\newcommand{\UU}{\mathbb{U}}
\newcommand{\VV}{\mathbb{V}}
\newcommand{\WW}{\mathbb{W}}
\newcommand{\XX}{\mathbb{X}}
\newcommand{\YY}{\mathbb{Y}}
\newcommand{\ZZ}{\mathbb{Z}}
\newcommand{\calA}{\mathcal{A}}
\newcommand{\calB}{\mathcal{B}}
\newcommand{\calC}{\mathcal{C}}
\newcommand{\calD}{\mathcal{D}}
\newcommand{\calE}{\mathcal{E}}
\newcommand{\calF}{\mathcal{F}}
\newcommand{\calG}{\mathcal{G}}
\newcommand{\calH}{\mathcal{H}}
\newcommand{\calI}{\mathcal{I}}
\newcommand{\calJ}{\mathcal{J}}
\newcommand{\calK}{\mathcal{K}}
\newcommand{\calL}{\mathcal{L}}
\newcommand{\calM}{\mathcal{M}}
\newcommand{\calN}{\mathcal{N}}
\newcommand{\calO}{\mathcal{O}}
\newcommand{\calP}{\mathcal{P}}
\newcommand{\calQ}{\mathcal{Q}}
\newcommand{\calR}{\mathcal{R}}
\newcommand{\calS}{\mathcal{S}}
\newcommand{\calT}{\mathcal{T}}
\newcommand{\calU}{\mathcal{U}}
\newcommand{\calV}{\mathcal{V}}
\newcommand{\calW}{\mathcal{W}}
\newcommand{\calX}{\mathcal{X}}
\newcommand{\calY}{\mathcal{Y}}
\newcommand{\calZ}{\mathcal{Z}}
\newcommand{\out}{\mathrm{out}}
\newcommand{\fraka}{{\mathfrak a}}
\newcommand{\frakb}{{\mathfrak b}}
\newcommand{\frakc}{{\mathfrak c}}
\newcommand{\frakd}{{\mathfrak d}}
\newcommand{\frake}{{\mathfrak e}}
\newcommand{\frakf}{{\mathfrak f}}
\newcommand{\frakg}{{\mathfrak g}}
\newcommand{\frakh}{{\mathfrak h}}
\newcommand{\fraki}{{\mathfrak i}}
\newcommand{\frakj}{{\mathfrak j}}
\newcommand{\frakk}{{\mathfrak k}}
\newcommand{\frakl}{{\mathfrak l}}
\newcommand{\frakm}{{\mathfrak m}}
\newcommand{\frakn}{{\mathfrak n}}
\newcommand{\frako}{{\mathfrak o}}
\newcommand{\frakp}{{\mathfrak p}}
\newcommand{\frakq}{{\mathfrak q}}
\newcommand{\frakr}{{\mathfrak r}}
\newcommand{\fraks}{{\mathfrak s}}
\newcommand{\frakt}{{\mathfrak t}}
\newcommand{\fraku}{{\mathfrak u}}
\newcommand{\frakv}{{\mathfrak v}}
\newcommand{\frakw}{{\mathfrak w}}
\newcommand{\frakx}{{\mathfrak x}}
\newcommand{\fraky}{{\mathfrak y}}
\newcommand{\frakz}{{\mathfrak z}}
\newcommand{\scrC}{\mathscr{C}}
\newcommand{\scrS}{\mathscr{S}}
\newcommand{\frakA}{{\mathfrak A}}
\newcommand{\frakB}{{\mathfrak B}}
\newcommand{\frakC}{{\mathfrak C}}
\newcommand{\frakD}{{\mathfrak D}}
\newcommand{\frakE}{{\mathfrak E}}
\newcommand{\frakF}{{\mathfrak F}}
\newcommand{\frakG}{{\mathfrak G}}
\newcommand{\frakH}{{\mathfrak H}}
\newcommand{\frakI}{{\mathfrak I}}
\newcommand{\frakJ}{{\mathfrak J}}
\newcommand{\frakK}{{\mathfrak K}}
\newcommand{\frakL}{{\mathfrak L}}
\newcommand{\frakM}{{\mathfrak M}}
\newcommand{\frakN}{{\mathfrak N}}
\newcommand{\frakO}{{\mathfrak O}}
\newcommand{\frakP}{{\mathfrak P}}
\newcommand{\frakQ}{{\mathfrak Q}}
\newcommand{\frakR}{{\mathfrak R}}
\newcommand{\frakS}{{\mathfrak S}}
\newcommand{\frakT}{{\mathfrak T}}
\newcommand{\frakU}{{\mathfrak U}}
\newcommand{\frakV}{{\mathfrak V}}
\newcommand{\frakW}{{\mathfrak W}}
\newcommand{\frakX}{{\mathfrak X}}
\newcommand{\frakY}{{\mathfrak Y}}
\newcommand{\frakZ}{{\mathfrak Z}}
\newcommand{\bA}{{\mathbb A}}
\newcommand{\bB}{{\mathbb B}}
\newcommand{\bC}{{\mathbb C}}
\newcommand{\bD}{{\mathbb D}}
\newcommand{\bE}{{\mathbb E}}
\newcommand{\bF}{{\mathbb F}}
\newcommand{\bG}{{\mathbb G}}
\newcommand{\bH}{{\mathbb H}}
\newcommand{\bI}{{\mathbb I}}
\newcommand{\bJ}{{\mathbb J}}
\newcommand{\bK}{{\mathbb K}}
\newcommand{\bL}{{\mathbb L}}
\newcommand{\bM}{{\mathbb M}}
\newcommand{\bN}{{\mathbb N}}
\newcommand{\bO}{{\mathbb O}}
\newcommand{\bP}{{\mathbb P}}
\newcommand{\bQ}{{\mathbb Q}}
\newcommand{\bR}{{\mathbb R}}
\newcommand{\bS}{{\mathbb S}}
\newcommand{\bT}{{\mathbb T}}
\newcommand{\bU}{{\mathbb U}}
\newcommand{\bV}{{\mathbb V}}
\newcommand{\bW}{{\mathbb W}}
\newcommand{\bX}{{\mathbb X}}
\newcommand{\bY}{{\mathbb Y}}
\newcommand{\bZ}{{\mathbb Z}}
\newcommand{\mA}{{\mathcal A}}
\newcommand{\mB}{{\mathcal B}}
\newcommand{\mC}{{\mathcal C}}
\newcommand{\mD}{{\mathcal D}}
\newcommand{\mE}{{\mathcal E}}
\newcommand{\mF}{{\mathcal F}}
\newcommand{\mG}{{\mathcal G}}
\newcommand{\mH}{{\mathcal H}}
\newcommand{\mI}{{\mathcal I}}
\newcommand{\mJ}{{\mathcal J}}
\newcommand{\mK}{{\mathcal K}}
\newcommand{\mL}{{\mathcal L}}
\newcommand{\mM}{{\mathcal M}}
\newcommand{\mN}{{\mathcal N}}
\newcommand{\mO}{{\mathcal O}}
\newcommand{\mP}{{\mathcal P}}
\newcommand{\mQ}{{\mathcal Q}}
\newcommand{\mR}{{\mathcal R}}
\newcommand{\mS}{{\mathcal S}}
\newcommand{\mT}{{\mathcal T}}
\newcommand{\mU}{{\mathcal U}}
\newcommand{\mV}{{\mathcal V}}
\newcommand{\mW}{{\mathcal W}}
\newcommand{\mX}{{\mathcal X}}
\newcommand{\mY}{{\mathcal Y}}
\newcommand{\mZ}{{\mathcal Z}}
\newcommand{\Tate}{\mathrm{Tate}}
\newcommand{\rmH}{\mathrm{H}}
\newcommand{\al}{{\alpha}} \nc{\be}{{\beta}} \nc{\ga}{{\gamma}}
\newcommand{\ve}{{\varepsilon}} \nc{\Ga}{{\Gamma}} \nc{\la}{{\lambda}}
\newcommand{\La}{{\Lambda}}
\newcommand{\ad}{{\on{ad}}}
\newcommand{\Ad}{{\on{Ad}}}
\newcommand{\aff}{{\on{aff}}}
\newcommand{\Aff}{{\mathbf{Aff}}}
\newcommand{\Aut}{{\on{Aut}}}
\newcommand{\Bun}{{\on{Bun}}}
\newcommand{\cha}{{\on{char}}}
\newcommand{\Coh}{{\on{Coh}}}
\newcommand{\der}{{\on{der}}}
\newcommand{\Der}{{\on{Der}}}
\newcommand{\diag}{{\on{diag}}}
\newcommand{\dR}{{\on{dR}}}
\newcommand{\End}{{\on{End}}}
\newcommand{\Fl}{{\calF\ell}}
\newcommand{\Gal}{{\on{Gal}}}
\newcommand{\Gr}{{\on{Gr}}}
\newcommand{\Hom}{{\on{Hom}}}
\newcommand{\id}{{\on{id}}}
\newcommand{\Id}{{\on{Id}}}
\newcommand{\ind}{{\on{ind}}}
\newcommand{\Ind}{{\on{Ind}}}
\newcommand{\inv}{{\on{Inv}}}
\newcommand{\Iso}{{\on{Isom}}}
\newcommand{\Lie}{{\on{Lie}}}
\newcommand{\Nm}{{\on{Nm}}}
\newcommand{\Pic}{{\on{Pic}}}
\newcommand{\pf}{{\on{pf}}}
\newcommand{\pr}{{\on{pr}}}
\newcommand{\rec}{{\on{rec}}}
\newcommand{\res}{{\on{res}}}
\newcommand{\Iw}{\on{Iw}}
\newcommand{\Res}{{\on{Res}}}
\newcommand{\s}{{\on{sc}}}
\newcommand{\sgn}{{\on{sgn}}}
\DeclareMathOperator{\Spec}{Spec}
\newcommand{\St}{{\on{St}}}
\newcommand{\tr}{{\on{tr}}}
\newcommand{\loc}{\mathrm{loc}}
\newcommand{\Tr}{{\on{Tr}}}
\newcommand{\Mod}{{\mathrm{-Mod}}}
\newcommand{\ord}{\mathrm{ord}}
\newcommand{\Orb}{\mathrm{Orb}}
\newcommand{\Shv}{\mathrm{Shv}}
\newcommand{\sym}{\mathrm{sym}}
\newcommand{\sep}{\mathrm{sep}}
\DeclareMathOperator{\RHom}{RHom}
\newcommand{\ev}{\on{ev}}
\newcommand{\coev}{\on{coev}}
\newcommand{\un}{\on{un}}
\newcommand{\coun}{\on{coun}}
\newcommand{\Frob}{\mathrm{Frob}}
\newcommand{\GL}{{\on{GL}}}
\newcommand{\GSp}{{\on{GSp}}} \nc{\GU}{{\on{GU}}} \nc{\SL}{{\on{SL}}}
\newcommand{\SU}{{\on{SU}}} \nc{\SO}{{\on{SO}}}
\newcommand{\gl}{{\frakg\frakl}}
\newcommand{\Ql}{{\overline{\bQ}_\ell}}
\newcommand{\reg}{\mathrm{reg}}
\newcommand{\proj}{\mathrm{proj}}
\newcommand{\fg}{\frakg}
\newcommand{\fp}{\frakp}
\newcommand{\pFr}{{_\mathrm{p}\sigma}}
\newcommand{\rat}{\overline{\bQ}}
\newcommand{\triv}{{\bf{1}}}
\newcommand{\trun}{\mathrm{trun}}
\newcommand{\longto}{\longrightarrow}
\newcommand{\bGr}{{\overline{\Gr}}}
\newcommand{\wGr}{{\widetilde{\Gr}}}
\newcommand{\std}{\mathrm{std}}
\newcommand{\perf}{\mathrm{perf}}
\newcommand{\et}{\mathrm{et}}
\newcommand{\ppart}{(\!(t)\!)}
\newcommand{\pparu}{(\!(u)\!)}
\newcommand{\dm}{/\!\!/}
\newcommand{\xcoch}{\mathbb{X}_\bullet}
\newcommand{\xch}{\mathbb{X}^\bullet}
\newcommand{\wt}{\mathrm{wt}}
\newcommand{\Sat}{\on{Sat}}
\newcommand{\Sh}{\on{Sh}}
\newcommand{\Sht}{\on{Sht}}
\newcommand{\wSh}{\widetilde{\Sht}}
\newcommand{\Sph}{\on{Sph}}
\newcommand{\Fr}{\on{Frob}}
\newcommand{\Fp}{{^\sigma\mE}}
\newcommand{\und}{\underline}
\newcommand{\mmu}{{\mu_\bullet}}
\newcommand{\nnu}{{\nu_\bullet}}
\newcommand{\Hk}{{\on{Hk}}}
\newcommand{\lhk}{\Hk^{\on{loc}}}
\newcommand{\IC}{\on{IC}}
\newcommand{\bb}{{\mathbf{b}}}
\newcommand{\bp}{{\mathbf{p}}}
\newcommand{\cris}{\mathrm{cris}}
\newcommand{\MV}{{\bM\bV}}
\newcommand{\Mp}{{\on{Mp}}}
\newcommand{\leftone}{\leftarrow}
\newcommand{\rightone}{\to}
\newcommand{\bfB}{{\mathbf{B}}}
\newcommand{\rmP}{\mathrm{P}}
\newcommand{\bbzeta}{\boldsymbol{\zeta}}
\newcommand{\on}{\operatorname}


% 原文 `\s` 与 `\sgn` 等单 token 名在少数运行环境下与内核冲突，这里统一兜底
\providecommand{\s}{{\operatorname{sc}}}

% 原文用 xy-pic 画交换图；必须在 tikz 之后加载，否则 \xymatrix 与 tikz-cd 冲突
\usepackage[all]{xy}


\begin{document}
\pagestyle{plain}

\begin{titlepage}
\centering
\vspace*{2.6cm}
{\Huge 几何 Satake、范畴迹\\[0.3em]与志村簇的算术\par}
\vspace{1.2cm}
{\LARGE Geometric Satake, categorical traces,\\[0.2em] and arithmetic of Shimura varieties\par}
\vspace{1.0cm}
{\large 中文译注本\par}
\vspace{1.8cm}
{\large 朱歆文（Xinwen Zhu）\par}
\vspace{0.6cm}
{\large 加州理工学院数学系\par}
\vspace{1.2cm}
{\large 原文：arXiv:1810.07375，50 页，2018 年\par}
\vspace{0.3cm}
{\large 会议报告：Current Developments in Mathematics 2016\par}
\vspace{0.3cm}
{\large 译者：数译（DSH agent）\par}
\vspace{0.3cm}
{\large 版本：v1.0　　日期：2026 年 11 月\par}
\vspace{1.4cm}
\rule{0.6\textwidth}{0.4pt}\par
\vspace{0.6cm}
{\small 本译本仅供学术研读；原文版权归作者所有。\par}
\vfill
\end{titlepage}

\section*{译者说明}
\addcontentsline{toc}{section}{译者说明}

\textbf{原文性质。}这是朱歆文在 2016 年 Current Developments in Mathematics 会议上的报告，
属\textbf{综述性质}：定义、主要定理与证明思路给全，技术细节常以「略去」或引用文献带过。
全文六节：引言的路线图；$p$ 进群的几何 Satake；范畴迹构造；几何 Satake 的范畴迹；
Shimura 簇算术的应用；关于对应范畴的附录。主体是几何 Langlands 的一支，
面向的读者是熟悉代数几何与表示论、但未必做过 Langlands 纲领的研究生与研究者。

\textbf{编号约定。}原文是 amsart 体例：定理类环境由 subsection 重置、共用一套计数器，
印刷为「Theorem 2.1.1.」（节.小节.序号，名字在前）；公式另用一套计数器，
印刷为 (2.1.1)。两套\textbf{互相独立}——同一小节里定理 2.1.2 之后可以紧接公式 (2.1.3)。
译文完全照此：定理类为「定理 2.1.1.」，公式为 (2.1.1)，与原文逐条对应；交叉引用一律经
\verb|\ref| 生成，不手写编号。

\textbf{译注标记。}文中「译$n$」是译者所加，用独立脚注序列，不占用作者脚注；
作者脚注保留原编号。每条解释性译注都标来源：有外部依据的写\verb|[来源：…]|，
仅凭上下文推断的写\verb|[译者按（上下文推断，无外部依据）：…]|。

\textbf{术语约定。}核心术语首次出现时括注英文，其后只用中文；Shtuka 保留原文不译，
Hecke 叠、几何 Satake 等价、反常层、半小映射、带钉、Tate 轮等译法与同系列的
SY02-17《仿射 Grassmann 簇与几何 Satake 等价》保持一致，详见
\verb|资料卡/04-术语文献勘误.md|。

\textbf{一处技术说明。}原文第 3.1 节末尾有一段关于 $\infty$-范畴中范畴迹的注，
在源文件中被作者用 \verb|\quash| 宏包住（该宏定义为空），因此\textbf{PDF 里并不出现}。
译文同样不译此段：若译出，其后所有定理编号会整体错一位，与原文不再对应。

\newpage
\tableofcontents
\newpage

% ============ 正文 ============

% ch01  §1 引言（原文 368–384 行）

\section{引言}

几何 Langlands 纲领由 Drinfeld 与 Laumon 发起，是 Drinfeld 关于整体函数域上 $\GL_2$ 的 Langlands 对应的工作的延续与推广。在几何理论中，基本研究对象从约化群 $G$ 的自守形式空间，变为代数曲线上 $G$-丛的模空间上的层范畴。

近年来，几何 Langlands 纲领对经典 Langlands 纲领以及若干相关的算术问题，都有富有成果的应用。传统做法是套用 Grothendieck 的层-函数字典，把几何理论中研究的层范畴「去范畴化」，得到算术理论中研究的函数空间。如上所述，Drinfeld 对 $\GL_2$ 的 Langlands 对应的处理用的就是这个办法（\cite{Dr}）。另一个著名的例子是 Ngô 对基本引理的证明（\cite{Ng}）。

在这篇综述中，我们要说明从几何理论通向算术理论的\emph{另一条}通道，它同样经由一个迹构造，但性质不同。这一方法的抽象版本在 \S\ref{S: cat trace} 讨论。应当指出，这一想法已经隐含在 V. Lafforgue 关于整体函数域上 Langlands 对应的工作（\cite{La}）中；Gaitsgory 的一篇非正式文章（\cite{Ga}）也有说明。

几何 Langlands 纲领的基石之一，是几何 Satake 等价：它在 $G$ 的仿射 Grassmann 簇上某类反常层的范畴，与 Langlands 对偶群 $\hat G$ 的有限维代数表示的范畴之间，建立张量等价。这是经典 Satake 同构（经由层-函数字典）的巨大推广，也可看作 Langlands 对偶群的一种概念性构造。

对等特征局部域上的约化群，几何 Satake 是 Lusztig、Ginzburg、Beilinson–Drinfeld 与 Mirković–Vilonen 等人工作的结果（\cite{Lu,Gi1,BD,MV}），现已在表示论、数学物理与（算术）代数几何中有许多应用。新近的工作 \cite{Zh1} 又对 $p$ 进群建立了这一等价，我们将在 \S\ref{S: GSat} 回顾。我们借机在 \S\ref{SS: MotSat} 讨论几何 Satake 的动机版本；它在某种意义上为本文的算术应用提供了一个玩具模型。接着在 \S\ref{S: cattrSat}，我们把上述抽象的迹构造用于 $p$ 进群的几何 Satake，由此得到对 Shimura 簇的上同调与轮的研究的一些应用，见 \S\ref{S:app to Sh}。最后两节基于与肖梁的合作工作 \cite{XZ}。

\medskip

\noindent\textbf{致谢。}本文源自作者在 Current Developments in Mathematics 2016 会议上的报告；作者感谢主办方的邀请。\S\ref{S: cattrSat} 与 \S\ref{S:app to Sh} 基于与肖梁（Liang Xiao）的合作工作，作者感谢他的合作。此外，作者感谢 Pavel Etingof、Dennis Gaitsgory 与刘一峰（Yifeng Liu）的有益讨论。

\section{$p$ 进群的几何 Satake}
\label{S: GSat}

本节回顾文献 \cite{Zh1} 中的工作。我们也借此机会，用动机的语言重新表述几何 Satake 等价。这让 \cite{Zh1} 中各项构造背后的某些想法更显透明，同时给出 \cite{XZ} 中算术应用的一个玩具模型；后者将在后面几节讨论。

\subsection{回顾 $p$ 进群的几何 Satake}
\label{S: affGrass}

我们先简要回顾经典 Satake 同构。

设 $F$ 为非阿基米德局部域，$\mO$ 为其整数环，$k=\bF_q$ 为其剩余域。换言之，$F$ 是 $\bQ_p$ 的有限扩张，或者同构于 $\bF_q((\varpi))$。
设 $\sigma$ 为 $k$ 的几何 $q$-Frobenius。

我们假定 $G$ 是 $\mO$ 上的连通约化群。
于是 $G(F)$ 是拓扑群，单位元的一组开邻域基由 $K:=G(\mO)$ 的有限指数子群给出。
例如取 $G=\GL_2$、$F=\bQ_p$、$\mO=\bZ_p$，令
$$\Ga(p^n)=\big\{\begin{pmatrix}a& b\\ c& d\end{pmatrix}\mid a=d=1 \mod p^n,\  b=c=0 \mod p^n\big\}$$
为第 $n$ 级主同余子群。则
$$\GL_2(\bZ_p)=\Ga(1)\supset\Ga(p)\supset\Ga(p^2)\supset\cdots$$ 构成 $\GL_2(\bQ_p)$ 中单位矩阵的一组开邻域基。

在此拓扑下 $G(F)$ 局部紧，$K$ 是开紧子群。因此存在唯一的 Haar 测度，使 $K$ 的体积为 $1$。
经典球面 Hecke 代数是 $G(F)$ 上紧支、$G(\mO)$-双不变、取值于 $\bC$ 的函数构成的空间，配备卷积积
\begin{equation}
\label{E:conv prod alge}
(f*g)(x)=\int_{G(F)}f(y)g(y^{-1}x)dy,
\end{equation}
注意若 $f$ 与 $g$ 都取值于 $\bZ$，则 $f*g$ 亦然。因此取值于 $\bZ$ 的函数构成一个 $\bZ$-子代数，记作 $H_G$。

由定义，双陪集 $\{KgK\mid g\in G\}$ 的特征函数构成 $H_G$ 的一组 $\bZ$-基。我们回忆这些双陪集的一个显式参数化。

设 $T\subset G$ 为定义在 $\mO$ 上的极大环面。令
$$\xcoch(T)=\Hom(\bG_m,T)$$ 表示 $T$ 的单参数子群构成的群，$\sigma$ 在其上有作用。
$\xcoch(T)$ 中的元素通常称为余特征标\zy{译者按（上下文推断，无外部依据）：原文用同一记号 $\xcoch$ 兼指余特征标群与余权群，本译本随上下文分别译作「余特征标」与「余权」，符号照旧。}。$\sigma$-不变的余特征标 $\la$ 定义在 $F$ 上，并诱导 $\la: F^\times\to T(F)\subset G(F)$；因此，取定一致化元 $\varpi\in F$，我们得到映射
\begin{equation*}\label{cartan1}
\xcoch(T)^\sigma\to G(F), \quad \la\mapsto \la(\varpi).
\end{equation*}
Cartan 分解断言，此映射诱导一个典范双射
\begin{equation}\label{cartan2}
\inv:  K\backslash G(F)/K\cong \xcoch(T)^{\sigma}/W_0 
\end{equation}
它与一致化元的选取无关。这里 $W=N_G(T)/T$ 是 $T$ 的 Weyl 群，$\sigma$ 在其上有作用，而 $W_0=W^{\sigma}$。

\begin{ex}\label{coweightGLn}
设 $G=\GL_n$ 定义在 $\bQ_p$ 上，可取 $T\subset G$ 为对角矩阵群。此时有典范等同 $\xcoch(T)\simeq \bZ^n$，$\sigma$ 在其上平凡作用。上述映射 $\xcoch(T)\to G(F)$ 由
\[\mu=(m_1,m_2,\ldots,m_n)\mapsto \mu(p)=\begin{pmatrix}p^{m_1}&&&\\&p^{m_2}&&\\ &&\ddots&\\ &&& p^{m_n}\end{pmatrix}.\]
给出。此情形下 Weyl 群是 $n$ 个文字的置换群 $S_n$，它通过置换直和项作用在 $\bZ^n$ 上。于是 $\xcoch(T)/W$ 中每个元素在
$$\xcoch(T)^+=\{\mu=(m_1,\ldots,m_n)\in \bZ^n\mid m_1\geq m_2\geq\cdots\geq m_n\}$$
中有唯一代表元。
\end{ex}

我们固定一个系数环 $E$，为简单起见假定它是特征零的代数闭域。
由 $G$ 可以构造 $E$ 上另一个分裂连通约化群 $\hat G$，称为 Langlands 对偶群。$\hat G$ 最初的定义是组合式的，依赖连通约化群的分类。一个更概念化的构造来自下文将叙述的几何 Satake（见注 \ref{R: birth of the dual group}），故此处不回顾 $\hat G$ 的原初定义。眼下只需知道：$\hat G$ 配有 Borel 子群 $\hat B$、极大环面 $\hat T\subset \hat B$，其特征标群 $\Hom(\hat T,\bG_m)$ 等于 $\xcoch(T)$；并且 $\hat G$ 的不可约表示的同构类由 $\xcoch(T)/W$ 参数化。此外，$\sigma$ 的作用诱导 $(\hat G,\hat B, \hat T)$ 的一个典范有限阶自同构。
下面列出一些应当记住的例子（忽略 $(\hat B,\hat T)$ 与 $\sigma$ 的作用）。

\medskip

\begin{center}
\begin{tabular}{| c | c | c | c | c | c |}
\hline 
$G$ & $\GL_n$ & $\on{SL}_n$ & $\on{SO}_{2n+1}$ & $\on{SO}_{2n}$ & $E_8$  \\ \hline 
$\hat G$ & $\GL_n$ & $\on{PGL}_{n}$ & $\on{Sp}_{2n}$ & $\on{SO}_{2n}$ & $E_8$ \\ \hline 
\end{tabular}
\end{center}

\medskip


现在考虑 $\hat G$ 在自身上由
\begin{equation}
\label{E:twist conj}
c_\sigma:\hat G\times \hat G\to\hat G, \quad (g,h)\mapsto c_\sigma(g)(h):=gh\sigma(g)^{-1},
\end{equation}
给出的 $\sigma$-扭共轭，它等价于 $\hat G$ 在旁集 $\hat G\sigma\subset \hat G\rtimes\langle\sigma\rangle=:{^L}G$ 上的通常共轭作用。令 $R({^L}G)$ 表示 ${^L}G$ 的有限维代数表示范畴 $\on{Rep}({^L}G)$ 的 Grothendieck $K$-环。
令 $\mathbf J= E[\hat G\sigma]^{\hat G}$ 表示 $\hat G$ 上 $\sigma$-扭共轭不变函数构成的空间。对 ${^L}G$ 的表示 $V$，令 $\chi_V$ 为 $V$ 的特征标在 $\hat G\sigma$ 上的限制。
则把 $[V]$ 送到 $\chi_V$ 的映射 $R({^L}G)\otimes E\to \mathbf J$ 是满射。

我们固定 $q^{1/2}\in E$ 的一个选取。
经典 Satake 同构（更确切地说，Langlands 的重新解释）给出典范同构
\begin{equation}\label{Sat isom}
\Sat^{cl}: H_G\otimes E\cong \mathbf J.
\end{equation}
它可以经 Grothendieck 的层-函数字典，由下文将讨论的几何 Satake 推出。故此处不讨论其证明（细节见 \cite[\S 5.6]{Zh2} 与 \cite[\S 3.5]{XZ}）。我们转而给出一个例子。

\begin{ex}
设 $G=\GL_2$。则 $\hat{G}=\GL_2$，$\sigma$ 在其上平凡作用，故 ${^L}G=\hat{G}$。作为代数，
$$H_{\GL_2}=\bZ[T_p,S_p^{\pm 1}],$$
其中 $T_p$ 是 $K\begin{pmatrix}p&\\&1\end{pmatrix}K$ 的特征函数，$S_p$ 是 $\begin{pmatrix}p&\\&p\end{pmatrix}K$ 的特征函数。
另一方面，已知
$$R(\GL_2)=\bZ[\tr,\det{}^{\pm 1}],$$
其中 $\tr$ 是通常的迹函数（标准表示 $\on{Std}$ 的特征标），$\det$ 是通常的行列式函数（$\wedge^2\on{Std}$ 的特征标）。于是，在选取 $p^{1/2}\in E$ 之后，Satake 同构由
$$\Sat^{cl}(T_p)=p^{\frac{1}{2}}\cdot \tr,\quad \Sat^{cl}(S_p)=\det.$$
给出。注意
\[T_p* T_p= (T_{p^2}+S_p)+pS_p,\quad \tr\cdot \tr= \chi_{\on{sym}^2\on{Std}}+\det,\]
其中 $T_{p^2}$ 是 $K\begin{pmatrix}p^2&\\&1\end{pmatrix}K$ 的特征函数，$\chi_{\on{sym}^2\on{Std}}$ 是标准表示的对称平方的迹函数。
由此推出
\begin{equation}\label{upper tr}
\Sat^{cl}(T_{p^2}+S_p)=p\cdot\chi_{\on{sym}^2\on{Std}}.
\end{equation}
\end{ex}

\begin{rmk}
注意由定义可知，\eqref{Sat isom} 中的两个环都有一组以 $\xcoch(T)^\sigma/W_0$ 为指标的自然基。即每个 $\mu\in \xcoch(T)^\sigma/W_0$ 给出特征函数 $1_{K\mu(\varpi)K}\in H_G$，以及函数 $\chi_{_{V_\mu}}\in \mathbf J$，其中 $V_\mu$ 是 ${^L}G$ 的「最高权 $\mu$」不可约表示。
然而，上例表明同构 $\Sat^{cl}$ 一般不把 $1_{K\mu(\varpi)K}$ 送到 $\chi_{_{V_\mu}}$。
\end{rmk}

\begin{rmk}
\label{R: norm Sat}
当然，可以定义同构 $\Sat'^{cl}: H_{\GL_2}\to R(\GL_2)$ 把 $T_p$ 送到 $\tr$（它在 $\bZ$ 上有定义，且与 $p^{1/2}$ 的选取无关）。这与 Satake 同构的归一化问题有关（见 \cite[\S 8]{Gr}）。更根本地，它与 Langlands 对偶群的\textbf{正确}定义有关。本文不讨论这一问题（注 \ref{R:Tate twist} 除外）。从几何 Satake 角度所作的详细讨论可参看 \cite[\S 5.5]{Zh2}。
\end{rmk}

\medskip

现在我们解释几何的做法：我们不再把 $G(F)$ 看作拓扑群、考虑其上 $K$-双不变的紧支函数空间，而是把 $G(F)$ 看作某个代数几何对象，研究其上的 $K$-双等变层范畴。等价而更方便的做法是，把 $G(F)/K$ 看作代数几何对象，研究其上的 $K$-等变层范畴。更确切地说，我们将构造 $k$ 上一个无穷维代数簇 $\Gr_G$，称为 $G$ 的仿射 Grassmann 簇，使 $G(F)/K$ 是它的 $k$-点集。
要看出这为何可行，我们先分析 $G(F)/K$ 的某些子集，它们可以自然地与某个代数簇的 $k$-点集等同。

\begin{ex}\label{minu Sch}
设 $G=\GL_n$ 定义在 $F$ 上。则商 $G(F)/K$ 可等同于 $F^n$ 中的格构成的集合。这里所谓格，指 $F^n$ 的有限自由 $\mO$-子模 $\Lambda$，且满足 $\Lambda\otimes_{\mO}F=F^n$。例如 $\La_0:=\mO^n\subset F^n$ 是格，通常称为标准格。双射由
\begin{equation}\label{coset vs lattice}
\GL_n(F)/K\simeq \big\{\mbox{格 } F^n\mbox{ 中的格}\big\},\quad gK\mapsto \Lambda:=g\La_0.
\end{equation}
给出。

考虑 $\omega_i=1^i0^{n-i}\in \bZ^{n}$，并如例 \ref{coweightGLn} 那样把它看作 $\GL_n$ 的余权。则容易看出，
在等同 \eqref{coset vs lattice} 下，集合 $K\omega_i(\varpi)K/K$ 可等同于满足
\[\big\{\varpi\mO^n\subset \Lambda \subset \mO^n,\quad  \dim_{k}\mO^n/\Lambda=i\big\}.\]
的格 $\Lambda\subset F^n$ 的集合。映射
$$\Lambda\mapsto (k^n=\mO^n/\varpi\mO^n\to \mO^n/\La)$$
把此集合等同于 $k^n$ 的 $i$ 维商空间（等价地，$k^n$ 的 $(n-i)$ 维子空间）的集合。总之，$K\omega_i(\varpi)K/K$ 与 Grassmann 簇 $\Gr(n-i,n)$（固定 $n$ 维空间中 $(n-i)$ 维平面的簇）的 $k$-点集之间有典范双射。自然期望 $\Gr(n-i,n)$ 是 $\Gr$ 的子簇。
\end{ex}

\begin{ex}\label{quasiminuSch}
设 $G=\GL_2$ 定义在 $F$ 上。依据 Satake 同构（特别见 \eqref{upper tr}），考虑集合
\begin{equation}\label{Gr2}
K\begin{pmatrix}\varpi^2 & \\ &1 \end{pmatrix}K/K\sqcup \begin{pmatrix}\varpi & \\ &\varpi \end{pmatrix} K/K.
\end{equation} 是合理的。在等同 \eqref{coset vs lattice} 下，它可等同于满足
\begin{equation}\label{Gr2latt}
\{\La\subset \mO^2\mid \on{length}\mO^2/\La=2\}.
\end{equation}
的格 $\La\subset F^2$ 的集合。
若 $F=k((\varpi))$，可把 $\mO^2/\La$ 看作 $k^4=\mO^2/\varpi^2\mO^2$ 的 $2$ 维商空间。此外，乘以 $\varpi$ 在 $\mO^2/\varpi^2\mO^2$ 上诱导幂零自同态，它保持 $\mO^2/\La$。容易看出，$\mO^2/\varpi^2\mO^2$ 的每个在此幂零自同态下稳定的 $2$ 维商空间都这样得到。因此 \eqref{Gr2} 可等同于 $\Gr(2,4)$ 的某个闭子簇的 $\bF_q$-点集。

若 $F=\bQ_p$，情形更复杂，因为我们不能把 $\mO^2/\La$ 看作 $k=\bF_p$ 上的向量空间。为克服这一困难，自然考虑另一个集合
\begin{equation}\label{wGr2latt}
\{\La\subset\La'\subset \mO^2\mid \dim\mO^2/\La'=\La/\La'=1\}.
\end{equation}
有一个自然满射 $m$ 把 \eqref{wGr2latt} 送到 \eqref{Gr2latt}，即忘掉 $\La'$。容易验证，除 $\La=p\La_0$ 之外，$m^{-1}(\La)$ 只含一个元素。
 
与前面类似的推理表明，\eqref{wGr2latt} 等同于 Hirzebruch 曲面 $\bP(\mO_{\bP^1}(-1)+\mO_{\bP^1}(1))$ 的 $\bF_p$-点集。此外，集合 $p^{-1}(\La_0)$ 构成该曲面上一条 $(-2)$-曲线的 $\bF_p$-点集。沿这条曲线收缩，所得射影簇的 $\bF_p$-点集自然等同于 \eqref{Gr2latt}。
\end{ex}

现在我们给出 $G$ 的仿射 Grassmann 簇的精确刻画。本节余下部分，我们允许 $F$ 稍更一般，即假定
$F$ 是关于离散赋值完备的局部域，整数环为 $\mO$，剩余域 $k$ 是\textbf{完全}的，特征 $p>0$\footnote{在等特征情形，即 $\cha F=\cha k$ 时，此假设并非必要。此处加上它是为了统一处理。}。设 $\varpi\in\mO$ 为一致化元。

回忆 $k$-代数 $R$ 称为完全的，指 Frobenius 自同态 $R\to R, \ r\mapsto r^p$ 是双射。
对完全 $k$-代数 $R$，令 $W(R)$ 为 $R$ 的（$p$-典型）Witt 向量环。它是元素形如序列 $(r_0,r_1,r_2,\ldots)\in R^{\bN}$ 的环，加法与乘法由某些（复杂的）多项式给出。到第一个分量的投影 $W(R)\to R, \ (r_0,r_1,\ldots)\mapsto r_0$ 是满环同态，带有乘性（但\textbf{非加性}）的截面 $R\to W(R), \ r\mapsto [r]=(r,0,0,\ldots)$，称为 $r$ 的 Teichm\"uller 提升。若 $R$ 是完全环，$W(R)$ 中每个元素可唯一地写成 $\sum_{i\geq 0} [a_i]p^i$，故 $W(R)$ 可视为「以 $p$ 为变量、系数在 $R$ 中的幂级数环」。
例如，$W(\bF_p)=\bZ_p$。我们定义 $R$ 上系数在 $\mO$ 中的 Witt 向量环为
\begin{equation}\label{ramified Witt vector}
W_\mO(R):= W(R)\hat\otimes_{W(k)}\mO:= \underleftarrow\lim_n W_{\mO,n}(R),\quad W_{\mO,n}(R):=W(R)\otimes_{W(k)} \mO/\varpi^n.
\end{equation}
特别地，若 $R=\bar k$，我们有时把 $W_\mO(\bar k)$ 记作 $\mO_L$，把 $W_{\mO}(\bar k)[1/\varpi]$ 记作 $L$。
注意若 $\cha F=\cha k$，则 $W_\mO(R)\simeq R[[\varpi]]$，且 $r\mapsto [r]$ 是环同态。我们有时写
\begin{equation}\label{discs}
D_{n,R}=\Spec W_{\mO,n}(R),\quad D_R=\Spec W_\mO(R),\quad D_R^*=\Spec W_\mO(R)[1/\varpi],
\end{equation}
并把它们看作由 $\Spec R$ 参数化的（去心）圆盘族。
给定 $D_R$ 上两个 $G$-torsor $\mE_1$ 与 $\mE_2$，从 $\mE_1$ 到 $\mE_2$ 的一个修正（modification）是同构 $\mE_1|_{D_R^*}\simeq \mE_2|_{D_R^*}$。我们通常记作 $\mE_1\dashrightarrow\mE_2$，以示该同构并非定义在整个 $D_R$ 上。

我们需要在完全代数几何中工作。一些基础可参看 \cite[附录 A]{Zh1}、\cite{BS} 与 \cite[\S A.1]{XZ}。

令 $\Aff_k^\pf$ 表示完全 $k$-代数构成的范畴。我们通过 $\Aff_k^\pf$ 上的预层来定义 $G$ 在 $k$ 上的仿射 Grassmann 簇 $\Gr_G$，
\begin{equation}
\label{E:affGrass def}
\Gr_G(R)= \left\{(\mE,\beta)\left | \begin{array}{l}\mE \mbox{ 是 } D_R\mbox{ 上的 } G\mbox{-torsor，且 }  \\
 \beta: \mE\dashrightarrow \mE^0 \mbox{ 是一个修正}\end{array}\right.\right\},
\end{equation}
其中 $\mE^0$ 表示平凡 $G$-torsor。若群 $G$ 明确，为简便起见把 $\Gr_G$ 写作 $\Gr$。
\begin{thm}\label{affGrass}
仿射 Grassmann 簇 $\Gr$ 可由子函子的归纳极限表示：$\Gr=\underrightarrow\lim X_i$，其中 $X_i\to X_{i+1}$ 为闭浸入，$X_i$ 是射影簇的完全化。
\end{thm}
当 $\cha F=\cha k$ 时，这是经典结果，归功于 Beauville–Laszlo、Faltings 等人（参看 \cite{BL, Fa}）。在混合特征情形（即 $\cha F\neq \cha k$），这作为猜想出现于 \cite[附录 B]{Zh1}，并由 Bhatt–Scholze 证明（\cite{BS}）。此前 \cite{Zh1} 证明了较弱的结论：$X_i$ 是紧合代数空间的完全化；这已足以在此设定下表述几何 Satake。

\begin{rmk}
完全 $k$-概形范畴是完全 $k$-代数上预层范畴的全子范畴，见 \cite[引理 A.10]{Zh1}。故上述定理有意义。
\end{rmk}
\begin{rmk}读者或许会问，我们为何只把仿射 Grassmann 簇 $\Gr$ 看作 $\Aff_k^\pf$ 上的预层，而不像通常那样看作 $k$-代数全范畴 $\Aff_k$ 上的预层。在等特征情形这确实可行，也正是 \cite{BL, Fa} 所考虑的设定。但
在混合特征情形，由于对非完全的特征 $p$ 环取 Witt 向量环存在某些病态行为，我们不知道 $\Gr$ 作为 $\Aff_k$ 上预层的正确刻画。另一方面，过渡到完全化不改变概形的拓扑，故若只关心与 $\Gr$ 有关的拓扑问题（即 $\Gr$ 上的反常层），我们不会丢失任何信息。
\end{rmk}
\begin{rmk}
\label{R:scholzediamond}
若 $F=k((t))$，则存在定义在圆盘 $D$ 上的 Beilinson–Drinfeld Grassmann 簇（事实上，Beilinson–Drinfeld Grassmann 簇定义在 $D$ 关于 $k$ 的所有笛卡尔幂之上），使得上述定义的 $\Gr$ 恰是闭点 $s\in D$ 处的纤维。在混合特征情形，Scholze 在发明称为 diamond 的新几何结构之后，近期构造了类似的几何对象（参看 \cite{SW}）。借助这些对象，像 $\Spec \bQ_p\times \Spec \bQ_p$ 这样的纤维积才真正有意义，也才可以在这些纤维积上定义类似的 Beilinson–Drinfeld Grassmann 簇。
\end{rmk}

仿射 Grassmann 簇还有一个有用的解释，即把它看作齐性空间。
设 $H$ 为 $\mO$ 上有限型的仿射群概形。以 $L^+H$（相应地 $LH$）记 $H$ 的\textbf{射流群}（相应地\textbf{环路群}）。作为预层，
\begin{equation*}\label{E:jet and loop}
L^+H(R)=H(W_\mO(R)),\quad LH(R)=H(W_\mO(R)[1/\varpi]).
\end{equation*}
它可由一个仿射群概形（相应地归纳概形）表示。对 $r\geq 0$，令 $L^rH$ 为第 $r$ 个射流群，即
\begin{equation*}\label{E:rjet}
L^rH(R)=H(W_{\mO,r}(R)).
\end{equation*} 
则 $L^rH$ 可由某个代数 $k$-群的完全化表示（通常的 Greenberg 实现），且 $L^+H=\underleftarrow\lim L^rH$。

现在可在 $\Gr$ 上定义空间 $\Gr^{(\infty)}$，它由 $(\mE,\beta)\in \Gr(R)$ 连同平凡化 $\epsilon:\mE\simeq \mE^0$ 构成。可以证明，它是 $\Gr$ 上的 $L^+G$-torsor，且有典范同构 $\Gr^{(\infty)}\cong LG$（见 \cite[引理 1.3]{Zh1} 或 \cite[命题 1.3.6]{Zh2}）。因此可把 $\Gr$ 实现为 fpqc 商
\begin{equation}
\label{E:gr as quotient}
\Gr\cong LG/L^+G.
\end{equation}
由此 $L^+G$ 通过左乘自然作用在 $\Gr$ 上。就模空间问题 \eqref{E:affGrass def} 而言，这一作用即 $\beta$ 与自同构 $L^+G=\Aut(\mE^0)$ 的自然复合。
\medskip

 
下面定义 $\Gr$ 的一些闭子簇，例 \ref{minu Sch} 与例 \ref{quasiminuSch} 所研究的对象作为特例包含在内。
首先设 $\mE_1$ 与 $\mE_2$ 是 $D_{\bar k}=\Spec \mO_L$ 上两个 $G$-torsor，$\beta: \mE_1\dashrightarrow \mE_2$ 为修正。
我们按如下方式给 $\beta$ 配一个元素
$G(\mO_L)\backslash G(L)/G(\mO_L)$：选取同构 $\epsilon_1:\mE_1\simeq \mE^0$ 与 $\epsilon_2:\mE_2\simeq \mE^0$，便得到平凡 $G$-torsor 的自同构 $\epsilon_2\beta\epsilon_1^{-1}\in \Aut(\mE^0|_{D_{\bar k}^*})$，从而得到 $G(L)$ 中一个元素。$\epsilon_1$ 与 $\epsilon_2$ 的另选会把此元素左乘、右乘 $G(\mO_L)$ 中的元素。因此，
$\beta$ 给出 $G(\mO_L)\backslash G(L)/G(\mO_L)$ 中一个良定义的元素。经由双射 \eqref{cartan2}，给 $\beta$ 配一个元素
\[\inv(\beta)\in \xcoch(T)/W.\]
我们称 $\inv(\beta)$ 为 $\beta$ 的相对位置。
\begin{ex}\label{invbetaGLn}
设 $G=\GL_n$。按通常方式把 $D$ 上的 $\GL_n$-torsor 与秩 $n$ 的有限投射 $\mO$-模等同。则 $\inv(\beta)$ 可如下描述。
如例 \ref{coweightGLn} 那样等同 $\xcoch(T)\simeq \bZ^n$。
给定秩 $n$ 的两个有限自由 $\mO$-模 $\La_1$ 与 $\La_2$，以及 $\beta: \La_1[1/\varpi]\simeq \La_2[1/\varpi]$，总存在 $\La_1$ 的基 $(e_1,\ldots,e_n)$ 与 $\La_2$ 的基 $(f_1,\ldots,f_n)$，使 $\beta$ 由
$$\beta(e_i)=p^{m_i}f_i$$ 
给出，且 $m_1\geq m_2\geq\cdots\geq m_n$。此外，序列 $(m_1,\ldots,m_n)$ 与基的选取无关。于是定义 $\inv(\beta)=(m_1,\ldots,m_n)$。
\end{ex}

现在设 $\mE_1$ 与 $\mE_2$ 是 $D_R$ 上两个 $G$-torsor，$\beta: \mE_1\dashrightarrow \mE_2$ 为修正。
则对每个几何点 $x\in \Spec R$，经基变换得到三元组 $(\mE_1|_{D_x},\mE_2|_{D_x},\beta_x: \mE_1|_{D_x}\dashrightarrow \mE_2|_{D_x})$。令 $\inv_x(\beta):=\inv(\beta_x)$。

我们如下定义（球面）Schubert 胞腔 $\mathring{\Gr}_{\mu}\subset \Gr\otimes \bar k$：
\[\mathring{\Gr}_\mu:=\left\{(\mE,\beta)\in \Gr\otimes \bar k\mid \inv(\beta)=\mu\right\}.\]
可以证明它在 $\Gr\otimes \bar k$ 中局部闭，并且是过 $\mu(\varpi)$ 的 $L^+G\otimes \bar k$-轨道。其闭包记作 $\Gr_\mu$，称为（球面）Schubert 簇\footnote{注意文献中有时把 Schubert 簇记作 $\Gr_{\leq\mu}$ 或 $\overline{\Gr}_\mu$，而把 Schubert 胞腔记作 $\Gr_\mu$。希望我们的记号不至于太令人困惑。}。它是一些 Schubert 胞腔之并。于是可在 $\xcoch(T)/W$ 上定义偏序（通常称为 Bruhat 序）如下：
\[\la\leq \mu \quad \mbox{ 当且仅当 } \quad \Gr_{\la}\subset \Gr_\mu.\]
\begin{rmk}
通常，Bruhat 序用 $G$ 对应的根数据以组合方式定义。对我们后文的讨论，上述定义已经够用。
\end{rmk}
\begin{ex}
设 $G=\GL_n$。如例 \ref{invbetaGLn} 那样把 $\xcoch(T)/W$ 与 $\xcoch(T)^+$ 等同。
则 Bruhat 序可显式描述如下：$\mu_1=(m_1,\ldots,m_n)\leq \mu_2=(l_1,\ldots,l_n)$ 当
$$m_1+\cdots+m_j\leq l_1+\cdots+l_j, \ \ \ j=1,\ldots, n, \quad \mbox{ 且 } \ \ m_1+\cdots+m_n=l_1+\cdots+l_n.$$ 
\end{ex}

Schubert 簇 $\Gr_\mu$ 一般是定义在 $\bar k$ 上的射影簇\footnote{事实上，它定义在 $\mu$ 的定义域上。}。它完全光滑，当且仅当 $\Gr_\mu=\mathring{\Gr}_\mu$；此时 $\mu$ 是 $\xcoch(T)/W$ 中关于 Bruhat 序的极小元。这样的 $\mu$ 称为\textbf{极小}（minuscule）。

\begin{ex}
设 $G=\GL_n$。令 $\omega_i=1^i0^{n-i}$ 如例 \ref{minu Sch}，并令 $\omega_i^*=\omega_{n-i}-\omega_{n}$。
注意 $\inv(\beta)=\omega_i$ 当且仅当 $\beta$ 延拓为真正的映射 $\La_1\to \La_2$，使得 $\varpi \La_2\subset \La_1$，且 $\La_2/\La_1$ 是 $i$ 维 $k$-向量空间。类似地，$\inv(\beta)=\omega_i^*$ 当且仅当 $\beta^{-1}$ 诱导包含 $\varpi \La_1\subset \La_2\subset \La_1$，且 $\La_1/\La_2$ 的维数为 $i$。
不难看出 $\omega_i,\omega_i^*$ 是极小余权，且 $\Gr_{\omega_i}\cong \Gr_{\omega_{n-i}^*}\cong \Gr(n-i,n)$。取 $k$-点即回到例 \ref{minu Sch}。
\end{ex}

接下来我们简要讨论几何 Satake 的通常表述，更详细的讨论参看 \cite{Zh2}（至少在等特征情形）。
动机版本将在下一小节讨论。
为简单起见，本小节余下部分假定 $k=\bar{k}$，并把 $L^+G$ 记作 $K$、把 $L^nG$ 记作 $K_n$（于是与前面的记号相反，本小节余下部分中 $K$ 是投射代数群而非投射有限群）。令 $K^{(n)}=\ker(K\to K_n)$ 表示第 $n$ 级同余子群。我们固定系数环 $E$，使 $\on{char} k$ 在其中可逆。层均指 $E$-层。

我们先简要回忆仿射 Grassmann 簇上反常层范畴 $\on{P}_{K}(\Gr)$ 的定义，它通常称为\textbf{Satake 范畴}，记作 $\Sat_G$。
首先，可以选取表示 $\Gr=\underrightarrow\lim X_i$，使每个 $X_i$ 是 $K$-不变闭子概形（例如取 $X_i$ 为若干 Schubert 簇的有限并）。此外，$K$ 在 $X_i$ 上的作用经某个 $K_m$ 分解。
于是可定义
\[\Sat_G=\on{P}_{K}(\Gr)=\underrightarrow\lim\on{P}_{K}(X_i),\]
其中 $\on{P}_{K}(X_i)$ 定义为 $X_i$ 上 $K_m$-等变反常层的范畴（它与 $m$ 的选取无关，相差典范等价），连接函子 $\on{P}_{K}(X_i)\to \on{P}_{K}(X_{i'})$ 为沿闭浸入 $X_i\to X_{i'}$ 的前推。此范畴定义的详细讨论见 \cite[\S A.1.17]{Zh2}。

可以证明它是半单 Abel 幺半范畴，幺半结构由 Lusztig 的层卷积积给出。我们回忆其定义。令 $\Gr\tilde\times \Gr$ 为仿射 Grassmann 簇的扭积（也称卷积 Grassmann 簇）。就模空间问题而言，它定义为预层
\begin{equation*}\label{I:conv}
(\Gr\tilde\times\Gr)(R)=\left\{(\mE_1,\mE_2,\beta_1,\beta_2)\ \left|\ \begin{split}&\mE_1,\mE_2 \mbox{ 是 }  D_R\mbox{ 上的 } G\mbox{-torsor}, \\
&  \beta_1:\mE_1\dashrightarrow \mE^0, \beta_2:\mE_2\dashrightarrow  \mE_1\end{split}\right.\right\}.
\end{equation*}
存在自然映射
\begin{equation}\label{conv m0}
m:\Gr\tilde\times\Gr\to \Gr,\quad (\mE_1,\mE_2,\beta_1,\beta_2)\mapsto (\mE_2, \beta_1\beta_2).
\end{equation}
与自然投影
$$\pr_1:\Gr\tilde\times\Gr\to \Gr, \quad (\mE_1,\mE_2,\beta_1,\beta_2)\mapsto (\mE_1,\beta_1),$$
二者一起诱导同构
$(\pr_1,m):\Gr\tilde\times\Gr\cong \Gr\times \Gr$。
特别地，卷积 Grassmann 簇可表。映射 $m$ 通常称为卷积映射。

更一般地，存在 $r$ 重卷积 Grassmann 簇 $\Gr\tilde\times\cdots\tilde\times\Gr$，它分类 $G$-torsor 的修正链
\begin{equation}
\label{E:chain mod}
\mE_r\stackrel{\beta_r}{\dashrightarrow}\mE_{r-1}\stackrel{\beta_{r-1}}{\dashrightarrow}\cdots\stackrel{\beta_0}{\dashrightarrow}\mE_0=\mE^0.
\end{equation}
对每个 $j=1,\ldots,r-1$，存在到 $j$ 重卷积 Grassmann 簇的自然投影 $\pr_j$，即记住 $\mE_j\stackrel{\beta_j}{\dashrightarrow}\cdots\stackrel{\beta_0}{\dashrightarrow}\mE_0=\mE^0$；还存在 $r$ 重卷积映射
\begin{equation}\label{conv m}
m:\Gr\tilde\times\cdots\tilde\times\Gr\to \Gr,
\end{equation}
把 \eqref{E:chain mod} 送到 $(\mE_r, \beta_1\cdots\beta_r)$。如前，$K$ 通过 $\mE_0=\mE^0$ 的自同构自然作用在 $\Gr\tilde\times\cdots\tilde\times\Gr$ 上，且有 $K$-torsor $(\Gr\tilde\times\cdots\tilde\times\Gr)^{(\infty)}\to\Gr\tilde\times\cdots\tilde\times\Gr$，它分类 \eqref{E:chain mod} 连同同构 $\epsilon:\mE_r\simeq \mE^0$。

利用同构 \eqref{E:gr as quotient}，容易看出 $\Gr\tilde\times\Gr\cong \Gr^{(\infty)}\times^K\Gr=LG\times^K\Gr$，
其卷积映射由作用映射 $LG\times\Gr\to \Gr$ 诱导。
这促使我们定义 $\Gr$ 的 $K$-不变闭子集的扭积。即：
对闭子集 $X\subset \Gr$，令 $X^{(\infty)}$ 表示它在 $\Gr^{(\infty)}$ 中的原像。则 $X^{(\infty)}\to X$ 是 $K$-torsor；对每个 $n$，令 $X^{(n)}=X^{(\infty)}/K^{(n)}\subset LG/K^{(n)}$，它是 $X$ 上的 $K_n$-torsor。
于是给定两个闭子集 $X_1, X_2$，其中 $X_2$ 是 $K$-不变的，我们把它们的扭积记作
\begin{equation}
\label{E-conv Sch}
X_1\tilde\times X_2:= X_1^{(\infty)}\times^K X_2,
\end{equation}
它是 $\Gr\tilde\times\Gr$ 的闭子集（因而可表）。若 $X_2$ 拟紧，则存在充分大的整数 $m$，使 $K$ 在 $X_2$ 上的作用经 $K_m$ 分解。
此时 $X_1^{(\infty)}\times^K X_2=X_1^{(m)}\times^{K_m}X_2$。注意
若 $X_i=\Gr_{\mu_i}$（$\mu_i\in \xcoch(T)/W$），我们也可把 $\Gr_{\mu_1}\tilde{\times}\Gr_{\mu_2}$ 描述为
\[\Gr_{\mu_1}\tilde{\times}\Gr_{\mu_2}\cong \left\{(\mE_1,\mE_2,\beta_1,\beta_2)\in\Gr\tilde\times\Gr\mid \inv(\beta_1)\leq \mu_1, \inv(\beta_2)\leq \mu_2\right\}.\]

更一般地，对 $r$ 重卷积 Grassmann 簇中的闭子集 $X_1$ 与 $s$ 重卷积 Grassmann 簇中的 $K$-不变闭子集 $X_2$，可用 \eqref{E-conv Sch} 中同一公式定义它们的扭积，其中 $X_1^{(\infty)}$ 理解为 $X_1$ 在 $(\Gr\tilde\times\cdots\tilde\times\Gr)^{(\infty)}$ 中的原像。
特别地，若 $\mmu=(\mu_1,\ldots,\mu_r)$ 是 $G$ 的支配余权序列，可定义 Schubert 簇的扭积
$$\Gr_{\mmu}:=\Gr_{\mu_1}\tilde\times\cdots\tilde\times\Gr_{\mu_r}\subset\Gr\tilde\times\cdots\tilde\times\Gr,$$
而若 $\nu_\bullet=(\nu_1,\ldots,\nu_s)$ 是另一序列，则 $\Gr_{\mmu}\tilde\times\Gr_{\nu_\bullet}\cong\Gr_{\mmu,\nu_\bullet}$。
令 $|\mu_\bullet|=\sum \mu_i$。则卷积映射 \eqref{conv m} 诱导
\begin{equation}\label{conv m2}
m: \Gr_{\mu_\bullet}\to \Gr_{|\mu_\bullet|}, \quad \mE_r\stackrel{\beta_r}{\dashrightarrow}\cdots\stackrel{\beta_0}{\dashrightarrow}\mE_0=\mE^0\mapsto (\mE_r,\beta_1\cdots\beta_r),
\end{equation}
已知它是半小映射。这意味着
\begin{equation}
\label{E:dim est}
\dim (\Gr_{\mmu}\times_\Gr\Gr_{\nu_\bullet})\leq \frac{1}{2}(d_{\mmu}+d_{\nu_\bullet}),
\end{equation}
其中记 $d_{\mmu}=\dim \Gr_{\mmu}$（见 \cite[命题 3.1.10]{XZ}）。

现在，对 $\mA_1,\mA_2\in \on{P}_{K}(\Gr)$，
以 $\mA_1\tilde\boxtimes\mA_2$ 记 $\mA_1$ 与 $\mA_2$ 在 $\Gr\tilde\times\Gr$ 上的「外扭积」。它是 $\Gr\tilde\times \Gr$ 上唯一的反常层，使其到 $\on{Supp}(\mA_1)^{(m)}\times \on{Supp}(\mA_2)$ 的拉回同构于 $\mA_1$ 到 $\on{Supp}(\mA_1)^{(m)}$ 的拉回与 $\mA_2$ 的外积。这里 $\on{Supp}(\mA_i)$ 是 $\mA_i$ 的支撑，它是 $K$-不变的；$m$ 充分大，使 $K$ 在 $\on{Supp}(\mA_2)$ 上的作用经 $K_m$ 分解。
例如若 $\mA_i=\IC_{\mu_i}$，则 $\IC_{\mu_1}\tilde\boxtimes\IC_{\mu_2}$ 典范同构于 $\Gr_{\mu_1}\tilde\times\Gr_{\mu_2}$ 的相交上同调层。
于是 $\mA_1$ 与 $\mA_2$ 的卷积积定义为
\begin{equation}\label{Lus conv}
\mA_1\star\mA_2:=m_!(\mA_1\tilde\boxtimes\mA_2),
\end{equation}
其中 $m:\Gr\tilde\times\Gr\to \Gr$ 是卷积映射（由 \eqref{conv m0} 定义）。
先验地，这是 $\Gr$ 上的 $K$-等变复形；但由卷积映射 \eqref{conv m2} 的半小性，它是反常层。

扭积 $\Gr\tilde\times\Gr$ 「拓扑同构于」乘积 $\Gr\times\Gr$。若 $F=\bC((\varpi))$，可用 $\Gr$ 上的解析拓扑使这句话精确化。即在解析拓扑下，$K$-torsor $\Gr^{(\infty)}\to \Gr$ 平凡。一般情形可从动机上理解。
于是当 $E=\Ql$（$\ell\neq \cha k$）时，应用 K\"unneth 公式，可给超上同调函子
$$\on{H}^*(\Gr_G,-) : \on{P}_{K}(\Gr_G,\Ql)\to \on{Vect}_{\Ql}$$ 配上典范幺半结构（见 \cite[\S 2.3]{Zh1} 与 \cite[\S 5.2]{Zh2}）。

下面这个定理通常称为几何 Satake 等价。我们假定 $E=\Ql$。

\begin{thm}\label{intro:geomSat}
幺半函子 $\on{H}^*$ 分解为两个函子的复合：一个是从 $\Sat_G$ 到 Langlands 对偶群 $\hat{G}$ 在 $\Ql$ 上有限维表示范畴 $\on{Rep}_{\Ql}(\hat{G})$ 的幺半范畴等价，另一个是从 $\on{Rep}_{\Ql}(\hat{G})$ 到 $\Ql$ 上有限维向量空间范畴 $\on{Vect}_{\Ql}$ 的遗忘函子。在该等价下，Schubert 簇 $\Gr_\mu$ 的相交上同调层 $\IC_\mu$ 对应于 $\hat G$ 的「最高权」为 $\mu$ 的不可约表示 $V_\mu$。
\end{thm}


\begin{rmk}
\label{R: birth of the dual group}
(1) 事实上，更典范的做法是把 $G$ 的对偶群 $\hat G$ \textbf{定义}为 Tannakian 范畴 $(\Sat_G,\on{H}^*)$ 的 Tannakian 群，并把 $V_\mu$ \textbf{定义}为 $\on{H}^*(\Gr,\IC_\mu)$，配以重言作用 $\hat G=\Aut^\otimes(\on{H}^*)$。本文余下部分采取这一观点。

(2) 如 \cite[\S 5.3]{Zh2} 所述，Tannakian 群典范地配有带钉 $(\hat G, \hat B, \hat T, \hat X)$。我们简要回忆 $(\hat B, \hat T)$ 的构造，细节参看「同上」。
上同调 $\on{H}^*$ 上的分次给出余特征标 $\bG_m\to \hat G$，可以证明它是正则的。于是它的中心化子给出极大环面 $\hat T$，而 $\hat B\supset \hat T$ 是 $\hat G$ 中使该余特征标为支配的唯一点 Borel 子群。
\end{rmk}

当 $F=k((t))$ 是等特征局部域时，此定理是 Lusztig、Ginzburg、Beilinson–Drinfeld 与 Mirkovi\'{c}-Vilonen 等人工作的结果（参看 \cite{Lu,Gi1,BD,MV}）。当 $F$ 是混合特征局部域时，此定理在 \cite{Zh1} 中证明。

定理最难的部分，是在 $\on{P}_{K}(\Gr_G)$ 的幺半结构上构造交换性约束，使 $\on{H}^*$ 成为张量函子。在等特征情形，该构造依赖把卷积积解释为融合积，这关键地用到 $k$ 上代数曲线笛卡尔幂上的 Beilinson–Drinfeld Grassmann 簇的存在性。如注 \ref{R:scholzediamond} 所述，混合特征情形中类似的几何对象已由 Scholze 近期构造。Scholze 近期还宣布，可以用与等特征情形类似的方式，在混合特征情形定义融合积并建立几何 Satake。

\cite{Zh1} 所用的方法不同。我们用经典 Gelfand 技巧的范畴版构造交换性约束。这一想法并不新，已见于 \cite{Gi1}。因此即使在混合特征情形，我们确实有交换性约束的候选。问题在于：若不使用融合解释，就难以验证它应当满足的性质（例如六边形公理）。

我们的新观察是：这些性质成立，等价于仿射 Hecke 代数的一个数值结果。即在 \cite{LV,Lu2} 中，Lusztig 与 Vogan 对带对合的 Coxeter 系 $(W,S)$ 引入了某些多项式 $P^\sigma_{y,w}(q)$，它们类似于通常的 Kazhdan–Lusztig 多项式 $P_{y,w}(q)$。随后 \cite{Lu2} 中猜想：若 $(W,S)$ 是仿射 Weyl 群，$y,w$ 是 $W$ 中某些元素，
\[P_{y,w}^\sigma(q)=P_{y,w}(-q).\]

更多细节见「同上」。此猜想纯属组合，但 Lusztig 与 Yun \cite{LY} 的证明是几何的，并且实际上用到了等特征几何 Satake！我们于是走向相反的方向：证明该公式蕴含上述交换性约束是正确的。

因此我们在混合特征情形对定理 \ref{intro:geomSat} 的证明，用到了等特征几何 Satake。一个有趣的问题是直接证明上述组合公式，那将给出几何 Satake 在等特征与混合特征情形下的纯局部证明。此外，类似的策略应当也适用于证明分歧 $p$ 进群的几何 Satake（其等特征对应版本已在 \cite{Zh0, Ri} 中建立）。


\subsection{动机 Satake 范畴}
\label{SS: MotSat}
为解释几何 Satake 应用于志村簇算术几何背后的某些想法，我们想用动机的语言重新表述几何 Satake。
本小节须假定读者熟悉（纯）动机的基本理论。经典参考是 \cite{Kl}。

本小节假定 $k$ 是完全域。
我们将构造（数值）动机范畴中一个 $\bQ$-线性、半单、超 Tannakian 的子范畴 $\Sat^m$。把 $\Sat^m$ 与 $\Ql$ 作张量积，可以恢复（通常 Satake 范畴的一个变体）。
为避免某些复杂的几何，我们主要关注 $\GL_n$ 情形，并在本小节末尾简述处理一般约化群所需的修改。



令 $\on{Var}_k$ 表示 $k$ 上光滑射影簇的范畴。
令 $\on{Mot}^{\on{rat}}_k$（相应地 $\on{Mot}_k^{\on{num}}$）表示 $k$ 上关于有理（相应地数值）等价关系的纯动机构成的 $\bQ$-线性范畴。则 $\on{Mot}_k^*$（$*=\on{rat},\on{num}$）是对称幺半的 $\bQ$-线性伪 Abel 范畴，而 $\on{Mot}_k^{\on{num}}$ 是超 Tannakian 半单 $\bQ$-线性 Abel 范畴（见 \cite{Ja}）。依 \cite{Ja}，$\on{Mot}_k^{*}$ 中的对象记作三元组 $(X,p,m)$，其中 $X$ 是 $k$ 上光滑射影簇，$p$ 是 $X$ 的零次自对应环中的幂等元（用上述某个适当等价关系定义），$m\in\bZ$。于是有自然函子
\[h:\on{Var}_k^{\on{op}}\to \on{Mot}^*_k,\quad X\mapsto h(X):=(X,\id,0).\]
照常，把 $(X,\id,m)$ 写作 $h(X)(m)$。
函子 $h$ 经 $\on{Var}_k$ 关于普遍同胚的局部化分解。
注意 $\cha k=p>0$，从 $\on{Var}_k$ 到完全 $k$-概形范畴 $\on{Sch}_k^\pf$ 的完全化函子经此局部化分解，而且实际上诱导该局部化范畴到 $\on{Sch}_k^\pf$ 的全嵌入（由 \cite[推论 A.16、命题 A.17]{Zh1}，另见 \cite[引理 3.8、命题 3.11]{BS}）。其本质像记作 $\on{Var}^\pf_k$，由作为光滑射影概形之完全化的 $k$ 上完全概形组成。
因此得到函子
$$h:(\on{Var}^\pf_k)^{\on{op}}\to\on{Mot}^*_k.$$

回忆给定序列 $\mmu=(\mu_1,\ldots,\mu_r)$，有 Schubert 簇的扭积 $\Gr_{\mmu}=\Gr_{\mu_1}\tilde\times\cdots\tilde\times\Gr_{\mu_r}$，其维数记作 $d_\mmu$。当且仅当每个 $\mu_i$ 极小，$\Gr_\mmu$ 才是某光滑射影簇的完全化，从而定义 $\on{Var}^\pf_k$ 中的对象。在 $\GL_n$ 情形，$\Gr_\mmu$ 是逐次以完全 Grassmann 簇为纤维的纤维化，
分类满足
$$\La_0=\mO^n, \quad  \varpi^{d_i+1}\La_{i-1}\subset \La_i\subset \varpi^{d_i}\La_{i-1} \mbox{ （对某个 } d_i\in\bZ\mbox{）}$$
的格链 $\{\La_i, i=0,\ldots, r\}$。
（这里及本小节余下部分，为简便只描述 $\Gr_\mmu$ 的 $k$-点。）
令 $\on{H}(\Gr_\mmu, \Gr_{\nu_\bullet})$ 为以 $\Gr_{\mmu}\times_\Gr\Gr_{\nu_\bullet}$ 的、维数\textbf{恰为} $\frac{1}{2}(d_{\mmu}+d_{\nu_\bullet})$ 的不可约分支为基的 $\bQ$-向量空间。我们需要如下几何事实。
\begin{lem}
\label{L:faithful realization}
设 $m,n,r$ 为三个整数，满足 $2(m-n)=d_{\mmu}-d_{\nu_\bullet}$ 与 $2(n-r)=d_{\nu_\bullet}-d_{\la_\bullet}$。
则自然映射
\[\on{H}(\Gr_{\mmu},\Gr_{\nu_\bullet})\to \Hom_{\on{Mot}^*_k}(h(\Gr_{\mmu})(m),h(\Gr_{\nu_\bullet})(n))\]
是单射。此外，存在唯一的 $\bQ$-线性映射
$$\on{H}(\Gr_{\mmu},\Gr_{\nu_\bullet})\otimes \on{H}(\Gr_{\nu_\bullet},\Gr_{\la_\bullet})\to \on{H}(\Gr_{\mmu},\Gr_{\la_\bullet})$$ 
使下图交换
{\scriptsize
\[
\begin{tikzcd}[row sep=2.2em, column sep=0.3em, thick]
\on{H}(\Gr_{\mmu},\Gr_{\nu_\bullet})\otimes \on{H}(\Gr_{\nu_\bullet},\Gr_{\la_\bullet}) \arrow[r] \arrow[d] & \on{H}(\Gr_{\mmu},\Gr_{\la_\bullet}) \arrow[d] \\
\Hom_{\on{Mot}^*_k}(h(\Gr_{\mmu})(m),h(\Gr_{\nu_\bullet})(n))\otimes\Hom_{\on{Mot}^*_k}(h(\Gr_{\nu_\bullet})(n),h(\Gr_{\la_\bullet})(r)) \arrow[r] & \Hom_{\on{Mot}^*_k}(h(\Gr_{\mmu})(m),h(\Gr_{\la_\bullet})(r))
\end{tikzcd}
\]
}
\end{lem} 
\begin{proof}
我们只给出梗概。
由于 $\Gr_\mmu\times\Gr_{\nu_\bullet}$ 允许胞腔分解（用半无限轨道的卷积积，例如见 \cite[\S 3.2.5, \S 3.2.6]{XZ}），此时有理等价、同调等价与数值等价三者一致。于是只需证明 $\Gr_{\mmu}\times_\Gr\Gr_{\nu_\bullet}$ 的这些 $\frac{1}{2}(d_{\mmu}+d_{\nu_\bullet})$ 维不可约分支的轮类，在 $\Gr_{\mmu}\times\Gr_{\nu_\bullet}$ 的 étale 上同调中线性无关。而这可由卷积映射 \eqref{conv m2}（它是半小映射）的通常分解定理推出。

给定单射性以及对应对合成的定义，要证图的交换性，只需注意在 $\Gr_{\mmu}\times\Gr_{\nu_\bullet}\times\Gr_{\la_\bullet}$ 中，来自 $\on{H}(\Gr_{\mmu},\Gr_{\nu_\bullet})$ 与 $\on{H}(\Gr_{\nu_\bullet},\Gr_{\la_\bullet})$ 的轮之相交是真交，然后应用维数估计 \eqref{E:dim est}。
\end{proof}



现在定义 $\bQ$-线性加性范畴 $\Sat_G^{0}$，并把\textbf{动机 Satake 范畴} $\Sat_G^m$ 定义为 $\Sat_G^{0}$ 的幂等完备化。我们指出，定义 $\Sat_G^0$ 的想法已见于 \cite[\S 4.3]{FKK}。
\begin{itemize}
\item 对象是数对 $(X,f)$，其中 $X$ 是形如 $\Gr_\mmu$ 的完全光滑射影概形的不交并（$\mmu$ 取遍极小余权序列），$f$ 是 $X$ 上取值于 $\bZ$ 的局部常值函数。若 $X=\Gr_\mmu$，则 $f$ 由一个整数 $m$ 给出；此时把该数对记作 $(\Gr_\mmu,m)$。若还有 $m=0$，有时简记 $(\Gr_\mmu,0)$ 为 $\Gr_\mmu$。
\item \[\Hom_{\Sat_G^0}((\Gr_{\mmu},m),(\Gr_{\nu_\bullet},n))=\left\{\begin{array}{ll}0 & 2(m-n)\neq d_\mmu-d_{\nu_\bullet}\\
\on{H}(\Gr_{\mmu},\Gr_{\nu_\bullet}) & 2(m-n)=d_\mmu-d_{\nu_\bullet} \end{array}\right.\]
若把 $(X,f)=\sqcup_i (X_i,m_i)$ 与 $(Y,g)=\sqcup (Y_j,n_j)$ 写成连通对象的不交并，则定义
$$\Hom_{\Sat_G^0}((X,f),(Y,g))=\bigoplus_{ij}\Hom_{\Sat_G^0}((X_i,m_i),(Y_j,n_j)).$$
\end{itemize}

由引理 \ref{L:faithful realization} 可知，这样定义确实得到一个范畴。自然函子
$$\Sat_G^0\to \on{Mot}_k^*: (X,f)\mapsto \oplus h(X_i)(m_i)$$
是忠实的（但非全的）加法函子。
以下，给定维数为 $\frac{1}{2}(d_{\mmu}+d_{\nu_\bullet})$ 的不可约分支 $Z\subset \Gr_{\mmu}\times_\Gr\Gr_{\nu_\bullet}$，我们用 $[Z]$ 记它在 $\Sat^0_G$ 中诱导的态射。注意，为使用引理 \ref{L:faithful realization}，我们确实需要把 $(\Gr_\mmu,m)$ 的各种「Tate 扭转」也纳入 $\Sat^0_G$ 的对象。

接下来，我们给 $\Sat_G^{0}$ 赋予幺半结构（它与 $\on{Mot}_k^*$ 中通常的张量积结构先验地不同）。
注意对角态射 $\Delta: \Gr_{\mmu}\to \Gr_{\mmu}\times_{\Gr}\Gr_{\mmu}$ 给出 $(\Gr_\mmu,m)$ 在 $\Sat_G^{0}$ 中的恒等态射 $1_{(\Gr_\mmu,m)}$。
我们在连通对象上定义张量积为
\[(\Gr_{\mmu},m)\star (\Gr_{\nu_\bullet},n):=(\Gr_{\mmu,\nu_\bullet},m+n),\]
再按线性性自然延拓到所有对象。
为在态射上定义张量积，只需定义 $1_{(\Gr_{\mmu},m)}\star [Z]$ 与 $[Z']\star 1_{(\Gr_{\nu_\bullet},n)}$，其中 $[Z]\in \Hom_{\Sat_G^{0}}((\Gr_{\nu_\bullet},n'), (\Gr_{\nu'_\bullet},n'))$（相应地 $[Z'] \in \Hom_{\Sat_G^{0}}((\Gr_\mmu,m), (\Gr_{\mu'_\bullet},m'))$）由不可约分支代表。我们定义
$$1_{(\Gr_{\mmu},m)}\star [Z]:=[\Delta\tilde\times Z],\quad [Z']\star 1_{(\Gr_{\nu_\bullet},n)}:=[Z'\tilde\times\Delta],$$
此处用到 $Z$ 是 $K$-不变这一事实，以构造由 \eqref{E-conv Sch} 定义的扭积。

\begin{lem}
\begin{enumerate}
\item 自然（非全）嵌入 $\Sat_G^0\subset\on{Mot}_k^*$ 带有典范的幺半结构，即存在典范同构
\[h(\Gr_{\mmu,\nu_\bullet})(m+n)\cong h(\Gr_{\mmu})(n)\otimes h(\Gr_{\nu_\bullet})(n),\]
它位于 $\on{Mot}^*_k$ 中，并与两个范畴各自的幺半结构相容。
\item $\Sat_G^{0}$ 中存在唯一的交换性约束，使上述幺半函子成为对称幺半函子。
\item 函子 $\Sat_G^0\to\on{Mot}^*_k$ 延拓为对称幺半函子
$$\Sat_G^m\to \on{Mot}_k^*.$$
\end{enumerate}
\end{lem}

由于 $\on{Mot}^*_k$ 幂等完备，而按定义 $\Sat_G^m$ 是 $\Sat_G^0$ 的幂等完备化，最后一条由前两条推出。
把 \cite[\S 2.3, \S 2.4]{Zh1} 中的（等变）上同调换成（等变）Chow 环（参看 \cite{EG}），可以构造映射 $A^*(\Gr_{\mmu})\otimes_{\bQ}A^*(\Gr_{\nu_\bullet})\to A^*(\Gr_{\mmu,\nu_\bullet})$ 与 $A^*(\Gr_{\mmu,\nu_\bullet})\to A^*(\Gr_{\nu_\bullet,\mmu})$，其中 $A^*(-)$ 表示有理 Chow 群。由于所有这些簇都容许胞腔分解，这些映射给出 (1) 与 (2) 中所求结构的候选。
过渡到上同调并利用 \cite[\S 2.3, \S 2.4]{Zh1}，即可看出它们确实是所需的结构。
我们不回顾全部细节，只举一个例子。

\begin{ex}
\label{R: comm via cycle}
设 $G=\GL_2$。
考虑 $\Gr_{\omega_1,\omega_1}$，它分类满足「$\La_i/\La_{i+1}$ 为一维」的格链 $\{\La_2\subset\La_1\subset\La_0=\mO^2\}$。如例 \ref{quasiminuSch} 所述，它同构于 $\bP(\mO_{\bP^1}(-1)\oplus\mO_{\bP^1}(1))$ 的完全化（详细讨论见 \cite[\S B.3]{Zh1}）。回忆其中有一条由条件 $\La_2=\varpi\La_0$ 定义的 $(-2)$-曲线，记作 $\Gr_{\omega_1,\omega_1}^{\omega_2}\subset \Gr_{\omega_1,\omega_1}$。（更确切地说，两种情形下我们都考虑由 $\bP(\mO_{\bP^1}(-1)\oplus\mO_{\bP^1}(1))$ 给出的去完全化；于是对应的约化闭子概形是 $\bP(\mO_{\bP^1}(-1)\oplus\mO_{\bP^1}(1))$ 中的一条 $(-2)$-曲线。）

同构 $h(\Gr_{\omega_1,\omega_1})\cong h(\Gr_{\omega_1})\otimes h(\Gr_{\omega_1})$ 可由轮
$$Z_1+Z_2\subset \Gr_{\omega_1,\omega_1}\times \Gr_{\omega_1}\times \Gr_{\omega_1}$$
实现，其中
$$Z_1=\{(z,\pr_1(z),\pr_1(z))\mid z\in \Gr_{\omega_1,\omega_1}\}, \quad Z_2=\{ (z, \pr_1(z), y)\mid z\in \Gr_{\omega_1,\omega_1}^{\omega_2}, y\in \Gr_{\omega_1}\}.$$
由交换性约束 $c$ 诱导的 $\Gr_{\omega_1,\omega_1}$ 在 $\Sat_G^0$ 中的自同构
\[c':\Gr_{\omega_1,\omega_1}= \Gr_{\omega_1}\star \Gr_{\omega_1}\stackrel{c}{\cong} \Gr_{\omega_1}\star \Gr_{\omega_1}=\Gr_{\omega_1,\omega_1}\]
由
\begin{equation}
\label{E:comm via cycle}
c'=[\Delta]+[\Gr_{\omega_1,\omega_1}^{\omega_2}\times \Gr_{\omega_1,\omega_1}^{\omega_2}]
\end{equation}
给出。这里对角轮 $\Delta$ 与 $\Gr_{\omega_1,\omega_1}^{\omega_2}\times \Gr_{\omega_1,\omega_1}^{\omega_2}$ 是 $\Gr_{\omega_1,\omega_1}\times_{\Gr}\Gr_{\omega_1,\omega_1}$ 的不可约分支，构成
$\Hom_{\Sat_G^{0}}(\Gr_{\omega_1,\omega_1},\Gr_{\omega_1,\omega_1})$ 的一组基。
\end{ex}

下面说明 $\Sat_G^m$ 与通常 Satake 范畴的关系。
设 $\Sat^T_G\subset \on{P}_K(\Gr)$ 为由 $\IC_\mu(i)$（$\mu\in\xcoch(T)/W$，$i\in\bZ$）张成的全子范畴，其中 $\IC_\mu$ 是 $\Gr_\mu$ 上的相交上同调层（它在 Schubert 胞腔上的限制为 $\Ql[d_\mu]$）。由 \cite[引理 5.5.14]{Zh2}（其论证在混合特征下同样成立，无须改动），可以对几何 Satake 作自举，得到如下交换图
\[
\begin{CD}
\Sat^T_G@>\cong>>\on{Rep}_{\Ql}(\hat G^T)\\
@VVV@VVV\\
\Sat_{G\otimes\mO_L}@>\cong>>\on{Rep}_{\Ql}(\hat G).
\end{CD}
\]
其中
\begin{itemize}
\item 群 $\hat{G}^T$ 是 $\hat{G}$ 与 $\bG_m$ 的半直积，$\bG_m$ 对 $\hat G$ 的作用定义如下（细节见 \cite[(5.5.10)]{Zh2}）：记 $\hat G_\ad$ 为 $\hat G$ 的伴随商，它通过共轭作用于 $\hat G$。则复合映射 $\bG_m\to \hat G\to\hat G_\ad$（第一个映射是注 \ref{R: birth of the dual group} (2) 中的余特征标）有平方根 $\rho_\ad:\bG_m\to \hat G_\ad$，后者诱导 $\bG_m$ 在 $\hat G$ 上的一个作用。注意 $\hat G^T$ 带有 Borel 子群 $\hat B^T:=\hat B\rtimes\bG_m$ 与极大环面 $\hat T^T=\hat T\rtimes \bG_m=\hat T\times \bG_m$。
\item 底部的等价即定理 \ref{intro:geomSat} 中的等价。
\item 左侧竖直函子分别是沿 $\Gr\otimes\bar k \to \Gr$ 的拉回函子与沿 $\hat G\to \hat G^T$ 的限制函子。
\end{itemize}

由卷积映射 \eqref{conv m2} 的半小性与分解定理，存在全忠实幺半函子
$$\Sat_G^0\otimes\Ql\to \Sat_G^T,\quad (\Gr_{\mmu},m)\mapsto m_*\Ql_{\Gr_{\mmu}}[d_{\mmu}](m)\in \on{P}_K(\Gr),$$
由 \cite[引理 2.16]{Zh1}，它延拓为等价
$$\Sat^m_G\otimes\Ql\cong \Sat^T_G.$$
因此，按 \cite[(2.4.4)]{Zh1} 对交换性约束作 Koszul 符号改变之后，得到对称幺半范畴的等价
\begin{equation*}
\label{E:geom sat Tate}
\Sat^m_G\otimes\Ql\cong \on{Rep}(\hat{G}^T),
\end{equation*}
我们把它看作动机版的几何 Satake。

\begin{rmk}
\label{R:Tate twist}

如上所述，为了把 $\Sat_G^0$ 定义为 $\on{Mot}^*_k$ 的子范畴，我们确实需要引入 Tate 扭转。因此 $\Sat_G^m\otimes\Ql$ 的 Tannakian 群不是原来的 Langlands 对偶群 $\hat{G}$，而是修正后的 $\hat G^T$。近年来的研究逐渐认识到，在 Langlands 对应的表述中，$\hat G^T$ 或许是更正确的对象（例如见 \cite{BG}）。在某些情形下，可以选取同构 $\hat{G}^T\simeq \hat{G}\times \bG_m$。例如当 $G=\GL_n$ 时，这样的同构可由
\begin{equation}
\label{E: dual group GLn}
\hat{G}^T:=\GL_n\rtimes\bG_m\to \GL_n\times \bG_m,\ (A,t)\mapsto (A\on{diag}\{t^{n-1},t^{n-2},\ldots,1\}, t)
\end{equation}
给出。关于 Langlands 对偶群不同版本的更详细讨论，参看 \cite[\S 5.5]{Zh2}。

该等价诱导 $\bZ[v,v^{-1}]$-代数的同构
$$K(\Sat^m_G)\cong K(\on{Rep}(\hat{G}^T)),$$
其中左端的 $\bZ[v,v^{-1}]$ 由全子范畴 $\{(\Gr_0,m), m\in\bZ\}$ 到 $\Sat_G^m$ 的嵌入诱导，右端的 $\bZ[v,v^{-1}]$ 由全子范畴 $\on{Rep}(\bG_m)\subset \on{Rep}(\hat G^T)$ 的嵌入诱导。于是取 $v=q$ 便得到\emph{典范的} Satake 同构；与 \eqref{Sat isom} 不同，它不依赖于任何选取。当 $\GL_2$ 时，该同构连同 \eqref{E: dual group GLn} 诱导出注 \ref{R: norm Sat} 中的归一化 Satake 同构 $\Sat'^{cl}$。
\end{rmk}

\begin{ex}\label{intersection by geom Sat}
下面是后文应用几何 Satake 的玩具模型。

仍取 $G=\GL_2$。设 $\Gr_{\omega_1,\omega_1^*}$（相应地 $\Gr_{\omega_1^*,\omega_1}$）是分类格链
\begin{equation}\label{Chain 2}
\{\La_2\supset\La_1\subset \La_0=\mO^2\}, \quad \quad \mbox{相应地 } \quad \{\La'_2\subset\La'_1\supset \La_0=\mO_2\}
\end{equation}
的模空间，其中 $\La_0/\La_{1}$ 与 $\La_2/\La_1$（相应地 $\La'_1/\La_0$ 与 $\La'_1/\La'_2$）都是 $k$ 上的一维空间。\zy{[译者按（上下文推断，无外部依据）：此处原文第二条链写作 $\La_0=\mO_2$，据上下文应为 $\La_0=\mO^2$。]}
注意有典范同构
\begin{equation}
\label{E: isomisom1}
\Gr_{\omega_1,\omega_1^*}\cong\Gr_{\omega_1,\omega_1}, \quad \La_2\supset\La_1\subset \La_0\mapsto \varpi\La_2\subset\La_1\subset\La_0,
\end{equation}
与
\begin{equation}
\label{E: isomisom2}
\Gr_{\omega_1^*,\omega_1}\cong\Gr_{\omega_1,\omega_1}, \quad \La'_2\subset\La'_1\supset \La_0\mapsto \varpi\La'_2\subset\varpi\La'_1\subset\La_0,
\end{equation}
故 $\Gr_{\omega_1,\omega_1^*}$ 与 $\Gr_{\omega_1^*,\omega_1}$ 都同构于 $\bP(\mO_{\bP^1}(-1)\oplus\mO_{\bP^1}(1))$ 的完全化。注意
$$\Gr_{\omega_1,\omega_1^*}^0:=\Gr_{\omega_1,\omega_1^*}\times_{\Gr}\Gr_0,\quad\quad \mbox{相应地 }\quad\Gr_{\omega_1^*,\omega_1}^0:=\Gr_{\omega_1^*,\omega_1}\times_{\Gr}\Gr_0$$
是 $\Gr_{\omega_1,\omega_1^*}$（相应地 $\Gr_{\omega_1^*,\omega_1}$）中分类 \eqref{Chain 2} 里满足 $\La_2=\La_0$（相应地 $\La'_2=\La_0$）的格链的闭子概形，因而同构于 $\bP^1$ 的完全化；如例 \ref{quasiminuSch} 所述，它是该曲面上的 $(-2)$-曲线。\zy{[译者按（上下文推断，无外部依据）：此处原文两个下标记作 $\omega_1^*$ 与 $\omega_2$，据上下文应分别为 $\omega_1^*$ 与 $\omega_1$，译文按后者书写。]}

我们考虑 $\Sat_G^{0}$ 中由以下映射复合而成的态射
\[\Gr_0\xrightarrow{\Gr_{\omega_1,\omega_1^*}^0} \Gr_{\omega_1,\omega_1^*}=\Gr_{\omega_1}\star\Gr_{\omega_1^*}\stackrel{c}{\cong}\Gr_{\omega_1^*}\star\Gr_{\omega_1}=\Gr_{\omega_1^*,\omega_1}\xrightarrow{\Gr_{\omega_1^*,\omega_1}^0}\Gr_0,\]
这里把 $\Gr_{\omega_1,\omega_1^*}^0$ 看作 $\Hom(\Gr_0, \Gr_{\omega_1,\omega_1^*})$ 中的元素，把 $\Gr_{\omega_1^*,\omega_1}^0$ 看作 $\Hom(\Gr_{\omega_1^*,\omega_1},\Gr_0)$ 中的元素，$c$ 是交换性约束。结合同构 \eqref{E: isomisom1} 与 \eqref{E: isomisom2}、交换性约束 \eqref{E:comm via cycle}，以及 $\Gr_{\omega_1,\omega_1}^{\omega_2}$ 是 $\Gr_{\omega_1,\omega_1}$ 中 $(-2)$-曲线这一事实，可以看出该复合等于乘 $2$。

另一方面，按几何 Satake，该映射可由 $\GL_2\times \bG_m$ 的有限维表示范畴中的映射
\[\mathbf{1}\to \on{Std}\otimes \on{Std}^*\cong \on{Std}^*\otimes \on{Std} \to \mathbf{1}\]
计算，而后者同样等于乘 $\dim \on{Std}=2$。总之，（卷积）仿射 Grassmann 簇中某些轮之间的相交数，可以借助 Langlands 对偶群的表示论算出。
\end{ex}

最后，我们简要说明如何把上述构造从 $\GL_n$ 推广到 $\mO$ 上的一般约化群 $G$。对一般的约化群，极小余权不够多。例如 $G=E_8$ 时没有非零的极小余权。
故一般情形下，我们把 $\Sat^0_G$ 定义为这样的范畴：对象是 $(\Gr_{\mmu},m)$ 的不交并，其中 $\mmu=(\mu_1,\ldots,\mu_r)$，每个 $\mu_i$ 极小或拟极小（拟极小余权的概念见例如 \cite[\S 2.2.2]{Zh1}）；$(\Gr_{\mmu},m)$ 与 $(\Gr_{\nnu},n)$ 之间的态射由与前面相同的公式给出。

我们需要说明这样确实得到一个范畴（即态射可以复合），并且它有到 $\on{Mot}_k^*$ 的忠实嵌入。
若 $\mu$ 拟极小，$\Gr_\mu$ 不是完全光滑的，而是某个部分旗簇 $G/P$ 上射影锥的完全化；在锥的奇点处爆破即给出奇点的一个显式消解（例如见 \cite[引理 2.12]{Zh1}）。记 $\pi:\wGr_\mu\to \Gr_\mu$ 为这个「消解」，由构造它是 $K$-等变的。不难写出幂等元 $p$，它可以由 $\wGr_\mu\times\wGr_\mu$ 中一个 $K$-不变的轮代表，使得 $(\wGr,p,0)$ 的 étale 实现给出 $\Gr_\mu$ 的相交上同调（由分解定理，它是 $\on{H}^*(\wGr_\mu)$ 的直和项）。
于是有与引理 \ref{L:faithful realization} 类似的引理，只需作显然的修改。重复前面的构造，即得动机 Satake 范畴 $\Sat^m_G$。

\section{范畴迹的构造}
近年来，范畴中心/范畴迹的概念在表示论中扮演了重要角色。
关于它们在特征层理论中的应用，我们参看 \cite{BN,BFO,Lu3,Lu4,Lu5,Lu6}；关于在稳定 Bernstein 中心中的应用，参看 \cite{BKV}。本文余下部分基于与肖梁的合作工作 \cite{XZ}，说明这些范畴构造的另一种视角与应用。

\subsection{范畴迹}
\label{S: cat trace}
为说明定义的动机，我们先简要回顾代数的迹（有时也称为余中心）的概念。
设 $E$ 是基交换环，$A$ 是 $E$-代数。从 $A$ 到 $E$-模 $V$ 的迹函数指的是满足 $f(ab)=f(ba)$ 的 $E$-线性映射 $f: A\to V$。
$A$ 的迹是一个 $E$-模 $\Tr(A)$，配以迹函数 $\tr: A\to \Tr(A)$，使得每个迹函数 $f:A\to V$ 都唯一地分解为 $A\xrightarrow{\tr} \Tr(A)\xrightarrow{\bar{f}}V$，其中 $\bar{f}:\Tr(A)\to V$ 是某个线性映射。显然，若把 $\Tr(A)$ 定义为 $A$ 模去由形如 $ab-ba$（$a,b\in A$）的元素张成的 $E$-子模所得的商，并令 $\tr$ 为该商映射，则 $(\Tr(A),\tr)$ 就是 $A$ 的迹。等价地，也可以把 $\Tr(A)$ 表示成 $E$-模范畴中下图之余极限
\begin{equation}
\label{E:HH}
\Tr(A)=\on{colim}(A\otimes A\substack{\longrightarrow\\[-1em] \longrightarrow}  A)
\end{equation}
其中两个映射 $A\otimes A\to A$ 分别由 $a\otimes b\mapsto ab$ 与 $a\otimes b\mapsto ba$ 给出。注意这个图正是计算 $A$ 的 Hochschild 同调的 bar 复形的前两项。

下面转到范畴情形，引入幺半范畴的范畴迹概念。
基本想法如下：幺半范畴是某类范畴构成的对称幺半 $2$-范畴中的代数对象；一般地，应当把对称幺半（高阶）范畴中代数对象的迹定义为其 bar 复形 \eqref{E:HH} 的余极限。借助稳健的高阶范畴论，这一想法在相当一般的设定下都可行。例如，我们参看 \cite{BN} 中精彩的论述。

对我们后来的应用而言，只需在通常的 $E$-线性范畴层面引入这一概念，并且重要的是显式地构造范畴迹。另一方面，我们还需要幺半自函子的扭范畴迹的概念。因此我们详细展开定义。我们将在本质小的 $E$-线性范畴构成的 $2$-范畴中工作。
设 $(\mC,\otimes,\mathbf{1})$ 是 $E$-线性幺半范畴，$\sigma: \mC\to \mC$ 是幺半自函子。对每个 $n$，记 $\mC^{\otimes n}$ 为这样的 $E$-线性范畴：对象是 $(X_i)_{i=1,\ldots,n}$，态射为 $\Hom_{\mC^{\otimes n}}((X_i)_i, (Y_i)_i)=\otimes_i \Hom_\mC(X_i,Y_i)$。则左 $\sigma$-扭范畴迹 $\Tr_\sigma(\mC)$ 定义为下图之 $2$-余极限
\[
\mC^{\otimes 3}\ \substack{\longrightarrow\\[-1em] \longrightarrow\\[-1em] \longrightarrow} \ \mC^{\otimes 2}\ \substack{\longrightarrow\\[-1em] \longrightarrow} \ \mC,
\]
其中函子 $\mC^{\otimes 2}\to \mC$ 分别由 $(X,Y)\mapsto X\otimes Y$ 与 $Y\otimes \sigma X$ 给出，而 $\mC^{\otimes 3}\to \mC^{\otimes 2}$ 分别由 $(X,Y,Z)\mapsto (X,Y\otimes Z)$、$(Y, Z\otimes \sigma X)$ 与 $(X\otimes Y,Z)$ 给出。

更具体地，设 $\mD$ 是（通常的）$E$-线性范畴。
从 $\mC$ 到 $\mD$ 的左 $\sigma$-扭\emph{迹函子}指的是一个函子 $F:\mC\to\mD$，连同典范同构族（对每个变元都是函子性的）
\[\al_{X,Y}: F(X\otimes Y)\cong F(Y\otimes \sigma X)\]
使得下图交换
\begin{equation}
\label{E:cond trace}
\xymatrix{
F((X\otimes Y)\otimes Z)\ar^{\al_{X\otimes Y,Z}}[r]\ar_\cong[d] & F(Z\otimes \sigma (X\otimes Y))\ar^{\cong}[r]&  F((Z\otimes \sigma X)\otimes \sigma Y)  \\
F((X\otimes (Y\otimes Z))\ar^{\al_{X,Y\otimes Z}}[r] & F((Y\otimes Z)\otimes \sigma X)\ar^{\cong}[r]&  F(Y\otimes (Z\otimes \sigma X)). \ar_{\al_{Y,(Z\otimes\sigma X)}}[u]
}
\end{equation}
注意由此可推出下面两个图也交换
\begin{equation}
\label{E:cond trace cor}
\xymatrix{
F(\mathbf{1}\otimes X)\ar_{\cong}[d]\ar^{\al_{\mathbf{1},X}}[r]& F(X\otimes \sigma\mathbf{1})\ar^\cong[d]&  F(X\otimes Y)\ar^{\al_{Y,\sigma X}\al_{X,Y}}[r]\ar_\cong[d]& F(\sigma X\otimes \sigma Y)\ar^\cong[d]\\
 F(X)\ar^-\cong[r]& F(X\otimes \mathbf{1})&  F((X\otimes Y)\otimes\mathbf{1})\ar^{\al_{X\otimes Y,\mathbf{1}}} [r]& F(\mathbf{1}\otimes\sigma(X\otimes Y)).
}\end{equation}
记 $\on{Func}^{\tr_\sigma}(\mC,\mD)$ 为左 $\sigma$-扭迹函子构成的范畴。则 $\Tr_\sigma(\mC)$（若存在）是满足下述性质的唯一范畴（在同构意义下唯一，且同构唯一）：对任意（通常的）$E$-线性范畴 $\mD$，
\[\on{Func}^{\tr_\sigma}(\mC,\mD)\cong \on{Func}(\Tr_\sigma(\mC),\mD),\]
其中 $\on{Func}(-,-)$ 表示 $E$-线性范畴之间的 $E$-线性函子范畴。

对我们的应用而言，需要 $\Tr_\sigma(\mC)$ 的同调空间的显式描述。通常，描述范畴余极限图的同调空间很困难。但在我们的情形，对 $\mC$ 作一点温和的假设即可做到。
\begin{prop}
\label{L:trace cat}
设 $\mC$ 是小 $E$-线性范畴，且其中每个对象都有左对偶对象。则 $(\Tr_\sigma(\mC),\tr)$ 存在。若 $\mC$ 是加性的，则 $\Tr_\sigma(\mC)$ 亦然。
\end{prop}
回忆在幺半范畴中，对象 $V$ 的左对偶指的是对象 $V^*$ 连同态射
$$\coev_V:\mathbf{1}\to V\otimes V^*,\quad \ev_V: V^*\otimes V\to \mathbf{1}$$
使得下面两个复合都等于恒等映射
\[V\cong \mathbf{1}\otimes V\xrightarrow{\coev_V\otimes\id_{V}}(V\otimes V^*)\otimes V\cong V\otimes (V^*\otimes V)\xrightarrow{\id_{V}\otimes\ev_V}V\otimes\mathbf{1}\cong V\]
\[V^*\cong V^*\otimes\mathbf{1}\xrightarrow{\id_{V^*}\otimes\coev_V} V^*\otimes (V\otimes V^*)\cong (V^*\otimes V)\otimes V^*\xrightarrow{\ev_V\otimes\id_{V^*}}\mathbf{1}\otimes V^*\cong V^*.\]
三元组 $(V^*,\coev_V,\ev_V)$ 在同构唯一的意义下唯一。

\begin{proof}
我们给出详细证明，因为它将启发下一小节中 $S$-算子的构造。为简化记号，以下所有公式与图表中都省略结合约束。

我们先定义范畴 $\Tr'_\sigma(\mC)$，其对象与 $\mC$ 的对象相同。为避免混淆，当 $X\in\mC$ 被看作 $\Tr'_\sigma(\mC)$ 的对象时，记作 $\widetilde{X}$。态射空间定义为
\begin{equation*}
\label{E:trace hom}
\Hom_{\Tr'_\sigma(\mC)}(\widetilde{X},\widetilde{Y}):=\big(\bigoplus_{(V,W)\in \mC\times \mC}\Hom(X, V\otimes W)\otimes_E \Hom(W\otimes \sigma V, Y)\big)/\sim,
\end{equation*}
其中等价关系由下述方式生成：若 $u: X\to V\otimes W$ 与 $v:W\otimes \sigma V\to Y$ 对应的元素记作 $S_{u,v}\in \Hom_{\Tr'_\sigma(\mC)}(\widetilde{X},\widetilde{Y})$，则当存在交换图
\[\xymatrix{
X\ar^-{u}[r] \ar_{u'}[dr]& V\otimes W\ar[d] & W\otimes \sigma V\ar[d]\ar^-{v}[r] & Y\\
& V'\otimes W' & W'\otimes \sigma V'\ar_{v'}[ur] &
}\]
（其中竖直映射由 $\Hom_\mC(V,V')\otimes_E\Hom_{\mC}(W,W')$ 中的元素诱导）时，令 $S_{u,v}=S_{u',v'}$。特别地，由下面两个交换图
\begin{equation}
\label{E:frac}
\xymatrix{
&V\otimes V^*\otimes X\ar^{\id_V\otimes u^*}[d]& V^*\otimes X\otimes \sigma V\ar^{v_u}[dr]\ar_{u^*\otimes \id_{\sigma V}}[d]&\\
X\ar^{\coev_V\otimes\id_X}[ur]\ar^{u}[r]\ar_{u_v}[dr]&V\otimes W\ar^{\id_V\otimes v^*}[d]&W\otimes \sigma V\ar^{v}[r]\ar_{v^*\otimes\id_{\sigma V}}[d]& Y\\
&V\otimes Y\otimes \sigma V^*& Y\otimes \sigma V^*\otimes \sigma V\ar_{\id_Y\otimes\ev_{\sigma V}}[ur]&
}
\end{equation}
可得 $S_{\coev_V\otimes\id_X, v_u}=S_{u,v}=S_{u_v,\id_Y\otimes\ev_{\sigma V}}$。
这里 $V^*$ 是 $V$ 的左对偶，$u^*: V^*\otimes X\to W$ 是 $u$ 的对偶，$v^*:W\to Y\otimes \sigma V^*$ 是 $v$ 的对偶，而
$u_v= (\id_V\otimes v^*)\circ u$，$v_u= v\circ(u^*\otimes \id_{\sigma V})$。

我们还需要说明 $\Tr'_\sigma(\mC)$ 中态射的复合。给定 $u_1: X\to V_1\otimes W_1$ 与 $v_1: W_1\otimes \sigma V_1\to Y$，再给定 $u_2: Y\to V_2\otimes W_2$ 与 $v_2: W_2\otimes \sigma V_2\to Z$，设 $S_{u_1,v_1}\in \Hom_{\Tr'_\sigma(\mC)}(\widetilde{X},\widetilde{Y})$ 与 $S_{u_2,v_2}\in \Hom_{\Tr'_\sigma(\mC)}(\widetilde{Y},\widetilde{Z})$ 是对应的元素。我们定义 $S_{u_2,v_2}\circ S_{u_1,v_1}$ 为 $S_{u,v}$，其中
\[u: X\xrightarrow{u_1} V_1\otimes W_1\xrightarrow{\id_{V_1}\otimes\coev_{V_2}\otimes\id_{W_1}}V_1\otimes V_2\otimes V^*_2\otimes W_1,\]
\[v: V^*_2\otimes W_1\otimes \sigma V_1\otimes \sigma V_2\xrightarrow{\id_{V^*_2}\otimes v_1\otimes\id_{\sigma V_2}}V^*_2\otimes Y\otimes \sigma V_2\xrightarrow{(v_2)_{u_2}} Z.\]

立即可验证该复合与 $(u_1,v_1)$ 的选取无关。另一方面，利用 \eqref{E:frac} 容易看出 $S_{u,v}=S_{u',v'}$，其中
\[
u': X\xrightarrow{(u_1)_{v_1}} V_1\otimes Y\otimes \sigma V^*_1\xrightarrow{\id_{V_1}\otimes u_2\otimes \id_{\sigma V^*_1}} V_1\otimes V_2\otimes W_2\otimes \sigma V^*_1,
\]
\[
v': W_2\otimes \sigma V^*_1 \otimes \sigma V_1\otimes \sigma V_2 \xrightarrow{\id_{W_2}\otimes\ev_{\sigma V_1}\otimes\id_{\sigma V_2}}W_2\otimes \sigma V_2\xrightarrow{v_2} Z.
\]

由这一表达式可验证该复合同样与 $(u_2,v_2)$ 的选取无关。
此外，立即可验证 $\widetilde{X}$ 的恒等态射可由典范同构 $X\cong  \mathbf{1}\otimes X$ 与 $X\otimes \sigma \mathbf{1}\cong X\otimes\mathbf{1}\cong X$ 代表。

至此我们已把 $\Tr'_\sigma(\mC)$ 定义成一个范畴。给定 $X$ 与 $Y$，有典范态射
\[\al_{X,Y}:\widetilde{X\otimes Y}\to \widetilde{Y\otimes \sigma X},\]
它由 $S_{u,v}$ 给出，其中 $u$ 与 $v$ 都取恒等映射。
特别地，对每个 $X$ 都有典范态射
$$\ga_{X}:\widetilde{X}\cong \widetilde{X\otimes\mathbf{1}}\xrightarrow{\al_{X,\mathbf{1}}} \widetilde{\mathbf{1}\otimes \sigma X}\cong \widetilde{\sigma X}.$$
由构造可验证族 $\{\al_{X,Y}\}$ 满足交换图 \eqref{E:cond trace}，从而 \eqref{E:cond trace cor} 也成立。
我们把 $\Tr_\sigma(\mC)$ 定义为 $\Tr'_\sigma(\mC)$ 关于乘法系 $\{\ga^n_X,\ X\in \mC, n\in \bZ_{\geq 0}\}$ 的局部化（于是 $\ga_X$ 在 $\Tr_\sigma(\mC)$ 中成为同构）。由此所有 $\al_{X,Y}$ 在 $\Tr_\sigma(\mC)$ 中都成为同构。
\end{proof}

\begin{rmk}
定义复合 $S_{u_2,v_2}\circ S_{u_1,v_1}$ 的想法可由下图说明
\begin{tiny}
\begin{equation*}
\label{E: comp in tr}
\xymatrix@C=0.8em@R=0.9em{
                     &                                                                                          &V_1\otimes V_2\otimes V^*_2\otimes W_1\ar@{-->}[r]&V_2\otimes V_2^*\otimes W_1\otimes \sigma V_1\ar[dr]\ar@{-->}[r]& V_2^*\otimes W_1\otimes \sigma V_1\otimes \sigma V_2\ar^-{\id_{V^*_2}\otimes v_1\otimes\id_{\sigma V_2}}[dr] &                                                    &\\
 X\ar^-{u_1}[r]&V_1\otimes W_1\ar^-{\id\otimes\coev_{V_2}\otimes\id}[ur]\ar@{-->}[r]& W_1\otimes \sigma V_1\ar[ur]\ar[r]                      &                 Y\ar[r]                                                   & V_2\otimes V_2^*\otimes Y    \ar@{-->}[r]                              & V_2^*\otimes Y\otimes \sigma V_2 \ar^-{(v_2)_{u_2}}[r]& Z. \\
}
\end{equation*}
\end{tiny}

也就是说，我们可以先把 $S_{u_2,v_2}$ 换成 $S_{\coev_{V_2}\otimes \id_Y,(v_2)_{u_2}}$。由于中间的三角形在 $\mC$ 中交换，且虚线箭头在 $\Tr_\sigma(\mC)$ 中应当是同构，$S_{\coev_{V_2}\otimes \id_Y,(v_2)_{u_2}}\circ S_{u_1,v_1}$ 必定就是证明中给出的那个元素。
\end{rmk}

\begin{rmk}
当然，还有右迹函子的对偶概念，即一个函子 $F:\mC\to \mD$ 连同函子性同构族 $\beta_{X,Y}: F(X\otimes Y)\cong F(\sigma Y\otimes X)$，满足与 \eqref{E:cond trace} 类似的条件；相应地有右 $\sigma$-扭迹范畴的概念。上述构造与下文的讨论都有对偶的版本。

为简化叙述，以下我们把左迹函子简称迹函子，把 $\mC$ 的左 $\sigma$-扭迹范畴简称 $\mC$ 的 $\sigma$-扭迹范畴。
\end{rmk}

下面是上述构造的一个具体例子。

\begin{ex}
\label{E: trace of rep}
设 $H$ 是特征零域 $E$ 上的分裂约化代数群，带有自同构 $\sigma: H\to H$。考虑由公式 \eqref{E:twist conj} 给出的 $H$ 在自身上的 $\sigma$-扭共轭，它等价于 $H$ 在旁集 $H\sigma\subset H\rtimes\langle\sigma\rangle$ 上的通常共轭。记 $\frac{H}{c_\sigma H}$ 为相应的商叠。
我们把 $H$ 的有限维表示范畴 $\on{Rep}_E(H)$ 看作分类叠 $\mathbf BH$ 上凝聚层的范畴。则拉回函子 $\pi^*: \on{Rep}_E(H)\to \Coh(\frac{H}{c_\sigma H})$ 带有 $\sigma$-扭迹函子结构，它诱导
\[\Tr_\sigma(\on{Rep}_E(H))\cong \on{Coh}_{fr}^H(H\sigma)\subset \on{Coh}(\frac{H}{c_\sigma H}),\]
其中 $\Coh_{fr}^H(H\sigma)$ 表示 $\on{Coh}(\frac{H}{c_\sigma H})$ 中由这些 $\widetilde V(:=\pi^*V)$ 张成的全子范畴。这可以这样看出。
经由下降，可以把 $\widetilde V$ 看作 $H$ 上的平凡向量丛 $H\times V$，配以对角的 $H$-等变结构。此外，对表示 $H\to \GL(V)$，记 $\sigma V$ 为 $V$ 的 $\sigma$-扭，即由 $H\xrightarrow{\sigma^{-1}}H\to\GL(V)$ 定义的 $H$ 的表示。我们形式地把恒等映射 $V\to \sigma V$ 记作 $\sigma$，于是 $h\cdot \sigma=\sigma\cdot \sigma^{-1}(h):V\to\sigma V$。
则对 $\Tr'_\sigma(\on{Rep}_E(H))$ 中的态射 $S_{u,v}$（由 $u:X\to V\otimes W$ 与 $v:W\otimes \sigma V\to Y$ 给出），我们定义 $\Coh_{fr}^H(H\sigma)$ 中的态射 $\widetilde{X}\to\widetilde{Y}$，它在 $h\in H$ 上纤维的限制为
\[
\widetilde{X}|_h=X\xrightarrow{u} V\otimes W\xrightarrow{h\sigma\otimes\id}\sigma V\otimes W\cong W\otimes \sigma V\xrightarrow{v} Y=\widetilde{Y}|_h.
\]
这给出自然函子 $\Tr_\sigma(\on{Rep}_E(H))\to \Coh_{fr}^H(H\sigma)$。按定义它是本质满的，而 Peter–Weyl 定理说明它也是全忠实的，因此是同构。

设 $\sigma=\id$。注意由于 $\on{Rep}_E(H)$ 是对称幺半的，恒等函子是迹函子，因此诱导函子 $\Tr_\sigma(\on{Rep}_E(H))\cong \Coh_{fr}^H(H\sigma)\to \on{Rep}_E(H)$。可验证它正是 $\frac{H}{cH}$ 上层的限制到 $H$ 的单位元处。
\end{ex}

我们继续关于范畴迹的一般讨论。设 $\mC$ 如命题 \ref{L:trace cat} 中那样。
先暂且假设 $\sigma$ 是恒等函子。此时把 $\Tr_\sigma(\mC)$ 记作 $\Tr(\mC)$。对 $\mC$ 中每个自同态 $f: V\to V$，记 $\coev_f:\mathbf{1}\to V\otimes V^*$ 为由 $\mathbf{1}\xrightarrow{\coev_V}V\otimes V^*\xrightarrow{f\otimes\id_{V^*}}V\otimes V^*$ 给出的态射。则有元素
\[S_{\coev_f,\ev_V}\in \End_{\Tr(\mC)}\widetilde{\mathbf{1}}.\]
特别地，当 $f=\id$ 时，记
\begin{equation}
\label{E:abs S}
S_V:=S_{\coev_{V},\ev_{V}}\in \End_{\Tr(\mC)}\widetilde{\mathbf{1}}.
\end{equation}
该元素只依赖于 $V$ 的同构类。由构造可得
\[S_{\mathbf{1}}=\id,\quad S_{V'}\cdot S_{V}=S_{V\otimes V'},\quad S_{V\oplus V'}=S_V+S_{V'}.\]
于是若 $\mC$ 是加性的，就有典范同态
\begin{equation}
\label{E: K to End}
K^{\oplus}(\mC)^{\on{op}}\to \End_{\Tr(\mC)}\widetilde{\mathbf{1}},
\end{equation}
其中 $K^{\oplus}(\mC)$ 表示 $\mC$ 的分裂 Grothendieck 环：作为 Abel 群，它由 $\mC$ 中对象生成，关系为 $[X]+[Y]=[X\oplus Y]$；环结构由张量积给出。

\begin{rmk}
注意若 $\mC$ 是对称幺半的，则 $\mC$ 的恒等函子由交换性约束给出自然的迹函子结构，从而分解为 $\mC\to \Tr(\mC)\to \mC$。显然诱导映射 $\End_{\Tr(\mC)}(\widetilde{\mathbf{1}})\to \End_\mC(\mathbf{1})$ 把 $S_{\coev_f,\ev_V}$ 送到 $f$ 的通常迹 $\tr_V(f)$，把 $S_V$ 送到 $V$ 的秩 $\on{rank}(V)$（这些概念的定义见例如 \cite[(1.7.3)]{DM}）。
\end{rmk}

上述构造有如下推广。设 $X\in \mC$，$V\in \mC$，并给定同构 $c_{X,V}: X\otimes V\simeq V\otimes X$。则由
\[X\xrightarrow{\coev_V\otimes\id_X}V\otimes V^*\otimes X, \quad V^*\otimes X\otimes V\xrightarrow{\id_{V^*}\otimes c_{X,V}} V^*\otimes V\otimes X\xrightarrow{\ev_V\otimes \id_X} X\]
得到元素 $S_{c_{X,V}}\in \End_{\Tr(\mC)}\widetilde{X}$。
若把 $X$ 在 $\mC$ 中的中心化范畴定义为这样的范畴 $\mZ_{\mC}(X)$：对象是偶 $\{V\in\mC, c_{X,V}: X\otimes V\simeq V\otimes X\}$，态射是 $\mC$ 中与 $c_{X,V}$ 和 $c_{X,V'}$ 相容的态射 $V\to V'$；则 $\mZ_{\mC}(X)$ 是幺半范畴，且当 $\mC$ 还加性时 $\mZ_{\mC}(X)$ 也加性。于是有态射
\begin{equation}
\label{E: K to End2}
K^{\oplus}(\mZ_{\mC}(X))^{\on{op}}\to \End_{\Tr(\mC)}(\widetilde X).
\end{equation}
特别地，若 $\mC$ 是加性双幺半范畴，则对每个 $X$ 都有自然幺半函子 $\mC\to \mZ_\mC(X)$。因此此时有
$$K^{\oplus}(\mC)^{\on{op}}\to \End_{\Tr(\mC)}\widetilde{X}.$$

现在处理一般的 $\sigma$。记 $\mC^\sigma$ 为 $\mC$ 中 $\sigma$-等变对象构成的范畴，即对象是偶 $(X,\phi)$，其中 $X\in\mC$，$\phi: \sigma X\simeq X$ 是同构；$(X,\phi)$ 与 $(X',\phi')$ 之间的态射是 $\mC$ 中与 $\phi$ 和 $\phi'$ 相容的态射。注意 $\mC^\sigma$ 自然地是幺半范畴，且若 $\mC$ 中每个对象都有左对偶，则 $\mC^\sigma$ 中亦然。
于是可以把上述讨论用于 $\mC^\sigma$。遗忘函子 $\mC^\sigma\to\mC$ 诱导函子
\[\Tr(\mC^\sigma)\to \Tr_\sigma(\mC).\]
特别地有
\begin{equation}
\label{E: K to End3}
K^{\oplus}(\mC^\sigma)^{\on{op}}\to \End_{\Tr_\sigma(\mC)}(\widetilde{\mathbf 1}).
\end{equation}
具体地，对 $(X,\phi: \sigma X\simeq X)\in\mC^\sigma$，$\widetilde{\mathbf 1}$ 在 $\Tr_\sigma(\mC)$ 中对应的自同态由 $S_{\coev_X, \ev_X(1\otimes\phi)}$ 给出。

下面假设 $E$ 是代数闭域，$\mC$ 是半单 Abel 范畴，且单位对象 $\mathbf{1}$ 不可约。则 $\mC^\sigma$ 是 Abel 幺半范畴。

\begin{prop}
\label{L: K iso End}
在上述假设下，映射 \eqref{E: K to End3} 诱导同构
\[
(K(\mC^\sigma)^{\on{op}}\otimes E)/I\cong \End_{\Tr_\sigma(\mC)}(\widetilde{\mathbf 1}),
\]
其中 $I$ 是 $K(\mC^\sigma)^{\on{op}}\otimes E$ 中由 $[(X,a\phi)]-a[(X,\phi)]$ 生成的理想，$(X,\phi)$ 取遍 $\mC^\sigma$ 中的不可约对象，$a\in E^\times$。
\end{prop}
结合下面的注 \ref{R: cat center}，Ostrik 有类似的结果 \cite[推论 2.16]{O}。
\begin{proof}
由于该结果后文不用，我们只给出证明梗概。首先，因 $E$ 代数闭且 $\mathbf{1}$ 不可约，可以说明：若对某个 $u:\mathbf{1}\to V\otimes W$ 与 $v: W\otimes \sigma V\to \mathbf{1}$ 有 $S_{u,v}\neq 0$，则 $S_{u,v}$ 是 $S_{\coev_{Y}, \ev_Y(1\otimes\psi)}$ 的线性组合，其中 $(Y,\psi)$ 不可分解。
其次可以说明，$\mC^\sigma$ 中的不可分解对象都形如 $(Y,\psi)=(X\otimes W, \phi\otimes f)$，其中 $(X,\phi)$ 在 $\mC^\sigma$ 中不可约，$W$ 是有限维 $E$-向量空间，$f$ 是 $W$ 的自同构。记 $f^{\on{ss}}$ 为 $f$ 的半单化，$(Y,\psi^{\on{ss}})=(X\otimes W,\phi\otimes f^{\on{ss}})$ 为 $(Y,\psi)$ 的半单化。
再说明 $S_{\coev_{Y}, \ev_Y(1\otimes\psi)}=S_{\coev_{Y},\ev_Y(1\otimes\psi^{\on{ss}})}$。由此 \eqref{E: K to End3} 诱导出所要的映射，且它是满射。单射性显然。
\end{proof}

\begin{rmk}回忆传统的范畴化/去范畴化是借助 $K$-环构造在幺半范畴与代数之间建立通道。上述命题提示了另一条经由范畴迹构造的通道。
\end{rmk}

\begin{rmk}
\label{R: cat center}
幺半范畴还有范畴中心（或称 Drinfeld 中心）的概念（见 \cite{BFO, BN, Lu3}）。给定 $(\mC,\otimes,\mathbf{1})$，其范畴中心 $\mZ(\mC)$ 是这样一个范畴：对象由 $(X\in \mC, \{\al_{X,Y}: X\otimes Y\simeq Y\otimes X\}_{Y\in \mC})$ 组成，并满足自然的相容条件。在某些情形下，有等价 $\mZ(\mC)\to \Tr(\mC)$。不过在我们的应用中，使用范畴迹的概念更为自然。
\end{rmk}

\subsection{有限 Hecke 范畴的范畴迹}

本小节中，我们在通常的代数几何（而非完全代数几何）中工作。我们研究 \S \ref{SS: cat of corr} 中所定义的范畴在如下情形下的表现。设 $G$ 是 $k$ 上的仿射代数群，带有自同态 $\sigma$；设 $K\subset G$ 是闭子群，且 $\sigma$ 限制在 $K$ 上仍为自同态。我们假设 $\sigma$ 诱导泛同胚，从而诱导 $G$ 的底层 \'etale 拓扑斯上的自同构。例如，若 $G$ 定义在有限域 $\bF_q$ 上，可取 $\sigma$ 为 $q$-Frobenius；另一个例子是 $\sigma=\id$。
把 $\sigma$-扭伴随作用记作
\[\Ad_\sigma: G\times G\to G,\quad \Ad_\sigma(g_1)(g_2)=\sigma(g_1)g_2 g_1^{-1}.\]
令
$$X=\frac{G}{\Ad_\sigma K},\quad Y=\frac{G}{\Ad_\sigma G},$$
分别为 $G$ 关于 $K$ 与 $G$ 的 $\sigma$-扭伴随作用的商叠。
注意 $f: X\to Y$ 紧合当且仅当 $K$ 是 $G$ 的抛物子群。

令 $C=X\times_YX$，并令 $\on{D}^{C}(X)$ 为 \S \ref{SS: cat of corr} 中由群胚 $C\rightrightarrows X$ 构造的范畴。
由引理 \ref{L: fully faithful}，它是 $G$ 上 $\Ad_\sigma G$-等变层范畴 $\on{D}(\frac{G}{\Ad_\sigma G})$ 的全子范畴，而后者是表示论中熟知的对象。
所以本小节并未构造新东西。不过，把 $G$ 上 $\Ad_\sigma G$-等变层表示成 $\on{D}^{C}(X)$ 中的对象，有助于我们重新理解表示论中许多已知的构造；而且在下一节我们转到仿射情形时，它将成为那里的构造的玩具模型——在仿射情形，$\on{D}^{C}(X)$ 的形式体系变得不可或缺。

作为热身，我们先用新的表述重写一些已知的结果与构造。

\begin{ex}
\label{Ex:springer}
设 $\sigma=\id$，$K=B$。由引理 \ref{L: fully faithful}，这是 $G$ 上 $\Ad G$-等变层范畴 $\on{D}(\frac{G}{\Ad G})$ 的全子范畴。

也可以考虑闭子叠 $\frac{B}{\Ad B}\subset \frac{G}{\Ad B}$ 及其上的常层 $\delta_{B}$。注意 $\frac{B}{\Ad B}\to \frac{G}{\Ad G}$ 沿 $G\to \frac{G}{\Ad G}$ 的基变换为
$\widetilde G:=G\times_{\Ad,B}B\to G$，通常称为 Grothendieck–Springer 变体。由经典的 Springer 理论，
$$\End_{\on{D}^{C}(X)}(\delta_{B})\cong E[W]$$
同构于 $G$ 的 Weyl 群 $W$ 的群代数。稍后我们会看到，这个 $W$ 在 $\delta_{B}$ 上的作用可由 \eqref{E: K to End} 应用于当前情形得到。更一般地，对每个 $w\in W$，可以把 $\frac{\overline{BwB}}{\Ad B}$ 上的相交上同调层看作 $\on{D}^{C}(X)$ 中的对象；在嵌入 $\on{D}^{C}(X)\subset \on{D}(\frac{G}{\Ad G})$ 下，它的幂等元将对应于 Lusztig 的幂幺特征层。

设 $U\subset B$ 为幂幺根。也可以考虑闭子叠 $\frac{U}{\Ad B}\subset \frac{G}{\Ad B}$ 及其上的常层 $\delta_{U}$。它对应于 $\on{D}(\frac{G}{\Ad G})$ 中的 Springer 层。
同样由经典的 Springer 理论，
$\End_{\on{D}^C(X)}(\delta_{U})\cong E[W]$。
\end{ex}

\begin{ex}
\label{Ex:DL}
设 $k=\bF_q$，$K=B$ 为 Borel 子群，$\sigma$ 为 $q$-Frobenius。由 Lang 定理，这是 $\on{D}^C(X)=\on{D}(\on{pt}/G(\bF_q))$ 的全子范畴。
考虑闭子叠 $\frac{B}{\Ad_\sigma B}\subset \frac{G}{\Ad_\sigma B}$。由 Lang 定理，对应
\[\frac{B}{\Ad_\sigma B}\leftarrow\frac{B}{\Ad_\sigma B}\times_{\frac{G}{\Ad_\sigma B}}\frac{B}{\Ad_\sigma B}\rightarrow \frac{B}{\Ad_\sigma B}\]
可等同于离散（Deligne–Mumford）叠之间的对应
\begin{equation}
\label{E:Hk for unit}
B(\bF_q)\backslash \on{pt}\leftarrow B(\bF_q)\backslash G(\bF_q)/B(\bF_q)\rightarrow \on{pt}/B(\bF_q).
\end{equation}
这里 $G(\bF_q)$ 与 $B(\bF_q)$ 看作 $\bF_q$ 上的离散代数群。
由此，对每个表示 $\rho: B(\bF_q)\to \GL(V)$，按通常的伴随构造有支撑在 $\frac{B}{\Ad_\sigma B}$ 上的局部系统 $\delta_\rho$，它进而定义 $\on{D}^C(X)$ 中的对象。若 $\rho$ 是平凡表示，我们把这个对象 $\delta_\rho$ 也记作 $\widetilde{\delta_e}$。
记 $H_W$ 为 $B(\bF_q)\backslash G(\bF_q)/B(\bF_q)$ 上 $E$ 值函数构成的 $E$-模，其代数结构由卷积（与 \eqref{E:conv prod alge} 比较）
\[(f*g)(x)=\sum_{y\in G(\bF_q)/B(\bF_q)} f(y)g(y^{-1}x)\]
给出。它通常称为附属于 $W$ 的、系数在 $E$ 中的 Iwahori–Hecke 代数。本质上由定义（以及例 \ref{Ex:examples of correspondences} (5)），有
\begin{equation}
\label{E:endo for unit}
\End(\widetilde{\delta_{e}})=H^{\on{op}}_W.
\end{equation}
类似地，若 $\rho=\on{Reg}$ 是 $B(\bF_q)$ 的正则表示，则
\begin{equation}
\label{E:endo for reg}
\End(\delta_{\on{Reg}})=E[G(\bF_q)]^{\on{op}},
\end{equation}
其中 $E[G(\bF_q)]$ 是 $G(\bF_q)$ 的群代数。

更一般地，对每个 $w\in W$，映射 $\overleftarrow{c}: \overrightarrow{c}^{-1}(\frac{BwB}{\Ad_\sigma B})\to \frac{G}{\Ad_\sigma B}$ 的纤维正是附属于 $w$ 的通常 Deligne–Lusztig 簇 $\on{DL}_w$。
此外有同构
\[  \frac{U}{\Ad_\sigma ((U \cap w^{-1}Uw)T^{w^{-1}\sigma})}\cong \frac{BwB}{\Ad_\sigma B},\quad u\mapsto u\dot{w},\]
其中 $T^{w^{-1}\sigma}=\{t\in T\mid \sigma(t)=wtw^{-1}\}$ 是有限环面，$\dot{w}$ 是 $w$ 到 $G$ 中 $T$ 的正规化子的提升。
因此，对该有限环面的任意特征标 $\theta$，有 $\frac{BwB}{\Ad_\sigma B}$ 上的局部系统 $\mL_{w,\theta}$；把它零延拓到 $\frac{G}{\Ad_\sigma B}$，就可看作 $\on{D}^{C}(X)$ 中的对象。
在全嵌入 $\on{D}^{C}(X)\subset \on{D}(\on{Rep}(G(\bF_q)))$ 下，它的像就是通常的 Deligne–Lusztig 表示。
\end{ex}

\begin{rmk}
\label{R:springer}
事实上可以验证，在上述两个例子中，若 $\cha E=0$，则 $\on{D}^C(X)$ 的幂等完备化就是 $\on{D}(Y)$。
\end{rmk}

下面我们说明 $\on{D}^{C}(X)$ 中态射的一般构造，其想法部分来自 \cite[\S 6]{XZ}，而后者推广了 V. Lafforgue 关于 $S$-算子的构造（\cite[\S 6]{La}）。我们将大量使用上同调对应的形式体系，\S \ref{Sec:cohomological correspondence} 中有简要回顾。

记 $X_r$ 为 $G^r$ 关于 $K^r$ 的作用
\begin{equation*}
(k_1,\ldots,k_r)\cdot (g_1,\ldots,g_r)=(\sigma(k_1)g_1k_2^{-1},k_2g_2k_3^{-1},\ldots,k_rg_rk_1^{-1}),\quad k_i\in K, \ g_i\in G
\end{equation*}
的商。特别地 $X_1=X=\frac{G}{\Ad_\sigma K}$。有时也把 $X_2$ 写作 $\frac{G\times^K G}{\Ad_\sigma K}$。
把第 $i$ 个与第 $(i+1)$ 个因子相乘的映射 $m_i:G^r\to G^{r-1}$ 诱导
\begin{equation}
\label{E: mult}
m_i: X_r\to X_{r-1},
\end{equation}
我们仍称之为乘法映射。
此外，映射 $(g_1,\ldots, g_r)\mapsto (\sigma(g_r),g_1,\ldots, g_{r-1})$ 诱导
\begin{equation}
\label{E: psigma}
\pFr: X_r\to X_r,
\end{equation}
我们称之为部分 $\sigma$-映射。注意 $(_\mathrm{p}\sigma)^r=\sigma$，这正是该名称的由来。
下面的等式是显然的。
\begin{equation*}
\label{E:mult vs pFr}
\pFr \cdot m_i=\left\{ \begin{array}{ll}m_{i-1}\cdot\pFr &  i>1\\ m_{r-1}\cdot (\pFr)^2 & i=1. \end{array}\right.
\end{equation*}

\begin{lem}
\label{L: decom of C}
存在典范同构
\[C=X\times_YX\cong X_2,\]
使得到第一个因子的左投影 $\overleftarrow{c}:X\times_YX\to X$ 等同于部分 $\sigma$-映射 \eqref{E: psigma} 与乘法映射 \eqref{E: mult} 的复合，
而到第二个因子的右投影 $\overrightarrow{c}:X\times_YX\to X$ 等同于乘法映射 \eqref{E: mult}。
\end{lem}
\begin{proof}
注意 $X\times_YX\cong K\backslash G\times^{K, \Ad_\sigma}G$，两个投影
$$\overleftarrow{c},\overrightarrow{c}:  K\backslash G\times^{K, \Ad_\sigma}G\to \frac{G}{\Ad_\sigma K}$$
分别由 $G\times G\to G, \ (g_1,g_2)\mapsto \sigma(g_1)g_2g_1^{-1}$ 与 $(g_1,g_2)\mapsto g_2$ 诱导。
而同构 $G\times G\cong G\times G,\ (g_1,g_2)\mapsto (g_2g_1^{-1}, g_1)$ 诱导
\[
 K\backslash G\times^{K, \Ad_\sigma}G \cong  X_2,\]
在此同构下，$\overleftarrow{c}$ 与 $\overrightarrow{c}$ 成为引理中所述的那些映射。
\end{proof}

\begin{construction}
\label{Con: S}
（V. Lafforgue 式的 $S$-算子。）
利用上述观察，我们将构造 $\on{D}^C(X)$ 中来自
\[\on{D}(K\backslash G/K)\to \on{D}(X)\to\on{D}^{C}(X)\]
的对象之间的态射，其中第一个函子是沿 $\frac{G}{\Ad_\sigma K}\to K\backslash G/K$ 的拉回。对对象 $\mF\in \on{D}(K\backslash G/K)$，记 $\widetilde\mF$ 为它在 $\on{D}(X)$（以及 $\on{D}^{X_2}(X)$）中的像。对 $\mF_1,\mF_2$，记
$\widetilde\mF_1\tilde\times\widetilde\mF_2$ 为 $\mF_1\boxtimes\mF_2$ 沿投影映射
$$\frac{G\times^KG}{\Ad_\sigma K}\to K\backslash G/K\times K\backslash G/K$$
的拉回。

考虑如下对应
\begin{equation*}
\label{E:Conv corr}
K\backslash G/K\times K\backslash G/K\xleftarrow{\pi} K\backslash G\times^KG/K\xrightarrow{m} K\backslash G/K,
\end{equation*}
其中 $m$ 由乘法 $G\times^K G\to G$ 诱导，$\pi$ 是自然投影。注意 $m$ 紧合而 $\pi$ 光滑（它是 $K$-主丛）。
与 Satake 范畴的卷积积之定义（见 \eqref{Lus conv}）类似，$\on{D}(K\backslash G/K)$ 上有卷积积
\begin{equation}
\label{Lus conv2}
\mF_1\star \mF_2=m_!\pi^*(\mF_1\boxtimes\mF_2)=m_*\pi^!(\mF_1\boxtimes\mF_2)\langle-2\dim K\rangle.
\end{equation}
配以该卷积积，$\on{D}(K\backslash G/K)$ 是幺半范畴，单位对象是支撑在 $K\backslash K/K$ 上的摩天大楼层，记作 $\delta_e$。

由伴随性，给出态射 $u: \mF\to \mG_1\star \mG_2$ 等价于给出上同调对应
\begin{equation*}
\label{E: creation corr}
m^*\mF\to\pi^!(\mG_1\boxtimes\mG_2)\langle-2\dim K\rangle,
\end{equation*}
为记号简便，它仍记作 $u$。对偶地，态射 $v:\mG_1\star\mG_2\to \mF$ 的数据等同于上同调对应
\begin{equation*}
\label{E: ann corr}
\pi^*(\mG_1\boxtimes\mG_2)\to m^!\mF,
\end{equation*}
仍记作 $v$。此外考虑态射
\[\on{sw}(\id\times \sigma): K\backslash G/K\times K\backslash G/K\to K\backslash G/K\times K\backslash G/K,\]
其中 $\on{sw}: (K\backslash G/K)^2\to (K\backslash G/K)^2$ 交换两个因子。于是有拉回上同调对应（见例 \ref{Ex:examples of correspondences} (2)）
\begin{equation}
\label{E: Frob corr}
\Gamma_{\on{sw}(\id\times\sigma)}^*: (K\backslash G/K\times K\backslash G/K, \mF_1\boxtimes \mF_2)\to (K\backslash G/K\times K\backslash G/K, \mF_2\boxtimes \sigma^*\mF_1).
\end{equation}

现在给定 $u: \mF\to \mG_1\star \mG_2$，把上同调对应的拉回形式体系 \eqref{ASS:smooth pullback correspondence}--\eqref{ASS:smooth pullback correspondence3} 用于下面的交换图（其中左方块是笛卡儿的）
\[\xymatrix{
\frac{G}{\Ad_\sigma K}\ar[d]&\ar_-m[l]\frac{G\times^KG}{\Ad_\sigma K}\ar@{=}[r]\ar[d]&\frac{G\times^KG}{\Ad_\sigma K}\ar[d]\\
K\backslash G/K&\ar_-m[l] K\backslash G\times^KG/K\ar^-\pi[r]& K\backslash G/K\times K\backslash G/K,
}\]
便得到上同调对应 $\widetilde u: m^*\widetilde{\mF}\to \widetilde{\mG_1}\tilde\times\widetilde{\mG_2}$。对偶地，态射 $v: \mG_1\star\mG_2\to \mF$ 诱导 $\widetilde v: \widetilde\mG_1\tilde\times\widetilde\mG_2\to m^!\widetilde\mF$。
此外，由下面的笛卡儿图
\[
\xymatrix{
K\backslash G/K\times K\backslash G/K\ar^-{\on{sw}(\id\times\sigma)}[rr]\ar[d]&& K\backslash G/K\times K\backslash G/K\ar[d]\\
\frac{G\times^KG}{\Ad_\sigma K}\ar^{\pFr}[rr]&& \frac{G\times^KG}{\Ad_\sigma K},
}\]
可见 \eqref{E: Frob corr} 的拉回给出 $\Ga_{\pFr}^*: (\frac{G\times^KG}{\Ad_\sigma K}, \widetilde{\mF_1}\tilde\times\widetilde{\mF_2})\to (\frac{G\times^KG}{\Ad_\sigma K}, \widetilde{\mF_2}\tilde\times \sigma^*\widetilde{\mF_1})$。

于是，若给定
\[u: \mF_1\to \mG_1\star \mG_2,\quad v: \mG_2\star \sigma^*\mG_1\to \mF_2,\]
则复合上同调对应得到
\[
\begin{aligned}
S_{u,v}:=\widetilde v \circ  \Ga_{{_p}\sigma}^* \circ \widetilde u:\ &( \frac{G}{\Ad_\sigma K},\widetilde{\mF_1})\to (\frac{G\times^KG}{\Ad_\sigma K},\widetilde{\mG_1}\tilde\times\widetilde{\mG_2})\\
&\to (\frac{G\times^KG}{\Ad_\sigma K}, \widetilde{\mG_2}\tilde\times\widetilde{\sigma^*\mG_1})\to (\frac{G}{\Ad_\sigma K},\widetilde{\mF_2}).
\end{aligned}
\]由引理 \ref{L: decom of C} 可知 $S_{u,v}\in \Hom_{\on{D}^{C}(X)}(\widetilde{\mF_1},\widetilde{\mF_2})$。这些算子称为 $S$-算子。
\end{construction}

下面是上述构造的一个直接推论。

\begin{cor}
\label{C:trace str on tilde}
函子 $\on{D}(K\backslash G/K)\to \on{D}(X)\to \on{D}^{C}(X)$ 具有自然的 $\sigma$-扭迹函子结构。因此复合 $\on{D}(K\backslash G/K)\to \on{D}(\frac{G}{\Ad_\sigma K})\to \on{D}(\frac{G}{\Ad_\sigma G})$（即所谓的极限圆变换）也具有自然的 $\sigma$-扭迹函子结构。
\end{cor}
第二个论断已见于 \cite{BFO,BN,Lu3}。
\begin{proof}
存在典范同构 $\widetilde{\mF_1\star\mF_2}\cong\widetilde{\mF_2\star \sigma^*\mF_1}$，它由
\[u=\id: \mF_1\star\mF_2=\mF_1\star\mF_2,\quad v=\id: \mF_2\star \sigma^*\mF_1=\mF_2\star \sigma^*\mF_1\]
诱导。
\end{proof}

若 $\mF\in \on{D}_c^b(K\backslash G/K)$，则 $\mF$ 关于由卷积积 \eqref{Lus conv2} 给出的 $\on{D}(K\backslash G/K)$ 上的幺半结构是可对偶化的。事实上，令
\[\mF^*=\iota^*\bD(\mF)\langle2\dim K\rangle,\]
其中 $\bD$ 是 $\on{D}_c^b(K\backslash G/K)$ 上的 Verdier 对偶，
$$\iota: K\backslash G/K\to K\backslash G/K,\quad  g\mapsto g^{-1}.$$
则 $\mF^*$ 同时是 $\mF$ 的左对偶与右对偶。
具体地说，利用投影 $\on{pt}\to \mathbf BK=K\backslash K/K$，可以把例 \ref{Ex:examples of correspondences} (4) 中的上同调对应看作上同调对应
\[\coev'_{\mF}\in \on{Corr}_{K\backslash G/K}((K\backslash K/K, E), (K\backslash G/K\times K\backslash G/K, \mF\boxtimes \bD \mF)),\]
\[\ev'_{\mF}\in\on{Corr}_{K\backslash G/K}((K\backslash G/K\times K\backslash G/K, \bD\mF\langle2\dim K\rangle\boxtimes \mF), (K\backslash K/K,E)).\]
有下面的交换图
\[
\xymatrix{
K\backslash K/K\ar[d]&\ar[l] K\backslash G/K\ar^-\Delta[r]\ar^{\Delta'}[d]& K\backslash G/K\times K\backslash G/K\ar^{\id\times\iota}[d]\\
K\backslash G/K&\ar[l] K\backslash G\times^KG/K\ar[r]& K\backslash G/K\times K\backslash G/K,
}\]
其中 $\Delta'$ 由 $G\to G\times^KG,\ g\mapsto (g,g^{-1})$ 诱导，使得左边的交换方块是笛卡儿的。注意右边的交换方块按定义 \ref{D: basechangeable} 也是可基变换的。
于是按 \eqref{ASS:pushforward correspondence} 的意义把 $\coev'_{\mF}$ 与 $\ev'_\mF$ 沿 $\Delta'$ 前推，得到
\begin{equation*}
\label{E:unandcoun}
\coev_{\mF}: \delta_e\to \mF\star \mF^*,\quad \ev_{\mF}:\mF^*\star\mF\to \delta_e,
\end{equation*}
使 $\mF^*$ 成为 $\mF$ 的左对偶。类似地，有态射 $\coev'_{\mF}:\delta_e\to \mF^*\star \mF$ 与 $\ev'_{\mF}:\mF\star\mF'\to \delta_e$，使 $\mF^*$ 成为 $\mF$ 的右对偶。
由该构造得到特别元素
\begin{equation}
\label{E: SLa}
S_{\mF}:=S_{\coev_\mF,\ev_\mF}\in \End_{\on{D}^{X_2}(X)}(\widetilde{\delta_e}).
\end{equation}
这正是 V. Lafforgue 原来构造的 $S$-算子。

由命题 \ref{L:trace cat} 与引理 \ref{C:trace str on tilde} 可知，函子 $\on{D}_c^b(K\backslash G/K)\to \on{D}^{C}(X)$ 经过
\[
\Tr_\sigma(\on{D}_c^b(K\backslash G/K))\to \on{D}^{C}(X)
\]
分解。

命题 \ref{L:trace cat} 的证明中定义的元素 $S_{u,v}$ 恰好映到构造 \ref{Con: S} 中构造的 $S$-算子。特别地，元素 \eqref{E:abs S} 映到 \eqref{E: SLa}。

\begin{ex}
现在把上述一般讨论具体到几个特殊情形。
我们考虑 $G$ 约化、$K=B$ 的情形。回忆 $\on{D}_c^b(B\backslash G/B)$ 通常称为有限 Hecke 范畴。
若 $E=\Ql$ 且 $k$ 是有限域，则对 $w\in W$，记 $\IC_w$ 为 $B\backslash BwB/B$ 上局部系统 $\Ql[\ell(w)](\ell(w)/2)$ 到其闭包的中间延拓，其中 $\ell(w)=\dim BwB/B$ 是 $W$ 上通常的长度函数。（这里我们固定了半 Tate 扭转。）
于是 $\on{D}_c^b(B\backslash G/B)$ 中由 $\{\IC_w[d](\frac{d}{2}),\ w\in W, d\in \bZ\}$ 张成的子范畴是半单幺半子范畴。它的 Grothendieck 环同构于 $\bZ[v,v^{-1}]$ 上的有限 Hecke 代数 $\mathbb H_W$（其中 $\bZ[v,v^{-1}]\subset \mathbb H_W$ 的嵌入由 $\{\on{IC}_0(i), \ i\in\bZ\}$ 到 Hecke 范畴的嵌入诱导）。当 $v=1$ 时，$\mathbb H_W$ 特化为 $W$ 的群代数；当 $v=\sqrt{q}$ 时，它特化为前面提到的 Hecke 代数 $H_W$。
由引理 \ref{L: K iso End} 可知
\begin{itemize}
\item 在例 \ref{Ex:springer} 的情形，$\delta_B$ 上有 $\bH^{\on{op}}_W$ 的自然作用。但显然 $S_{\IC_w(i)}=S_{\IC_w}$。因此 $\bH^{\on{op}}_W$ 的作用经过 $W$ 在 $\delta_B$ 上的作用分解。事实上，我们的构造不过是把 Ginzburg 用对应构造 Springer 作用的做法重新表述了一遍。
\item 在例 \ref{Ex:DL} 的情形，映射 \eqref{E: K to End3} 诱导 $H_W^{\on{op}}$ 在 $\widetilde{\delta_e}$ 上的典范作用。另一方面，
由 \eqref{E:Hk for unit} 有 $\End\widetilde{\delta_e}=H_W^{\on{op}}$。于是得到态射 $H^{\on{op}}_W\to H^{\on{op}}_W$。下面的命题 \ref{P:S=T finite} 说明该映射是恒等映射。
\end{itemize}
\end{ex}
\begin{rmk}
映射 \eqref{E: K to End2}（及其 $\sigma$-扭版本）应当使人能够推断 $\on{D}^C(X)$ 中更一般对象（例如 Deligne–Lusztig 表示）的自同态环的信息。
\end{rmk}

设 $\mF\in \on{D}_c^b(B\backslash G/B)$，并设它带有典范同构 $\sigma^*\mF\cong \mF$（Weil 结构）。记 $f_\mF$ 为经由 Grothendieck 的层-函数字典得到的 $B(\bF_q)\backslash G(\bF_q)/B(\bF_q)$ 上的函数。由 \eqref{E:Hk for unit}，该函数给出 $\End(\widetilde{\delta_e})$ 中的元素，记作 $S'_\mF$。
\begin{prop}
\label{P:S=T finite}
在映射 $\End_{\Tr_{\sigma}(\on{D}^b_c(B\backslash G/B))}(\widetilde{\mathbf{1}})\to \End (\widetilde{\delta}_e)$ 下，元素 $S_\mF$ 映到 $S'_{\mF}$。
\end{prop}
\begin{proof}
我们给出该命题的详细证明。
\cite[\S 6.3]{XZ} 中用了一个类似但更复杂的论证来建立它的仿射版本（该版本将在下文作为定理 \ref{T: trace of geom Sat} 的一部分陈述）。

记 $\mZ=m^{-1}(\frac{B}{\Ad_\sigma B})\subset \frac{G\times^BG}{\Ad_\sigma B}$。考虑下图
\begin{small}
\begin{equation}\label{E: cal coh corr to local}
\xymatrix@C=15pt{
\on{pt}&\ar[l] G/B\ar^-{\Delta}[r]& G/B\times G/B\ar@{}[drr]|{\mathrm{(A)}}&& G/B\times G/B\ar_-{\on{sw}(\id\times \sigma)}[ll]\ar@{}[dr]|{\mathrm{(B)}}&\ar_-\Delta[l] G/B\ar[r]& \on{pt}\\
\on{pt}\ar[d]\ar@{=}[u]&\ar[l] \mZ_1\ar[r]\ar[d]\ar[u]& G\times G\ar[d]\ar[u]\ar@{}[drr]|{\mathrm{(C)}}&&\ar_-{\on{sw}(\id\times \sigma)}[ll]G\times G\ar[d]\ar[u]\ar@{}[dr]|{\mathrm{(D)}}&\ar[l]\mZ_2\ar[r]\ar[d]\ar[u]&\on{pt}\ar[d]\ar@{=}[u]\\
\frac{B}{\Ad_\sigma B} & \mZ\ar[l]\ar[r]& \frac{G\times^BG}{\Ad_\sigma B}&&\ar_{\pFr}[ll]\frac{G\times^BG}{\Ad_\sigma B}&\ar[l] \mZ\ar[r]& \frac{B}{\Ad_\sigma B},
}
\end{equation}
\end{small}
其中
\begin{itemize}
\item 方块 $(C)$ 中两个竖直映射是自然投影，映射 $G\times G\to G\times G$ 由 $(g_1, g_2)\mapsto (\sigma(g_2), g_1)$ 给出，使得 $(C)$ 是笛卡儿的；
\item 簇 $\mZ_2$ 的定义使得图 $(D)$ 是笛卡儿的，即 $\mZ_2=\{(g_1,g_2)\in G\times G,\ g_1g_2\in B\}$；
\item 图 $(A)$ 中左边的竖直映射 $G\times G\to G/B\times G/B$ 由 $(g_1,g_2)\mapsto (g_1B, g_2^{-1}B)$ 给出，右边竖直映射由 $(g_1,g_2)\mapsto (g_1^{-1}B,g_2B)$ 给出，使得 $(A)$ 与 $(B)$ 都是笛卡儿的；
\item 簇 $\mZ_1=G$，映射 $\mZ_1\to G\times G$ 由 $g\mapsto (g,g^{-1})$ 给出，使得复合 $\mZ_1\to G\times G\to \frac{G\times^BG}{\Ad_\sigma B}$ 经过 $\mZ$ 分解，且所得态射 $\mZ_1 \to \mZ$ 光滑；而映射 $\mZ_1\to G/B$ 是自然投影（因而也光滑）。
\end{itemize}
把每一行取纤维积，得到交换图
\begin{equation*}
\xymatrix{
\on{pt}&\ar[l] (G/B)(\bF_q)\ar[r] & \on{pt}\\
\on{pt}\ar@{=}[u]\ar[d]&\ar[l] \mW \ar[u]\ar[d]\ar[r] & \on{pt}\ar@{=}[u]\ar[d]\\
B(\bF_q)\backslash \on{pt}&\ar[l] B(\bF_q)\backslash G(\bF_q)/B(\bF_q)\ar[r]& \on{pt}/B(\bF_q),
}\end{equation*}
其中 $\mW:=\mZ_1\times_{G\times G, \on{sw}(\id\times \sigma)}\mZ_2\cong\{g\in G\mid \sigma(g)^{-1}g\in B\}$；它通过 $g\mapsto\sigma(g)$ 映入 $\mZ_1=G$，通过 $g\mapsto(\sigma(g)^{-1},g)$ 映入 $\mZ_2$。映射 $\mW\to (G/B)(\bF_q)$ 把 $g$ 送到 $gB$，故其纤维都同构于 $B$。映射 $\mW\to B(\bF_q)\backslash G(\bF_q)/B(\bF_q)$ 把 $g$ 送到它在 $B(\bF_q)\backslash (G/B)$ 中的像（该像自然落在 $B(\bF_q)\backslash G(\bF_q)/B(\bF_q)$ 中）。

按定义，$S_{\mF}: (B(\bF_q)\backslash \on{pt},\widetilde{\delta_e})\to (\on{pt}/B(\bF_q),\widetilde{\delta_e})$ 是支撑在上图底行的上同调对应。另一方面，有支撑在顶行的上同调对应
$$S'_\mF:=\on{ev}_{\mF}\circ \Ga_{\on{sw}(\id\times\sigma)}^* \circ\on{coev}_{\mF}: (\on{pt},E)\to(\on{pt},E),$$
按 Braverman–Varshavsky 的迹公式（见 \cite[引理 A.2.22]{XZ}），它由 $B(\bF_q)\backslash G(\bF_q)/B(\bF_q)$ 上的函数 $f_\mF$ 给出。

由于 $\mW\to (G/B)(\bF_q)$ 与 $\mW\to B(\bF_q)\backslash G(\bF_q)/B(\bF_q)$ 都是等维光滑的（维数为 $\dim B$），可以用 \eqref{ASS:smooth pullback correspondence}--\eqref{ASS:smooth pullback correspondence3} 中的形式体系把 $S_\mF$ 与 $S'_\mF$ 拉回到中间一行。只须证明它们的拉回相等。

注意 $S_{\mF}$ 与 $S'_{\mF}$ 都是三个上同调对应的复合。
由于交换方块 (A)–(D) 都是笛卡儿的，引理 \ref{AL:pullback compatible with composition} 说明光滑拉回与上同调对应的复合可交换。因此只剩验证在图 \eqref{E: cal coh corr to local} 中
\begin{itemize}
\item 支撑在 $\mZ$ 上的上同调对应 $\widetilde{\coev_\mF}: (\frac{B}{\Ad_\sigma B}, \widetilde{\delta_e})\to (\frac{G\times^BG}{\Ad_\sigma B}, \widetilde{\mF}\tilde\times \widetilde{\mF^*})$ 到中间一行的拉回，等于支撑在 $G/B$ 上的 $\on{un}_{\mF}: (\on{pt},E)\to (G/B\times G/B, \mF\boxtimes \bD\mF)$ 的拉回；关于 $\ev_\mF$ 与 $\on{coun}_{\mF}$ 的对偶论断同样成立；
\item 对应 $\Ga_{\pFr}^*: (\frac{G\times^BG}{\Ad_\sigma B}, \widetilde{\mF}\tilde\times\widetilde{\mF^*})\to (\frac{G\times^BG}{\Ad_\sigma B}, \widetilde{\mF^*}\tilde\times\widetilde{\mF})$ 到中间一行的拉回，等于 $\Ga_{\on{sw}(\id\times\sigma)}^*: (G/B\times G/B, \mF\boxtimes \bD \mF)\to (G/B\times G/B, \bD \mF\boxtimes \mF)$ 的拉回。
\end{itemize}
\end{proof}

\section{几何 Satake 的范畴迹}
\label{S: cattrSat}

本节转到仿射情形。基本想法依旧，但技术细节更繁复。

设 $F$ 是剩余域 $k=\bF_q$ 有限的局部域，$q$-Frobenius 记作 $\sigma$。
回忆 \S \ref{S: affGrass} 中的记号，特别是 \eqref{ramified Witt vector} 与 \eqref{discs}。我们假设 $G$ 是 $\mO$ 上的连通约化群。为简便，记 $K=L^+G$。

由上一节的讨论可见，要研究 $p$ 进群的表示论，重要的是研究 $LG$ 上 $\Ad LG$-等变层或 $\Ad_\sigma LG$-等变层的范畴。然而仿射情形的困难之一在于：商 $Y=\frac{LG}{\Ad LG}$ 或 $Y=\frac{LG}{\Ad_\sigma LG}$ 远非代数叠，因此如何理解其上的层范畴并不清楚。我们的办法（应当等价于 \cite{Ga} 中概述的办法）是用 \S \ref{SS: cat of corr} 中引入的范畴 $\on{D}^\Hk(X)$ 代替 $\on{D}(Y)$，其中 $X=\frac{LG}{\Ad K}$ 或 $X=\frac{LG}{\Ad_\sigma  K}$，而 $\Hk=\Hk(X):=X\times_YX$\footnote{更确切地说，需要按注 \ref{R: DC enh} (2) 所述把 $\on{D}(Y)$ 定义为 $\Shv^{\Hk}(X)$ 的同伦范畴。但本文不需要这个更精细的版本。}。我们将会看到，虽然 $X$ 与 $\Hk$ 都不是代数叠，它们可以被有限维代数叠逼近，因此 $\on{D}^{\Hk}(X)$ 有意义。用 $\on{D}^\Hk(X)$ 代替 $\on{D}(Y)$ 的另一个好处是：前者经由上同调对应实现，与整体模空间（Shimura 簇或整体 Shtuka 的模空间）相关。我们还要指出，在 Fargues–Scholze 处理 $p$ 进群局部 Langlands 对应的办法中，他们用 Fargues–Fontaine 曲线上 $G$-丛的模空间上的层范畴替代所寻求的范畴 $\on{D}(\frac{LG}{\Ad_\sigma LG})$。它应当经由近循环函子与我们的范畴 $\on{D}^\Hk(X)$ 相关。不过，作者并不清楚如何在 Fargues–Scholze 的办法中触及 $\on{D}(\frac{LG}{\Ad LG})$。

一旦定义了合适的范畴，我们就可以取几何 Satake 对应的范畴迹，这将是下一节算术应用的关键要素。本节只讨论 $Y=\frac{LG}{\Ad_\sigma LG}$ 的情形；不过几乎所有构造也适用于 $Y=\frac{LG}{\Ad LG}$ 的情形，见注 \ref{R:affine springer}。

\subsection{局部 Shtuka 的模空间}

由于
$X=\frac{LG}{\Ad_\sigma K}$、$Y=\frac{LG}{\Ad_\sigma LG}$ 与 $X\times_YX$ 都不是代数叠，要把先前的讨论推广到仿射情形，第一步是使 \S \ref{SS: cat of corr} 中定义的范畴在此情形下有确切含义。结果表明，此时的 $\frac{LG}{\Ad_\sigma K}$ 可以解释为局部 Shtuka 的模空间。

设 $R$ 是完美 $k$-代数。若 $\mE$ 是 $D_R$ 上的 $G$-主丛，则它沿 $q$-Frobenius $\sigma:D_R\to D_R$ 的拉回记作 ${^\sigma}\mE$。
$\Spec R$ 上的局部 $r$ 重 $G$-Shtuka 是指 $D_{R}$ 上一列 $G$-主丛 $\mE_1,\ldots, \mE_r$，连同一串修正
\[\mE_r\dashrightarrow \mE_{r-1}\dashrightarrow\cdots\dashrightarrow \mE_0:={^\sigma}\mE_r.\]
若
\[\inv(\mE_i\dashrightarrow \mE_{i-1})\leq \mu_i,\]
则称其奇点被 $\mmu=(\mu_1,\ldots,\mu_r)$ 所界。我们把局部 $r$ 重 $G$-Shtuka 的模空间 $\on{Sht}^{r,\loc}$ 定义为如下预叠：它对每个 $R$ 赋予 $\Spec R$ 上局部 $r$ 重 $G$-Shtuka 的群胚。它是闭子预叠 $\on{Sht}^{\loc}_{\mmu}$ 之并（$\mmu$ 取遍），后者由奇点被 $\mmu=(\mu_1,\ldots,\mu_r)$ 所界的局部 Shtuka 组成。特别地，当 $r=1$ 时，$\on{Sht}^{r,\loc}$ 记作 $\Sht^\loc$，称为局部 $G$-Shtuka 的模空间，它是闭子预叠 $\Sht^\loc_\mu$ 之并。

由于 $W(R)$ 上每个 $G$-主丛在 $R$ 上 \'etale 局部可平凡化（例如见 \cite[引理 1.3]{Zh1}），有自然同构
\begin{equation}
\label{R:pre of Sht}
\Sht^\loc\cong \frac{LG}{\Ad_\sigma K},\quad \mbox{它限制为同构 } \Sht^\loc_\mu\cong \frac{\Gr_\mu^{(\infty)}}{\Ad_\sigma K},\ \ \ \mu\in \xcoch(T)/W.
\end{equation}
具体地说，记 $\Sht^{\loc,\Box}$ 为 $\Sht^\loc$ 上的 $K$-主丛，它对局部 $G$-Shtuka $(\beta:\mE\dashrightarrow {^\sigma}\mE)$ 连同平凡化 $\epsilon: \mE\simeq \mE^0$ 分类。则复合 $\mE^0\xrightarrow{\epsilon^{-1}}\mE\stackrel{\beta}{\dashrightarrow}{^\sigma}\mE\xrightarrow{\sigma(\epsilon)}\mE^0$ 定义元素 $g\in LG$，并诱导同构 $\Sht^{\loc,\Box}\cong LG$；在该同构下，$K$ 在 $\Sht^{\loc,\Box}$ 上的作用等同于 $\Ad_\sigma$ 在 $LG$ 上的作用。此外，若 $\inv(\beta)\leq \mu$，则 $g\in \Gr_{\mu}^{(\infty)}$。取商即得上述同构。

\begin{ex}\label{no singularity}
若 $\mmu=\mu=0$，则 $\Sht^{\loc}_{0}=\mathbf B G(\mO)$ 是投射有限群 $G(\mO)$ 的分类叠。
\end{ex}

\begin{ex}\label{loc Sht and pdiv}
若 $\mO=\bZ_p$，则 $R$ 上奇点被余权 $\omega_i$（如例 \ref{minu Sch} 中所定义）所界的 $\GL_n$-Shtuka 范畴，等价于 $R$ 上维数 $i$、高度 $n$ 的 $p$-可除群的范畴。该函子把 $p$-可除群 $\bX$ 送到它的 Dieudonné 模，而 Gabber 的一个定理断言这是等价。注意来自 $p$-可除群的 $\GL_n$-Shtuka 非常特殊：它们的奇点必须被极小余权所界。
\end{ex}

\begin{rmk}
我们刚才定义的是奇点在闭点 $s\in D$ 处的局部 Shtuka。也可以定义奇点在一般点 $\eta\in D$ 处的局部 Shtuka，甚至让奇点沿 $D$ 移动。
在混合特征情形，奇点沿 $D$ 的局部 Shtuka 与 Breuil–Kisin 模密切相关（\cite{SW}）。
\end{rmk}

\begin{rmk}
在等特征情形，可以定义整体 $G$-Shtuka，其中 $D$ 被有限域上的整体代数曲线替代。这正是 Drinfeld 最初的发明，它极大地推广了椭圆模的概念。不过目前尚不清楚数域情形下是否存在整体 $G$-Shtuka 的类似物。一些猜测可参见 \cite{S}。
\end{rmk}

同构 \eqref{R:pre of Sht} 显然推广为同构
$$\frac{LG\times^KLG\times^K\cdots\times^KLG}{\Ad_\sigma K}\cong \on{Sht}^{r,\loc},$$
在此同构下，与 \eqref{E: psigma} 类似的映射也有模空间的解释，通常称为\emph{部分 Frobenius 映射}。
固定如上的一列支配余权 $\mmu=(\mu_1,\ldots,\mu_r)$，并令 $\mu_0=\sigma(\mu_r)$。则部分 Frobenius 映射定义为
\begin{equation*}
\label{hpf}
\xymatrix@R=0pt{
F_\mmu\colon \on{Sht}_{(\mu_1,\ldots,\mu_r)}^{r,\loc} \ar[r] & \on{Sht}^{r,\loc}_{(\sigma(\mu_{r}),\mu_1,\ldots,\mu_{r-1})} \\
(\mE_r\stackrel{\beta_r}{\dashrightarrow} \mE_{r-1}\stackrel{\beta_{r-1}}{\dashrightarrow}\cdots\stackrel{\beta_2}{\dashrightarrow} \mE_1\stackrel{\beta_1}{\dashrightarrow} {^\sigma}\mE_r) \ar@{|->}[r] & (\mE_{r-1}\stackrel{\beta_2}{\dashrightarrow}\cdots\stackrel{\beta_r}{\dashrightarrow} \mE_1\stackrel{\beta_1}{\dashrightarrow} {^\sigma}\mE_r\stackrel{\sigma(\beta_r)}{\dashrightarrow} {^\sigma}\mE_{r-1}),
}
\end{equation*}
它是预叠的同构（即完美 $k$-代数上函子的同构）。

我们定义如下预叠的对应
\begin{equation}
\label{E;Hk Shtuka}
\xymatrix{&\Hk^{s,t}(\Sht^\loc)\ar_{\overleftarrow{h}^\loc}[dl]\ar^{\overrightarrow{h}^\loc}[dr]&\\
\Sht^{s,\loc}&&\Sht^{t,\loc},
}
\end{equation}
其中 $\Hk^{s,t}(\Sht^\loc)$ 是如下预叠：它对每个完美 $k$-代数 $R$，分类 $D_R$ 上 $G$-主丛修正的如下交换图
\begin{equation*}\label{qi for diff type}
\xymatrix@C=35pt{
\mE'_s \ar@{-->}[r]^{\beta'_s} \ar@{-->}[d]_\beta & \cdots  \ar@{-->}[r]^{\beta'_2} & \calE'_1 \ar@{-->}[r]^{\beta'_1} & {}^\sigma\mE'_s \ar@{-->}[d]^{\sigma(\beta)}\\
\mE_t \ar@{-->}[r]^{\beta_t} & \cdots  \ar@{-->}[r]^{\beta_2} & \calE_1 \ar@{-->}[r]^{\beta_1} & {}^\sigma\mE_t,
}
\end{equation*}
使得顶行（相应地底行）给出 $\Sht^{s,\loc}$（相应地 $\Sht^{t,\loc}$）的 $R$-点。
我们把这样的图称为从顶行到底行的\emph{Hecke 对应}。当 $s=t=1$ 时，把 $\Hk^{s,t}(\Sht^\loc)$ 简记作 $\Hk(\Sht^\loc)$。
于是与 \eqref{R:pre of Sht} 类似，有典范同构
$$\Hk(\Sht^\loc)\cong \frac{LG}{\Ad_\sigma K}\times_{\frac{LG}{\Ad_\sigma LG}} \frac{LG}{\Ad_\sigma K},$$
其中 $\overleftarrow{h}^\loc$ 与 $\overrightarrow{h}^\loc$ 是自然投影。
对两列支配余权 $\la_\bullet$ 与 $\mmu$，令
\[\Sht^{\loc}_{\la_\bullet\mid\mmu}:=(\overleftarrow{h}^\loc)^{-1}(\Sht^\loc_{\la_\bullet})\cap (\overrightarrow{h}^\loc)^{-1}(\Sht^\loc_{\mmu}),\]
有时把 $\overleftarrow{h}^\loc$ 与 $\overrightarrow{h}^\loc$ 在 $\Sht^{\loc}_{\la_\bullet\mid\mmu}$ 上的限制记作 $\overleftarrow{h}^\loc_{\la_\bullet}$ 与 $\overrightarrow{h}^\loc_\mmu$。
若进一步要求 $\beta$ 的相对位置 $\leq \nu$，就得到 $\Sht_{\la_\bullet|\mmu}^{\loc}$ 的闭子预叠 $\Sht_{\la_\bullet|\mmu}^{\nu,\loc}$。

\begin{rmk}
\label{R:fiber of hloc}
注意态射 $\overleftarrow{h}^\loc: \Hk(\Sht^\loc)\to\Sht^\loc$ 的纤维正是仿射 Grassmann 簇。对每个 $\mu$，态射
$\overleftarrow{h}^\loc: (\overrightarrow{h}^\loc)^{-1}(\Sht^{\loc}_\mu)\to \Sht^\loc$ 的纤维是仿射 Deligne–Lusztig 簇（定义可参见 \cite[\S 3.1]{Zh1}）。
\end{rmk}

由例 \ref{no singularity} 已经可见，$\Sht^\loc_\mu$ 不是代数叠，因此 \S \ref{SS: cat of corr} 中的一般构造不能直接用于当前情形。为了在此情形定义范畴，我们需要用有限维逼近替代这些模空间。
这里只勾勒构造，细节参见 \cite[\S 5]{XZ}。构造的群论描述见注 \ref{R: grp theoretical Sht}。

回忆 $K$ 作用在卷积 Grassmann 簇上。给定支配余权序列 $\mmu=(\mu_1,\ldots,\mu_r)$，称整数 $m$ 为 $\mmu$-大，若它 $\geq \langle\sum_i\mu_i,\al_h\rangle$，其中 $\al_h$ 是 $G$ 的最高根。
容易看出，给定非负整数对 $(m,n)$（允许 $(m,n)=(\infty,\infty)$），若 $m-n$ 是 $\mmu$-大，则 $K$ 在 $\Gr_\mmu^{(n)}$ 上的作用经过 $K_{m}$ 的作用分解（参看 \cite[引理 3.1.7]{XZ}）。此时 fpqc 商叠
$\Hk^{\loc(m)}_\mmu:=[K_m\backslash \Gr_\mmu]$ 称为 $m$-限制局部 Hecke 叠。
由于
\[
[K_m\backslash\Gr_\mmu]\cong [K_m\backslash\Gr^{(n)}_\mmu/K_n]\xrightarrow{t_\leftone\times t_\rightone} \bfB K_n\times \bfB K_m,
\]
我们看到，在 $\Hk_{\mmu}^{\loc(m)}$ 之上，第一个映射 $t_\rightone$ 定义 $K_m$-主丛
\begin{equation*}\label{torsor2}
\Gr_\mmu \to [K_m\backslash\Gr_\mmu],
\end{equation*}
第二个映射 $t_\leftone$ 定义 $K_n$-主丛
\begin{equation*}\label{torsor1}
[K_m\backslash\Gr_{\mmu}^{(n)}]\to  [K_m\backslash\Gr_\mmu].
\end{equation*}
把这两个\emph{典范}主丛分别记作 $\calE_{\rightone}|_{D_{m}}$ 与 $\calE_{\leftone}|_{D_n}$。

对 $m\geq n$，记
$$\res^m_n: \bfB K_m\to \bfB K_n$$
为由商映射 $\pi_{m,n}:K_m\to K_n$ 诱导的映射。显然，当 $m\leq m'\leq m''$ 时
\begin{equation}
\label{E:ex and q compatibility}
\res_{m}^{m'} \circ \res_{m'}^{m''} = \res_{m}^{m''},
\end{equation}
这里允许 $m''=\infty$。
于是对 $m\leq m'$、$n\leq n'$（允许 $(m',n')=(\infty,\infty)$），若 $m'-n'$ 与 $m-n$ 都是 $\mmu$-大，则有
限制映射
$$\res^{m'}_m: \Hk^{\loc(m')}_{\mmu}\to \Hk^{\loc(m)}_{\mmu},$$ 使得下图交换且右方块是笛卡儿的。
\begin{equation*}
\label{E:restriction and torsor commutative}
\xymatrix{
\bfB K_{n'} \ar[d]^{\res^{n'}_n} &
\Hk_\mmu^{\loc(m')} \ar[l]_-{t_\leftone} \ar[r]^-{t_\rightone} \ar[d]^{\res^{m'}_{m}} &  \bfB K_{m'} \ar[d]^{\res^{m'}_{m}}
\\
\bfB K_n & \Hk_\mmu^{\loc(m)} \ar[l]_-{t_\leftone} \ar[r]^-{t_\rightone} &  \bfB K_m.
}
\end{equation*}
此外，这些限制映射满足自然的相容性条件 \eqref{E:ex and q compatibility}。

现在固定支配余权序列 $\mmu$，并选整数对 $(m,n)$ 使 $m-n$ 为 $\mmu$-大。
对完美 $k$-代数 $R$，$\Spec R$ 上奇点被 $\mmu$ 所界的 $(m,n)$-限制局部迭代 $G$-Shtuka 由以下数据组成：
\begin{itemize}
\item $ \Hk^{\loc(m)}_\mmu$ 的一个 $R$-点，以及
\item $\Spec R$ 上 $K_n$-主丛的同构
$$\psi:{^\sigma} (\mE_\leftone|_{D_{n}}) \simeq(\mE_\rightone|_{D_{m}})|_{D_{n}}.$$
\end{itemize}
把分类所有 $(m,n)$-限制局部 $G$-Shtuka 的预叠定义为 $\Sht^{\loc(m,n)}_\mmu$。换言之，$\Sht_\mmu^{\loc(m,n)}$ 定义为 $\Hk_\mmu^{\loc(m)}$ 与一个 Frobenius 图的纤维积
\begin{equation}
\label{E:mS restricted as a Cartesian product}
\xymatrix@C=90pt{
\Sht_\mmu^{\loc(m,n)} \ar[r]^{\varphi^{\loc(m,n)}} \ar[d]
&  \Hk_\mmu^{\loc(m)} \ar[d]^-{ t_\leftone\times \res^{m}_{n}\circ t_\rightone }
\\ \bfB K_n \ar[r]^{1\times \sigma} & \bfB K_n \times \bfB K_n,
}
\end{equation}
特别地，$\Sht_\mmu^{\loc(m,0)}=\Hk_\mmu^{\loc(m)}$，且 $\Sht_\mmu^{\loc(\infty,\infty)}=\Sht_\mmu^\loc$。
注意 $\varphi^{\loc(m,n)}$ 是完全光滑的态射，相对维数为 $n \dim G$。

设 $(m',n')$ 与 $(m,n)$ 是两对非负整数，满足 $m\leq m'$、$n\leq n'$，且 $m'-n'$ 与 $m-n$ 都是 $\mmu$-大（允许 $(m',n')=(\infty,\infty)$）。
则有\emph{限制态射}
\begin{equation}\label{trunc}
\res^{m',n'}_{m,n}: \Sht^{\loc(m',n')}_\mmu\longto \Sht_\mmu^{\loc(m,n)},
\end{equation}
它是相对维数 $(m-n)-(m'-n')$ 的完全光滑态射，与 \eqref{E:mS restricted as a Cartesian product} 相容，并满足
\begin{equation}
\label{true-to-trun}
\res^{m_2,n_2}_{m_1,n_1}\circ\res^{m_3,n_3}_{m_2,n_2}=\res^{m_3,n_3}_{m_1,n_1}.
\end{equation}

\begin{ex}
\label{Ex:mStau mn}
当 $\mu=0$，更一般地当 $\mu=\tau\in\xcoch(Z_G)$ 时，有
$$\Sht^{\loc(n,n)}_\tau\cong \bfB G(\mO/\varpi^n).$$
限制映射 $\res_{n,n}: \Sht_\tau^\loc\cong\bfB G(\mO)\to \Sht_\tau^{\loc(m,n)}\cong\bfB G(\mO/\varpi^n)$ 由自然映射 $G(\mO)\to G(\mO/\varpi^n)$ 诱导。
\end{ex}

\begin{rmk}
\label{Ex:Sht and GZip}
如例 \ref{loc Sht and pdiv} 所述，局部 Shtuka 的模空间可以看作带 $G$-结构的 $p$-可除群模空间的推广。自然期待：限制局部 Shtuka 的模空间可以看作带 $G$-结构的截断 Barsotti–Tate 群模空间的推广。结果表明两者确实密切相关，但关系比非限制情形略复杂。如 \cite[引理 5.3.6]{XZ} 所述，若 $\mu$ 极小且 $(m,n)=(2,1)$，则有自然完全光滑映射
\begin{equation}
\label{E:Sht-Zip}
\Sht_\mu^{\loc(2,1)}\to G\on{-Zip}^\pf_\mu.
\end{equation}
这里 $G\on{-Zip}^\pf_\mu$ 是 \cite{PWZ} 中定义的带 $G$-结构的 $F$-zip 模空间的完全化，它同胚于带 $G$-结构的 $1$-截断 Barsotti–Tate 群的模空间。

我们描述 $G=\GL_n$、$\mu=\omega_i$ 时的映射 \eqref{E:Sht-Zip}。此时 $G\on{-Zip}_{\omega_i}(R)$ 分类四元组 $(M, C^\bullet M, D_\bullet M, \phi_0,\phi_1)$，其中 $M$ 是秩 $n$ 的有限投射 $R$-模，$M=C^0M\supset C^1M\supset C^2M=0$ 是递减滤过，$D_{-1}M=0\subset D_0M\subset D_1M$ 是递增滤过，$\phi_0:\sigma^*\on{gr}_C^0M\cong \on{gr}_0^DM$ 是秩 $i$ 的有限投射 $R$-模之间的同构，$\phi_1: \sigma^*\on{gr}_C^1M\cong \on{gr}_1^DM$ 是秩 $n-i$ 的有限投射 $R$-模之间的同构。上述类型的 $F$-zip 与高度 $n$、维数 $i$ 的 $1$-截断 Barsotti–Tate 群之间的关系见 \cite[\S 9.3]{PWZ}。

注意给出 $\Sht_{\omega_i}^{\loc(2,1)}$ 的 $R$-点，就有一一对应的有限投射 $R$-模映射
$$\mE_\leftone|_{R}\xrightarrow{\beta} \mE_\rightone|_{R}\cong \sigma^*(\mE_\leftone|_{R}),$$ 其余核是秩 $i$ 的有限投射 $R$-模。
于是 \eqref{E:Sht-Zip} 把 $\Sht_{\omega_i}^{\loc(2,1)}$ 的 $R$-点送到 $G\on{-Zip}_{\omega_i}$ 的 $R$-点，后者由 $M=\mE_\rightone|_{R}$、$C^1M=\on{Im} \beta$、$D_0M=\ker(\sigma^*\beta)$ 表示。同构 $\phi_1: \sigma^*(\on{gr}_C^1M)=\sigma^*\on{Im}\beta\cong \on{gr}^D_1M=\frac{M}{\ker\sigma^*\beta}$ 是典范的。
典范同构 $\phi_0: \sigma^*(\on{gr}_C^0M)=\sigma^*\frac{M}{\on{Im}\beta}\cong \ker(\sigma^*\beta)=\on{gr}^0_DM$ 可由 \cite[(1.3.2)]{Zh1} 推出。
\end{rmk}

各种限制局部迭代 Shtuka 模空间之间存在 Hecke 对应，它们经限制态射 \eqref{trunc} 与对应 \eqref{E;Hk Shtuka} 相容。此外，这些 Hecke 对应可以复合，且复合满足结合律。我们将在命题 \ref{P:pushforward is local} 中总结这些结构，它使我们能够在下一小节定义范畴 $\on{P}^{\Hk}(\Sht^\loc_{\bar k})$。

引入下面的术语会很方便。
给定支配余权序列 $\mmu=(\mu_1,\ldots,\mu_t)$，若非负整数四元组 $(m_1,n_1,m_2,n_2)$ 满足
\begin{enumerate}
\item
$m_1-m_2=n_1-n_2$ 是 $\mu_t$-大（等价地，是 $\sigma(\mu_t)$-大），
\item
$m_2-n_1$ 是 $(\mu_1, \dots, \mu_{t-1})$-大，
\end{enumerate}
则称其为 $\mmu$-\emph{可接受}的。特别地，$m_1-n_1$ 是 $\mmu$-大。我们把 $(\infty,\infty,\infty,\infty)$ 看作对任意 $\mmu$ 都可接受。

下面的命题总结了 \cite[\S 5.3]{XZ} 的大部分内容。
为简便，我们只考虑奇点被单个余权所界的限制局部 Shtuka 的模空间。

\begin{proposition}
\label{P:pushforward is local}
\begin{enumerate}
\item 设 $\mu_1$、$\mu_2$ 是支配余权，$\nu$ 是支配余权。
设 $(m_1,n_1,m_2,n_2)$ 是 $(\mu_1+\nu, \nu)$-可接受且 $(\mu_2+\nu,\nu)$-可接受的四元组。
则存在 Hecke 对应
\begin{equation}\label{deompf1}
\Sht_{\mu_1}^{\loc(m_1,n_1)}\xleftarrow{\overleftarrow{h}_{\mu_1}^{\loc(m_1,n_1)}} \Sht_{\mu_1\mid \mu_2}^{\nu,\loc(m_1,n_1)}\xrightarrow{\overrightarrow{h}_{\mu_2}^{\loc(m_2,n_2)}} \Sht_{\mu_2}^{\loc(m_{2},n_{2})},
\end{equation}
其中 $\Sht_{\mu_1|\mu_2}^{\nu,\loc(m_1,n_1)}$ 是代数叠，且与 $(m_1,n_1)$ 的选取无关。

此外，若 $(m'_1,n'_1,m'_2,n'_2)\geq (m_1,n_1,m_2,n_2)$ 是满足同样条件的另一个四元组（允许 $(m'_1,n'_1,m'_2,n'_2)=(\infty,\infty,\infty,\infty)$），
则下图交换，且左方块是笛卡儿的。
\begin{equation*}
\label{E:cycle construction diagram full}
\xymatrix@C=17pt{
\Sht_{\mu_1}^{\loc(m'_1,n'_1)}\ar_{\res^{m'_1,n'_1}_{m_1,n_1}}[d]
&\Sht_{\mu_1|\mu_2}^{\nu,\loc(m'_1,n'_1)}\ar_{\res}[d]\ar[l]\ar[r]& \Sht_{\mu_2}^{\loc(m'_2,n'_2)}\ar^{\res^{m'_2,n'_2}_{m_2,n_2}}[d]
\\
\Sht_{\mu_1}^{\loc(m_1,n_1)}&\Sht_{\mu_1|\mu_2}^{\nu,\loc(m_1,n_1)}\ar[l]\ar[r]& \Sht_{\mu_2}^{\loc(m_2,n_2)}.
}
\end{equation*}

\item 设 $\mu_1,\mu_2,\mu_3,\nu_1,\nu_2$ 是支配余权，$(m_1,n_1,m_2,n_2,m_3,n_3)$ 是六元组，使得每个四元组 $(m_i,n_i,m_{i+1},n_{i+1})$ 都是 $(\mu_i+\nu_i,\nu_i)$-可接受且 $(\mu_{i+1}+\nu_i,\nu_i)$-可接受的。则存在自然的完全紧合态射
\begin{equation*}
\label{E:composition restricted}
\on{Comp}_{\mu_1, \mu_2, \mu_3}^{\mathrm{loc}(m_1,n_1)}: \Sht_{\mu_1|\mu_2}^{\nu_1,\loc(m_1,n_1)} \times_{\Sht_{\mu_2}^{\loc(m_2,n_2)}} \Sht_{\mu_2|\mu_3}^{\nu_2,\loc(m_2,n_2)} \longto \Sht_{\mu_1|\mu_3}^{\nu_1+\nu_2,\loc (m_1,n_1)}.
\end{equation*}

此外，若 $(m'_1,n'_1,m'_2,n'_2,m'_3,n'_3)\geq (m_1,n_1,m_2,n_2,m_3,n_3)$ 是满足同样条件的另一个六元组（允许 $(m'_1,n'_1,m'_2,n'_2,m'_3,n'_3)=(\infty,\infty,\infty,\infty,\infty,\infty)$），则下图交换。
{\footnotesize
\begin{equation*}
\label{E:projection localization truncated}
\xymatrix@C=5pt{
&&\Sht_{\mu_1|\mu_3}^{\nu_1+\nu_2,\loc(m'_1,n'_1)}
\ar[ddd]\ar[dll]\ar[dr]
\\
\Sht^{\loc(m'_1,n'_1)}_{\mu_1} \ar[d]^{\res^{m'_1,n'_1}_{m_1,n_1}} &\ar[l]\ar[d] \Sht_{\mu_1\mid\mu_2}^{\nu_1,\loc(m'_1,n'_1)} \times_{\Sht_{\mu_2}^\loc(m'_2,n'_2)} \Sht_{\mu_2\mid\mu_3}^{\nu_2,\loc(m'_2,n'_2)} \ar[l] \ar[rr] \ar[ur]_-{\mathrm{Comp}^{\loc(m'_1,n'_1)}}&& \Sht^{\loc(m_3,n_3)}_{\mu_3} \ar[d]_{\res^{m'_3,n'_3}_{m_3,n_3}}
\\
\Sht_{\mu_1}^{\loc(m_1,n_1)}
&\ar[l]\Sht_{\mu_1|\mu_2}^{\nu_1,\loc(m_1,n_1)} \times_{\Sht_{\mu_2}^{\loc(m_2,n_2)}} \Sht_{\mu_2|\mu_3}^{\nu_2,\loc(m_2,n_2)} \ar[rd]^-{\mathrm{Comp}^{\loc(m_1,n_1)}} \ar[rr]
&&\Sht_{\mu_3}^{\loc(m_3,n_3)}
\\
&& \Sht_{\mu_1|\mu_3}^{\nu_1+\nu_2, \loc(m_1,n_1)}\ar[ull]\ar[ru].
}
\end{equation*}}
此外，中间的梯形是笛卡儿的。

\item 设 $\mu_1, \mu_2, \mu_3, \mu_4, \nu_1, \nu_2, \nu_3$ 是支配余权，$(m_1, n_1, m_2,n_2, m_3,n_3, m_4,n_4)$ 是非负整数八元组，使得 $(m_i,n_i,m_{i+1},n_{i+1})$ 是 $(\mu_i+\nu_i,\nu_i)$-可接受且 $(\mu_{i+1}+\nu_i,\nu_i)$-可接受的。则
\[\on{Comp}_{\mu_1,\mu_3, \mu_4}^{\mathrm{loc}(m_1,n_1)}\circ (\on{Comp}_{\mu_1,\mu_2, \mu_3}^{\mathrm{loc}(m_1,n_1)} \times \id)=\on{Comp}_{\mu_1,\mu_2, \mu_4}^{\mathrm{loc}(m_1,n_1)}\circ (\id \times \on{Comp}_{\mu_2,\mu_3, \mu_4}^{\mathrm{loc}(m_2,n_2)}).\]
\end{enumerate}
\end{proposition}

\begin{remark}
\label{R: grp theoretical Sht}
我们给出限制局部 $G$-Shtuka 模空间的群论解释。
给定 $(m,n)$-限制局部 $G$-Shtuka，把 $\calE_\rightone|_{D_{m}}$ 与 $\calE_\leftone|_{D_{n}}$ 都平凡化，则同构 $\psi$ 定义 $K_n$ 中的元素。
注意 $\Gr_\mmu^{(n)}$ 正是分类 $\Hk^{\loc(m)}_\mmu$ 上 $\calE_\rightone|_{D_{m}}$ 与 $\calE_\leftone|_{D_{n}}$ 两者的平凡化的模空间。于是
\[
\Sht_{\mmu}^{\loc(m,n)}\cong \frac{\Gr^{(n)}_{\mmu}\times^{K_n,\sigma} K_n}{\Ad K_m},
\]
其中右端表示 $\Gr_{\mmu}^{(n)}\times K_n$ 关于 $K_m\times K_n$ 的如下作用的商：
$$(k_m,k_n)\cdot (x,k)=(k_mxk_n^{-1},\sigma(k_n)kk_m^{-1}), \quad (k_m,k_n)\in K_m\times K_n, \ (x,k)\in \Gr_\mmu^{(n)}\times K_n.$$
特别地，$\Gr_{\mu}^{(n)}$ 的 Frobenius 自同构诱导同构
\[\sigma: \Sht_\mu^{\loc(m,n)}\cong \frac{\Gr_{\mu}^{(n)}}{\Ad_\sigma K_m}.\]
（在等特征情形，这些空间有典范的去完全化，且 $\Gr_{\mu}^{(n)}$ 的 Frobenius 自同态诱导 $\Sht_\mu^{\loc(m,n)}$ 与 $\frac{\Gr_{\mu}^{(n)}}{\Ad_\sigma K_m}$ 之间的泛同胚。）
因此，后文定义 $\Sht^\loc_{\bar k}$ 上反常层范畴时，既可以用 $\Sht_{\mmu}^{\loc(m,n)}$，也可以用 $\frac{\Gr^{(n)}_{\mmu}\times^{K_n,\sigma} K_n}{\Ad K_m}$。不过，后者（从群论上看）似乎更易描述，而前者才是与 Shimura 簇（或整体 Shtuka 的模空间）相联系所需要的。对 $\Sht_{\mu_1\mid \mu_2}^{\nu,\loc(m_1,n_1)}$ 也有类似的群论描述，细节参见 \cite[注 5.3.15]{XZ}。
\end{remark}

\subsection{范畴 $\on{P}^\Hk(\Sht^\loc_{\bar k})$}
有了上一小节的准备，就可以定义范畴 $\on{D}^\Hk(\Sht^\loc_{\bar k})$；由引理 \ref{L: fully faithful}，它可以看作 $LG$ 上 $\Ad_\sigma LG$-等变层范畴（的某个子范畴）的替代物。事实上，我们将主要关注子范畴 $\on{P}^{\Hk}(\Sht^\loc_{\bar k})\subset \on{D}^\Hk(\Sht^\loc_{\bar k})$，它将成为几何 Satake 的范畴迹的接收方。

首先，局部 Shtuka 模空间（基变换到 $\bar k$）上反常层的范畴 $\mathrm{P}(\Sht_{\bar k}^\loc)$ 容易定义。
设 $\mmu$ 是支配余权序列。对非负整数对 $(m',n')$ 与 $(m,n)$，若 $m \leq m'$、$n\leq n'$，且 $m'-n'$ 与 $m-n$ 都是 $\mmu$-大（$m'\neq \infty$），则层的平移 $!$-拉回函子
$$\Res^{m',n'}_{m,n}=(\res^{m',n'}_{m,n})^![((m'-n')-(m-n))\dim G]:\on{P}(\Sht^{\loc(m,n)}_{\mu})\to \on{P}(\Sht^{\loc(m',n')}_{\mu})$$
是反常精确的。由 \eqref{true-to-trun}，有函子的典范同构
\[\Res^{m',n}_{m,n}\circ\Res^{m',n'}_{m',n}\cong \Res^{m',n'}_{m,n}.\]
若 $m\geq 1$，则函子 $\Res^{m',n}_{m,n}$ 诱导范畴的等价。另一方面，$\res^{m',n'}_{m',n}$ 是 $\Ker\big( \Aut(\calE_\rightone|_{D_{n'}}) \to \Aut(\calE_\rightone|_{D_{n}}) \big)$ 下的主丛；若 $n\geq 1$，后者又是某个幂幺（非常值）概形的完全化，因此 $\Res^{m',n'}_{m',n}$ 全忠实。

于是可以把 $\Sht^\loc$ 上反常层的范畴定义为（直和意义下的）滤过极限
\begin{equation}
\label{E:definition of P(Sloc)}
\mathrm{P}(\Sht_{\bar k}^\loc): =\varinjlim_{(\mu, m,n)} \mathrm{P}(\Sht_\mu^{\loc(m, n)}),
\end{equation}
其中极限取遍三元组 $\{(\mu, m, n) \in \xcoch(T)/W \times \ZZ^2_{\geq 0} \mid m-n \mbox{ 是 }\mu\mbox{-大}\}$（带积偏序），
连接函子由全忠实函子的如下复合给出
\[
\mathrm{P}(\Sht_{\mu}^{\loc(m, n)})
\xrightarrow{\Res^{m',n'}_{m,n}}
\mathrm{P}(\Sht_{\mu}^{\loc(m', n')})
\xrightarrow{i_{\mu,\mu',*}}
\mathrm{P}(\Sht_{\mu'}^{\loc(m' n')}),
\]
其中 $i_{\mu, \mu'}: \Sht_{\mu}^{\loc(m', n')} \to \Sht_{\mu'}^{\loc(m', n')}$ 是自然闭嵌入。
注意由紧合或光滑基变换，这些连接函子满足自然的相容性条件，因此该极限确实有意义。

注意由紧合或光滑基变换，这些连接函子满足自然的相容性条件，因此该极限确实有意义。
把上述余极限中的 $\on{P}(\Sht_{\mu}^{\loc(m, n)})$ 换成 $\on{D}(\Sht_{\mu}^{\loc(m, n)})$，也就定义了 $\on{D}(\Sht_{\bar k}^\loc)$，它带有反常 $t$-结构。

\begin{rmk}
\label{R: sheaves on Shtloc}
由 \eqref{R:pre of Sht} 与注 \ref{R: grp theoretical Sht}，可用更群论化的语言解释上述定义。即写作
\[\frac{LG}{\Ad_\sigma K}=\underrightarrow\lim_{\mu}\underleftarrow\lim_n \frac{\Gr_\mu^{(n)}}{\Ad_\sigma K},\]
从而定义
\[\on{P}(\Sht^\loc_{\bar k})=\underrightarrow\lim_{\mu,m,n}\on{P}(\frac{\Gr^{(n)}_\mu}{\Ad_\sigma K_m}).\]
\end{rmk}

对每个支配余权 $\mu$ 以及使 $m-n$ 为 $\mu$-大的数对 $(m,n)$，有自然的拉回函子
\begin{equation*}
\label{E:functor P(Hk) to P(S)fin}
\Phi^{\loc(m,n)}:=\Res^{m,n}_{m,0}: \on{P}_{K\otimes \bar k}(\Gr_\mu)\cong \on{P}(\Hk_\mu^{\loc(m)})\to \on{P}(\Sht_\mu^{\loc(m,n)}),
\end{equation*}
由 \eqref{true-to-trun} 与紧合光滑基变换，它与上述连接态射可交换。
取余极限，就得到良定义的函子
\begin{equation}
\label{E:functor P(Hk) to P(S)}
\Phi^{\loc}: \mathrm{P}_{K\otimes\bar k}(\Gr\otimes\bar k) \to \mathrm{P}(\Sht_{\bar k}^\loc).
\end{equation}

\begin{rmk}
\label{R: more object}
注意 $\mathrm{P}(\Sht_{\bar k}^\loc)$ 中有许多对象并不来自 $\on{P}_{K\otimes \bar k}(\Gr\otimes\bar k)$ 在上述函子 $\Phi^{\loc}$ 下的像。例如如例 \ref{Ex:mStau mn} 所述，$\Sht^{\loc(n,n)}_0=\bfB G(\mO/\varpi^n)$。因此有限群 $G(\mO/\varpi^n)$ 的每个表示 $\rho$ 都在 $\Sht^{\loc(n,n)}_0$ 上定义局部系统，从而定义 $\mathrm{P}(\Sht_{\bar k}^\loc)$ 中的对象。只要 $n>0$，该对象就不在 $\Phi^\loc$ 的本质像中。更一般地，对每个 $\mu$ 与 $G(\mO)$ 的每个表示 $\rho$，都有 $\mathrm{P}(\Sht_{\bar k}^\loc)$ 中的对象 $\mF_{\mu,\rho}$，来自近循环函子构造 \cite{GL}。
\end{rmk}

现在我们定义范畴 $\rmP^{\Hk}(\Sht^\loc_{\bar k})$。它与 $\rmP(\Sht_{\bar k}^\loc)$ 有相同的对象，只需按下述方式说明态射空间及其复合。
称 $\rmP(\Sht_{\bar k}^\loc)$（等价地 $\rmP^\mathrm{Hk}(\Sht^\loc_{\bar k})$）中的对象 $\mF$ 为连通的，若它来自余极限 \eqref{E:definition of P(Sloc)} 中某个 $(\mu,m, n)$ 对应的 $\rmP(\Sht_\mu^{\loc(m,n)})$。此时记 $\mF_\mu^{(m,n)}\in \rmP(\Sht_\mu^{\loc(m,n)})$ 为代表 $\mF$ 的唯一（在同构唯一意义下）对象。由于 $\rmP(\Sht_{\bar k}^\loc)$ 的每个对象都是连通对象的直和，只需定义两个连通对象 $\calF_1, \calF_2$ 之间的态射空间，再按线性性推广到一般对象。

回忆对七元组 $(\mu_1,\mu_2,\nu,m_1,n_1,m_2,n_2)$，若 $(m_1,n_1,m_2,n_2)$ 是 $(\mu_1+\nu,\nu)$-可接受且 $(\mu_2+\nu,\nu)$-可接受的，则由命题 \ref{P:pushforward is local} 有对应 \eqref{deompf1}。给定满足类似可接受条件的另一个七元组 $(\mu'_1,\mu'_2,\nu',m'_1,n'_1,m'_2,n'_2)\geq (\mu_1,\mu_2,\nu,m_1,n_1,m_2,n_2)$，有下图
\begin{equation*}
\label{E: loc Hecke corr changing level}
\begin{CD}
\Sht_{\mu_1}^{\loc(m_1,n_1)}@<\res^{m'_1,n'_1}_{m_1,n_1}<< \Sht_{\mu_1}^{\loc(m'_1,n'_1)}@>i_{\mu_1,\mu'_1}>> \Sht_{\mu'_1}^{\loc(m'_1,n'_1)}\\
@AAA@AAA@AAA\\
\Sht_{\mu_1\mid\mu_2}^{\nu,\loc(m_1,n_1)}@<\res<<\Sht_{\mu_1\mid\mu_2}^{\nu,\loc(m'_1,n'_1)}@>i>>\Sht_{\mu'_1\mid\mu'_2}^{\nu',\loc(m'_1,n'_1)}\\
@VVV@VVV@VVV\\
\Sht_{\mu_2}^{\loc(m_2,n_2)}@<\res^{m'_2,n'_2}_{m_2,n_2}<<\Sht_{\mu_2}^{\loc(m'_2,n'_2)}@>i_{\mu_2,\mu'_2}>>\Sht_{\mu'_2}^{\loc(m'_2,n'_2)}
\end{CD}
\end{equation*}
其中向左的箭头是等维的、完全光滑的，且相对维数相同；向右的箭头是闭嵌入。注意由命题 \ref{P:pushforward is local}，左上方块是笛卡儿的，而右上方块中所有箭头都完全紧合。

现在，对 $\on{P}^\Hk(\Sht^\loc_{\bar k})$ 的两个连通对象 $\mF_1,\mF_2$，把 $\mathrm{Mor}_{\rmP^\Hk(\Sht^\loc_{\bar k})} (\calF_1, \calF_2)$ 定义为
\begin{equation*}
\label{E: hom space of PcorrSht}
\varinjlim_{(\mu_1,\mu_2,\nu,m_1,n_1,m_2,n_2)} \mathrm{Corr}_{\Sht_{\mu_1|\mu_2}^{\nu, \loc(m_1,n_1)}} \big( (\Sht_{\mu_1}^{\loc(m_1,n_1)}, \calF_{1,\mu_1}^{(m_1,n_1)}),\,( \Sht_{\mu_2}^{\loc(m_2,n_2)},\calF_{2,\mu_2}^{(m_2,n_2)}) \big),
\end{equation*}
其中余极限取遍所有（带偏序的）满足上述可接受条件、且使 $\mF_{i,\mu_i}^{(m_i,n_i)}$ 代表 $\mF_i$ 的七元组，余极限的连接映射由 $i_!\circ\res^![((m'-n')-(m-n))\dim G]$ 给出，其中 $i_!$ 表示按 \eqref{ASS:pushforward correspondence} 沿 $i$ 的上同调对应前推，$\res^!$ 表示按 \eqref{ASS:smooth pullback correspondence2} 的上同调对应拉回。如 \cite[\S 5.4.1]{XZ} 所述，这些连接映射可以复合，该余极限良定义。

接着说明态射的复合。由线性性，只需对连通对象定义，并可设所给态射 $c_1: \calF_1 \to \calF_2$ 与 $c_2: \calF_2 \to \calF_3$ 都实现为对应（而非对应的线性组合）
{\footnotesize
\begin{align*}
c^{(\mu_1,\mu_2,\nu_1,m_1,n_1,m_2,n_2)}_1 \in \mathrm{Corr}_{\Sht_{\mu_1|\mu_2}^{\nu_1,\loc(m_1,n_1)}}\big( (\Sht_{\mu_1}^{\loc(m_1, n_1)}, \calF_{1,\mu_1}^{(m_1,n_1)}),\ (\Sht_{\mu_2}^{\loc(m_2, n_2)}, \calF_{2,\mu_2}^{(m_2,n_2)}) \big)&, \quad\textrm{以及}\\
c^{(\mu_2,\mu_3,\nu_2,m_2,n_2,m_3,n_3)}_2 \in \mathrm{Corr}_{\Sht_{\mu_2|\mu_3}^{\nu_2,\loc(m_2,n_2)}}\big( (\Sht_{\mu_2}^{\loc(m_2, n_2)},\ \calF_{2,\mu_2}^{(m_2,n_2)}) \,\ (\Sht_{\mu_3}^{\loc(m_3, n_3)}, \calF_{3,\mu_3}^{(m_3,n_3)})\big)&,
\end{align*}}
其中各七元组满足上述条件。
按上同调对应复合的形式体系（见定义 \ref{AD:correspondences}），复合 $c^{(\mu_2,\mu_3,\nu_2,m_2,n_2,m_3,n_3)}_2 \circ
c^{(\mu_1,\mu_2,\nu_1,m_1,n_1,m_2,n_2)}_1$ 支撑在如下对应上
\[
\Sht_{\mu_1}^{\loc(m_1,n_1)} \leftarrow \Sht_{\mu_1|\mu_2}^{\nu_1,\loc(m_1, n_1)} \times_{ \Sht_{\mu_2}^{\loc(m_2,n_2)}}
\Sht_{\mu_2|\mu_3}^{\nu_2,\loc(m_2, n_2)}  \to  \Sht_{\mu_3}^{\loc(m_3,n_3)}.
\]

然后把复合 $c_2\circ c_1$ 定义为从 $\calF_1$ 到 $\calF_3$ 的上同调对应，使 $(c_2\circ c_1)^{(\mu_1,\mu_3,\nu_1+\nu_2,m_1,n_1,m_3,n_3)}$ 由 $c^{(\mu_2,\mu_3,\nu_2,m_2,n_2,m_3,n_3)}_2 \circ
c^{(\mu_1,\mu_2,\nu_1,m_1,n_1,m_2,n_2)}_1$ 沿完全紧合态射
\[
\Sht_{\mu_1|\mu_2}^{\nu_1,\loc(m_1, n_1)} \times_{ \Sht_{\mu_2}^{\loc(m_2,n_2)}}
\Sht_{\mu_2|\mu_3}^{\nu_2,\loc(m_2, n_2)} \xrightarrow{\mathrm{Comp}^{\loc(m_1,n_1)}}\Sht_{\mu_1|\mu_3}^{\nu_1+\nu_2,\loc(m_1, n_1)}
\]
的前推（按 \eqref{ASS:pushforward correspondence} 的意义）给出。利用命题 \ref{P:pushforward is local} 与 \S \ref{Sec:cohomological correspondence} 中的若干引理，可以证明 $c_2\circ c_1$ 良定义，并满足自然的结合律 $(c_3\circ c_2) \circ c_1 \cong c_3 \circ (c_2 \circ c_1)$。至此 $\on{P}^{\Hk}(\Sht^\loc_{\bar k})$ 的构造完成。

我们描述 $\on{P}^{\on{Hk}}(\Sht^\loc_{\bar k})$ 中某些对象的自同态环。
对每个 $n>0$，记 $\on{reg}_n$ 为 $G(\mO/\varpi^n)$ 的正则表示。它在 $\on{Sht}^{\loc(n,n)}_{0}=\bfB G(\mO/\varpi^n)$ 上定义局部系统
$\delta_{\on{reg}_n}^{(n,n)}$（见注 \ref{R: more object}）。记 $\delta_{\on{reg}_n}$ 为 $\on{P}^{\Hk}(\Sht^\loc_{\bar k})$ 中相应的对象。
注意 $\delta_{\on{reg}_0}$ 含于函子 \eqref{E:functor P(Hk) to P(S)} 的像中，我们也把它记作 $\widetilde{\delta_{e}}$。
下面的结论是 \eqref{E:endo for unit} 与 \eqref{E:endo for reg} 的仿射类似物。
\begin{prop}
\label{P: endo of unit}
有典范同构
$$\End_{\on{P}^{\Hk}(\Sht^\loc_{\bar k})}(\delta_{\on{reg}_n})\cong C_c(K(n)\backslash G(F)/K(n),E)^{\on{op}},$$ 其中 $K(n)=K^{(n)}(k)$ 是第 $n$ 级主同余子群，$C_c(K(n)\backslash G(F)/K(n),E)$ 表示相应的 Hecke 代数。
特别地，$\End_{\on{P}^{\Hk}(\Sht^\loc_{\bar k})}(\widetilde{\delta_{e}})$ 同构于 $E$ 值球面 Hecke 代数 $H_G\otimes E=C_c(G(\mO)\backslash G(F)/G(\mO),E)$。
\end{prop}

\begin{rmk}
完全相同的定义给出 $\on{D}^\Hk(\Sht^{\loc}_{\bar k})$。此外，若 $\on{Iw}\subset K$ 是 Iwahori 子群，则可类似地定义带 Iwahori 层级结构的局部 Shtuka 的模空间，记作 $\Sht^{\loc}_{\on{Iw}}$。同样可以定义 Hecke 对应 $\Hk(\Sht^{\loc}_{\on{Iw}})$ 与范畴 $\on{D}^\Hk(\Sht^{\loc}_{\on{Iw},\bar k})$。如注 \ref{R:springer} 所述，若 $\cha E=0$，则它在幂等完备化之后等价于 $\on{D}^\Hk(\Sht^{\loc}_{\bar k})$。
\end{rmk}

\subsection{几何 Satake 的范畴迹}
现在可以陈述下一节算术应用的关键要素之一了。设 $\hat G$ 是 $G$ 的 Langlands 对偶群，定义在 $\Ql$ 上，即经由几何 Satake（定理 \ref{intro:geomSat} 与注 \ref{R: birth of the dual group}）构造的 Tannakian 群。它带有 $k$ 的几何 $q$-Frobenius $\sigma$ 的作用。

回忆例 \ref{E: trace of rep}：拉回函子 $\on{Rep}(\hat G)\to \Coh(\frac{\hat G}{c_\sigma\hat G})$ 把 $\Coh_{fr}^{\hat G}(\hat G\sigma)$ 与 $\on{Rep}(\hat G)$ 的 $\sigma$-扭范畴迹等同起来，其中 $\Coh_{fr}^{\hat G}(\hat G\sigma)$ 是由那些 $\widetilde{V}$（$\hat G$ 的代数表示 $V$ 的拉回）张成的 $\Coh(\frac{\hat G}{c_\sigma\hat G})$ 的全子范畴。\zy{[译者按（上下文推断，无外部依据）：原文此处写作「$G$ 的代数表示」，据上文应为「$\hat G$ 的代数表示」；译文按下标改正，并登记为勘误。]}

有了以上准备，下面的定理就变得自然了。
\begin{thm}
\label{T: trace of geom Sat}
存在函子
\[S: \Coh_{fr}^{\hat G}(\hat G\sigma)\to \on{P}^{\Hk}(\Sht^\loc_{\bar k}),\]
使下图交换
\[
\begin{CD}
\on{Rep}(\hat G)@>\cong>>\on{P}_{K\otimes \bar k}(\Gr\otimes \bar k)\\
@VVV@VVV\\
\Coh_{fr}^{\hat G}(\hat G\sigma)@>>> \on{P}^{\Hk}(\Sht^\loc_{\bar k})
\end{CD}
\]
特别地，$S(\mO)=\widetilde{\delta_e}$，其中 $\mO$ 表示 $\frac{\hat G}{c_\sigma \hat G}$ 的结构层。

此外，映射
\[S: \mathbf J=\Ql[\hat G\sigma]^{\hat G}=\Gamma([\hat G\sigma/\hat G],\mO)\to \End(S(\mO))\cong H_G\otimes \Ql\]
与 Satake 同构 \eqref{Sat isom} 一致。
\end{thm}

\begin{rmk}
函子 $S$ 应当是全忠实的。这应当可由 \cite[\S 4]{XZ} 中对某些仿射 Deligne–Lusztig 簇的不可约分支的研究推出。不过，最好能有该事实的概念性证明。
\end{rmk}

\begin{rmk}
\label{R:affine springer}
本节的大部分讨论也适用于
$X=\frac{LG}{\Ad K}$ 与 $Y=\frac{LG}{\Ad LG}$ 的情形，其中 $K\subset LG$ 以通常的共轭作用在 $LG$ 上。例如有 Hecke 对应
\[\frac{LG}{\Ad K}\xleftarrow{\overleftarrow{h}}\Hk:= X\times_YX\xrightarrow{\overrightarrow{h}}\frac{LG}{\Ad K},\]
而 $\overleftarrow{h}: \overrightarrow{h}^{-1}(\frac{\Gr_\mu^{(\infty)}}{\Ad K})\to \frac{LG}{\Ad K}$ 的纤维就是所谓的广义（或群版本）仿射 Springer 纤维（与注 \ref{R:fiber of hloc} 比较）。
若 $m-n$ 是 $\mu$-大，则有空间
$\frac{\Gr_\mu^{(n)}}{\Ad K_m}$，以及它们之间的 Hecke 叠。可以类似地定义 $\on{D}^{\Hk}(\frac{LG}{\Ad K})$。定理 \ref{T: trace of geom Sat} 在此情形有对应物，给出仿射 Springer 作用。
\end{rmk}

\section{对 Shimura 簇的应用}
\label{S:app to Sh}
我们将按照 \cite[\S 7]{XZ}，在\emph{不同} Shimura 簇的模 $p$ 纤维之间构造上同调对应。本节需要读者对 Shimura 簇模 $p$ 几何的基础理论有一定了解，必要的背景可参阅 \cite[\S 7.1, \S 7.2]{XZ}。

\subsection{Shimura 簇模 $p$ 纤维之间的上同调对应}

设 $(G,X)$ 是（弱）Shimura 数据。固定开紧子群 $K\subset G(\bA_f)$，并记
\[\Sh_K(G,X)=G(\bQ)\backslash X\times G(\bA_f)/K\]
为相应的 Shimura 簇。该 Shimura 数据定义 $G_\bC$ 的一类共轭单参数子群 $\mu$，通常称为 Shimura 余特征标。
设 $(\hat G,\hat B,\hat T)$ 是 $G$ 的 Langlands 对偶群（见注 \ref{R: birth of the dual group}）。则 $\mu$ 决定 $\hat G$ 的一个最高权为 $\mu$ 的不可约表示。
回忆 $\Sh_K(G,X)$ 有定义在反射域 $E=E(G,X)\subset \bC$ 上的典范模型。表示 $V_\mu$ 可典范地延拓为 $\hat G\rtimes \Gal(\overline\bQ/E)$ 的表示。
若 $(G,X)$ 是 Hodge 型，则可把 $\Sh_K(G,X)$ 解释为带附加结构的阿贝尔簇的模空间。

现在取 $p>2$ 为关于 $(G,X,K)$$p$ 非分歧的素数，这里的意思是 $K_p\subset G(\bQ_p)$ 是超特殊开紧子群。
此外，假设 $(G,X)$ 是 Hodge 型。在这些假设下，由 Kisin（\cite{Ki}）与 Vasiu（\cite{Vasiu}）的工作，$\Sh_K(G,X)$ 在 $\mO_{E,(v)}$ 上有典范整模型 $\mathscr S_K(G,X)$，其中 $v$ 是 $E$ 上位于 $p$ 的素位，$\FF_v$ 是它的剩余域。记 $\overline{\mathscr S}$ 为其模 $p$ 纤维，$\Sh_\mu$ 为 $\overline{\mathscr S}$ 的完全化。Kisin 的构造还给出 $\Sh_\mu$ 上奇点被 $\mu$ 所界的 $G$-Shtuka $\mE_1\dashrightarrow \mE_0\cong {^\sigma}\mE_1$。即
\[\mE_0:=\mE_{\on{cris}},\]
其中 $\mE_{\on{cris}}$ 是 $\Sh_\mu$ 上通常的晶体 $G$-主丛（\cite[推论 7.1.14]{XZ}）。
于是得到预叠的态射
\[\loc_p: \Sh_{\mu}\to \Sht_{\mu}^\loc.\]
在 Siegel 情形，$\loc_p$ 是把阿贝尔簇送到其底层的 $p$-可除群的态射的完全化。

按 Serre–Tate 理论，该态射（在完全化之前）是形式 \'etale 的。然而形式 \'etaleness 对完美概形之间的态射不是好的概念，因此我们不用它。此外，$\Sht_\mu^\loc$ 不是代数叠，直接处理它很困难。所以我们把它与从 $\Sht_{\mu}^\loc$ 到限制局部 Shtuka 模空间的态射复合。设 $(m,n)$ 是非负整数对，使 $m-n$ 是 $\mu$-大。令
\[\loc_p(m,n):\Sh_{\mu}\xrightarrow{\loc_p} \Sht_{\mu}^\loc\xrightarrow{\res_{m,n}} \Sht_{\mu}^{\loc(m,n)}.\]

下面的结果（\cite[命题 7.2.4]{XZ}）使我们能用前面定义的局部 Shtuka 研究 Shimura 簇的特殊纤维。

\begin{prop}\label{smoothness}
态射 $\loc_p(m,n)$ 完全光滑。
\end{prop}
\begin{rmk}
借助注 \ref{Ex:Sht and GZip}，该命题细化了熟知的完全光滑态射 $\Sh_\mu\to G\on{-zip}^\pf_\mu$（例如见 \cite{Zha}），后者可用来定义 Shimura 簇上的 Ekedahl–Oort 分层。
\end{rmk}

下面讨论 Shimura 簇模 $p$ 纤维之间的上同调对应。

设 $(G_1,X_1)$ 与 $(G_2,X_2)$ 是两个 Shimura 数据。
设 $\theta: G_1\otimes\bA_f\simeq G_2\otimes\bA_f$ 是同构。固定开紧子群 $K_i\subset G_i(\bA_f)$，它充分小，使得 $\theta K_1=K_2$；
并假设 $K_{1,p}$（从而 $K_{2,p}$）超特殊。
设 $\underline G_i$ 是由 $K_{i,p}$ 确定的 $G_{i,\bQ_p}$ 在 $\bZ_{p}$ 上的整模型。同构 $\theta: \underline G_1\simeq \underline G_2$ 使我们能等同它们的 Langlands 对偶群 $(\hat G,\hat B,\hat T)$，其上 Frobenius $\sigma_p$ 以自同构作用。固定同构 $\iota: \bC\simeq \overline\bQ_p$。
记 $\{\mu_i\}$ 为由 $X_i$ 确定的 $G_{i\bC}$ 的 Hodge 余特征标所在共轭类；经同构 $\iota$，它们可看作 $(\hat G,\hat B,\hat T)$ 的支配特征标。
记 $V_{\mu_1}$ 与 $V_{\mu_2}$ 为 $\hat G$ 相应的最高权表示，记 $\widetilde{V_{\mu}}$ 与 $\widetilde{V_{\mu'}}$ 为 $\frac{\hat G}{c_{\sigma_p}\hat G}$ 上相应的向量丛（如例 \ref{E: trace of rep} 所述）。

设 $E_i\subset \bC$ 是 $(G_i,X_i)$ 的反射域。设 $v$ 是 $E=E_1E_2$ 上位于 $p$ 的素位，由同构 $\iota$ 确定。
记 $\Sh_{\mu_i,K_i}$ 为 $\Sh_{K_i}(G_i,X_i)$ 的典范整模型的模 $p$ 纤维（基变换到 $k_{v}$）。
假设存在完美归纳概形 $\Sh_{\mu_1|\mu_2}$，使下图交换且两个方块都是笛卡儿的
\begin{equation}
\label{E:global to local cartesian1}
\xymatrix{
\Sh_{\mu_1,K_1} \ar_{\loc_p}[d] & \ar_{\overleftarrow{h}_{\mu_1}}[l] \Sh_{\mu_1|\mu_2} \ar^{\overrightarrow{h}_{\mu_2}}[r] \ar[d] & \Sh_{\mu_2,K_2} \ar^{\loc_p}[d] \\
\Sht_{\mu_1}^\loc & \ar_{\overleftarrow{h}_{\mu_1}^\loc}[l] \Sht_{\mu_1|\mu_2}^{\loc} \ar^{\overrightarrow{h}_{\mu_2}^\loc}[r]& \Sht_{\mu_2}^\loc.
}
\end{equation}
注意 $\Sht_{\mu_1\mid\mu_2}^\loc$ 非空等价于如下条件
\begin{equation}
\label{E: exotic p}
\Hom_{\Coh^{\hat G}_{fr}(\hat G\sigma_p)}(\widetilde{V_{\mu_1}},\widetilde{V_{\mu_2}})\neq 0.
\end{equation}
但即便 $\Sht_{\mu_1\mid\mu_2}^\loc$ 非空，要保证 \eqref{E:global to local cartesian1} 存在，对 $(G_i,X_i)$ 仍有限制。具体地说，可以定义 $\Sh_{\mu_1\mid\mu_2}$ 使 \eqref{E:global to local cartesian1} 中左方块成为笛卡儿的。由于 $\overleftarrow{h}^\loc_{\mu_1}$ 归纳紧合（见注 \ref{R:fiber of hloc}），$\Sh_{\mu_1\mid\mu_2}$ 由归纳完美概形表示。但一般来说，并不存在从 $\Sh_{\mu_1\mid\mu_2}$ 到 $\Sh_{\mu_2}$ 的态射使右方块交换（且笛卡儿）。这里不讨论 \eqref{E:global to local cartesian1} 存在性所需的一般条件，只提及下面两个极端情形，更多讨论参见 \cite{XZ}。

\begin{thm}
\label{T: corr. Shimura}
图 \eqref{E:global to local cartesian1} 在如下情形存在。假定 $(G_i,X_i)$ 都是 Hodge 型且 $p>2$。
\begin{enumerate}
\item
当 $(G_1,X_1)=(G_2,X_2)$ 且 $\theta$ 是恒等映射时，$\Sh_{\mu_1\mid \mu_2}$ 存在，它是 $\Sh_K(G_1,X_1)$ 的 $p$ 幂 Hecke 对应的某个自然整模型的模 $p$ 纤维的完全化。
\item
当存在内扭曲\zy{[译者按（上下文推断，无外部依据）：uniformization 取「均化」，指用群论数据与 Rapoport–Zink 空间把 Shimura 簇的模 $p$ 纤维写成商。]} $\Psi: G_1\to G_2$，它在所有有限素位处等于 $\Ad_h \circ \theta$（对某个 $h\in G_2(\overline{\bA}_f)$），并把 $G_{2\bR}$ 实现为 $G_{1\bR}$ 的模中心紧的内形式（特别地，$\Sh_{\mu_2}$ 是有限集，通常称为 Shimura 集）时，$\Sh_{\mu_1\mid \mu_2}$ 存在，并可看作 $\Sh_{\mu_1}$ 的基本链的 Rapoport–Zink 均化。
\end{enumerate}
\end{thm}

\begin{rmk}
如前所述，若 $(G,X)=(G_1,X_1)=(G_2,X_2)$，则 $\Sh_{\mu\mid\mu}$ 是 Hecke 对应
\begin{equation*}
\label{E: pHecke corr}
\xymatrix@C=-20pt{
&G(\bQ)\backslash X\times G(\bA_f^p)/K^p\times G(\bQ_p)\times^{K_p}G(\bQ_p)/K_p\ar^{\overleftarrow{h}}[ld]\ar_{\overrightarrow{h}}[rd]&\\
G(\bQ)\backslash X\times G(\bA_f)/K && G(\bQ)\backslash X\times G(\bA_f)/K.}
\end{equation*}
的某个自然整模型的模 $p$ 纤维。
若 $(G_1,X_1)\neq (G_2,X_2)$ 但 \eqref{E: exotic p} 成立，则 $\Sh_{\mu_1\mid \mu_2}$ 可看作\emph{不同} Shimura 簇的模 $p$ 纤维之间的「异型 Hecke 对应」。这些对应不能提升到特征零，它们给出 Shimura 簇上一大类特征 $p$ 的轮（见下一小节）。
\end{rmk}

把 \eqref{E:global to local cartesian1} 与命题 \ref{P:pushforward is local} 结合，得到下面的交换图，其左方块是笛卡儿的
\[
\xymatrix{
\Sh_{\mu_1,K_1} \ar_{\loc_p(m_1,n_1)}[d] && \ar_{\overleftarrow{h}_{\mu_1}}[ll] \Sh_{\mu_1|\mu_2} \ar^{\overrightarrow{h}_{\mu_2}}[rr] \ar[d] && \Sh_{\mu_2,K_2} \ar^{\loc_p}[d] \\
\Sht_{\mu_1}^{\loc(m_1,n_1)} && \ar_{\overleftarrow{h}_{\mu_1}^{\loc(m_1,n_1)}}[ll] \Sht_{\mu_1|\mu_2}^{\loc(m_1.n_1)} \ar[rr]&& \Sht_{\mu_2}^{\loc(m_2,n_2)}.
}
\]

现在，把 $\on{P}^{\Hk}(\Sht_{\bar k}^\loc)$ 中 $(\Sht_{\mu_1}^{\loc(m_1,n_1)},\mF_1)$ 与 $(\Sht_{\mu_1}^{\loc(m_2,n_2)},\mF_2)$ 之间的态射，沿上图中的竖直映射按 \eqref{ASS:smooth pullback correspondence2} 的意义拉回，就得到 $(\Sh_{\mu_1}, \loc_p(m_1,n_1)^!\mF_1)$ 与 $(\Sh_{\mu_2},\loc_p(m_2,n_2)^!\mF_2)$ 之间的上同调对应。取紧支上同调并应用引理 \ref{AL:pushforward compatible with composition}，即得如下定理。
\begin{thm}
\label{introT: spectral action}
\begin{enumerate}
\item
设 $(G_i,X_i), i=1,2,3$ 是一族 Shimura 数据，连同满足自然上闭链条件的同构 $\theta_{ij}:G_i(\bA_f)\cong G_{j}(\bA_f)$。固定共同的层级结构 $K$。取 $p>2$ 为非分歧素数，并固定同构 $\iota: \bC\simeq \overline\bQ_p$。
此外，假设对每对 $(G_i,X_i),(G_j,X_j)$，笛卡儿图 \eqref{E:global to local cartesian1} 存在。记 $\{\mu_i\}$ 为 $(G_i,X_i)$ 的 Shimura 余特征标所在共轭类，$V_i=V_{\mu_i}$ 为 $\hat G$ 相应的最高权表示，$\widetilde{V_i}$ 为 $\frac{\hat G}{c_{\sigma_p}\hat G}$ 上相应的向量丛。记 $d_i=\dim \Sh_K(G_i,X_i)$，$\mH^p$ 为素-to-$p$ Hecke 代数。
设 $v\mid p$ 是由 $\iota$ 确定的、所有反射域之复合域上的素位。
则 $\mathbf J$ 自然作用在 $\on{H}^*_c(\Sh_{\mu_i,\overline\bF_v},\Ql(d_i/2))$ 上，并有
典范 $\mathbf J$-等变映射
\[S_{\theta_{ij}}:\Hom(\widetilde{V_i},\widetilde{V_j})\to \Hom_{\mathbf J\otimes\mH^p}\big(\on{H}_c^{*+d_i}(\Sh_{\mu_i,\overline\bF_v},\Ql(d_i/2)),\on{H}_c^{*+d_j}(\Sh_{\mu_j,\overline\bF_v},\Ql(d_j/2))\big),\]
它与两侧的自然复合相容。特别地，代数 $\End(\widetilde{V_{\mu}})$ 自然作用在 $\on{H}_c^{*}(\Sh_{\mu, \overline \FF_v},\Ql)$ 上。

\item 当 $\Sh_{\mu,\overline \FF_v}$ 是 Shimura 集时，$\mathbf J=\End(\widetilde{V_{\mu}})$ 在
$$\on{H}_c^{*}(\Sh_{\mu, \overline \FF_v},\Ql)\cong C_c(G(\bQ)\backslash G(\bA_f)/K, \Ql)$$
上的作用，在 Satake 同构 \eqref{Sat isom} 下与通常的 Hecke 代数作用一致。
\end{enumerate}
\end{thm}

\begin{rmk}
(1) 上述定理特别以几何方式给出了一些 Jacquet–Langlands 转移，这正是 David Helm（\cite{He}）最早在一些特殊情形下研究的内容。

(2) 定理中 $\mathbf J\subset \End(\widetilde{V_{\mu}})$ 在 $\on{H}_c^{*}(\Sh_{\mu},\Ql)$ 上的作用，是 V. Lafforgue 的 $S$-算子在 Shimura 簇情形的类比。然而，更大的代数 $\End(\widetilde{V_{\mu}})$ 在 $\on{H}_c^{*}(\Sh_{\mu},\Ql)$ 上的作用是新的，即便在函数域情形（此时用 Shtuka 的模空间代替 Shimura 簇）也是如此。
下面的猜想是 V. Lafforgue 的 $S=T$ 定理的类比。
\begin{conjecture}
\label{introconj: S=T}
在 Satake 同构 \eqref{Sat isom} 下，上述定理中 $\mathbf J$ 的作用与通常的 Hecke 代数在 $\on{H}_c^{*}(\Sh_{\mu, \overline \FF_v},\Ql)$ 上的作用一致。
\end{conjecture}

(3) 由 \cite[\S 6.3]{XZ2}，上述 $S=T$ 猜想蕴含 Hodge 型 Shimura 簇的同余关系猜想（即所谓的 Blasius–Rogawski 猜想）。
\end{rmk}

\begin{rmk}
\label{R: coh of Shimura}
下面是由 Drinfeld 对 V. Lafforgue 工作的解释所启发的、关于 $\overline{\scrS}$ 的上同调的（启发式）猜想性描述：
存在分次 $\Ql$-向量空间 $A^*$（只依赖于 $G\otimes\bA_f$），带有 $\mH_K$-作用以及与它交换的 $\Ql[\hat{G}\sigma_p]$-作用，使得
$$\on{H}_c^{*+d}(\overline{\scrS}_{\overline\bF_v}, \Ql(d/2)) = (A^*\otimes V_{\mu})^{\hat G}.$$
显然定理 \ref{introT: spectral action} 由该猜想推出。事实上，对 $\Sh_K(G,X)$ 的上同调也有类似的猜想性描述，我们打算另找机会详细讨论。
\end{rmk}

\subsection{某些 Shimura 簇模 $p$ 纤维上的 Tate 轮}
我们运用上述机制，对许多 Shimura 簇的模 $p$ 纤维验证 Tate 猜想的「一般」情形。

设 $(G,X,K)$ 如上一小节，$p>2$ 为非分歧素数。
设 $(\hat G, \hat B,\hat T)$ 是 $G_{\bQ_p}$ 的 Langlands 对偶群，配有注 \ref{R: birth of the dual group} 中的 Borel 子群与极大环面。如前固定同构 $\iota: \overline\bQ_p\simeq \bC$，并把 Shimura 余特征标 $\mu$ 看作 $(\hat G,\hat B,\hat T)$ 的支配权。设 $v\mid p$ 是由 $\iota$ 确定的 $E$ 的素位，$\overline{\scrS}$ 如前为 $\Sh_K(G,X)$ 在 $\mO_{E,(v)}$ 上的典范整模型的模 $p$ 纤维。

对 $\la\in\xch(\hat T)$ 与 $\hat G$ 的表示 $V$，记 $V(\la)$ 为 $V$ 的 $\la$-权子空间（关于 $\hat T$）。
定义如下格
\[\Lambda^\Tate_p=\Big\{\la\in \xch(\hat T)\; \Big|\; \sum_{i=0}^{m-1} \sigma_p^i(\la)\in \xcoch(Z_G)\Big\}\subset \xch(\hat T).\]
这里 $Z_G$ 表示 $G$ 的中心，$\xcoch(Z_G)$ 表示 $Z_G$ 的余权格，它是 $\xch(\hat T)$ 的自然子群。对 $\hat G$ 的表示 $V$，定义空间
\[V^\Tate=\bigoplus_{\la\in\Lambda^\Tate_p} V(\la).\]

我们特别关心条件 $V^\Tate_\mu\neq 0$，其中 $\mu$ 是 $G$ 的 Shimura 余特征标。
如 \cite{XZ} 的引言所述，\zy{[来源：Tate cycle 译作「Tate 轮」，沿用《仿射 Grassmann 簇与几何 Satake 等价》（SY02-17）的术语表；轮类映射即 cycle class map。]}在注 \ref{R: coh of Shimura} 关于 Shimura 簇上同调的猜想性描述之下，若 Hecke 模 $\pi_f$ 在 $p$ 处的 Satake 参数足够一般，则该向量空间维数的某个倍数 $a(\pi_f)$ 应当等于 $\overline{\scrS}$ 中间维数紧支上同调的 $\pi_f$-分支中 Tate 轮空间维数。
此外，该空间通常很大。例如当 $G$ 是二次数域上的奇酉（similitude）群、签名 $(i,n-i)$ 时，该空间在惰性素数处的维数为 $\begin{pmatrix}\frac{n+1}{2}\\ i\end{pmatrix}$。

我们还需要一个记号。对代数闭域上维数为 $d$ 的（不必不可约的）代数簇 $Z$，记 ${\rm H}^{\rm BM}_{2d}(Z)(-d)$ 为最高次 Borel–Moore 同调的 $(-d)$-Tate 扭转，它是由 $Z$ 的不可约分支张成的向量空间。现设 $X$ 是有限域 $\FF_q$ 上维数为 $d+r$ 的光滑簇，$Z \subseteq X_{\overline \FF_q}$ 是维数为 $d$ 的（不必不可约的）射影子簇。则有轮类映射
\[\on{cl}: {\rm H}^{\rm BM}_{2d}(Z)(-d)\to \bigcup_{j\geq 1}{\rm H}_{\et,c}^{2r}(X_{\overline \FF_q}, \bQ_\ell(r))^{\sigma_q^j}.\]

\cite{XZ} 的主要定理如下。$\overline{\mathscr S}$ 上有自然的 Newton 分层。
记 $\overline{\mathscr S}_b$ 为基本 Newton 层，$\mH_K$ 通过（上同调）对应作用在其上（这些对应由上一小节构造的 $S$-算子给出）。
\begin{thm}
\label{T:main theorem}
假设 $(G,X)$ 是 Hodge 型，且 $G$ 的中心 $Z_G$ 连通。设 $K\subset G(\bA_f)$ 是开紧子群，记 $2d=\dim\Sh_K(G,X)$。
设 $p>2$ 是关于 $(G,X,K)$ 的非分歧素数，且 $r:=V_{\mu^*}^\Tate\neq 0$。
则：
\begin{enumerate}
\item $\overline{\scrS}$ 的基本 Newton 层 $\overline{\mathscr S}_b$ 是纯 $d$ 维的。特别地，$d$ 总是整数。此外，存在 $\mH_K$-等变同构
$${\rm H}^{\rm BM}_{2d}(\overline{\mathscr S}_{b, \overline \FF_v})(-d)\cong C(G'(\bQ)\backslash G'(\bA_f)/K)^{\oplus r}.$$
这里 $G'$ 是 $G$ 的内形式，它在 $\bA_f$ 上与 $G$ 同构，且在无穷处模中心紧；$C(G'(\bQ)\backslash G'(\bA_f)/K)$ 表示有限集 $G'(\bQ)\backslash G'(\bA_f)/K$ 上的函数空间。

\item
设 $\pi_f$ 是 $\calH_K$ 的不可约模，记
$${\rm H}^{\rm BM}_{2d}(\overline{\mathscr S}_{b, \overline \FF_v})[\pi_f]=\Hom_{\mH_K}(\pi_f,{\rm H}^{\rm BM}_{2d}(\overline{\mathscr S}_{b, \overline \FF_v})(-d))\otimes \pi_f$$
为 $\pi_f$-等型分支。则轮类映射
\[\on{cl}:{\rm H}^{\rm BM}_{2d}(\overline{\mathscr S}_{b, \overline \FF_v})(-d)\to {\rm H}^{2d}_{\et,c}\big(\overline{\mathscr S}_{\overline \FF_v},\overline\bQ_\ell(d)\big)\]
限制在 ${\rm H}^{\rm BM}_{2d}(\overline{\mathscr S}_{b, \overline \FF_v})[\pi_f]$ 上时，若 $\pi_{f,p}$（即 $\pi_f$ 在 $p$ 处的分量）的 Satake 参数是 $V_\mu$-一般的，则它是单射。

\item 假设 $\Sh_K(G,X)$ 是四元数 Shimura 簇或 Kottwitz 算术簇。则当 $\pi_{f,p}$ 的 Satake 参数强 $V_\mu$-一般时，$\on{cl}$ 的 $\pi_f$-等型分支满射到 $T^d(\pi,\Ql)\otimes\pi_f$。特别地，对这些 $\pi_f$，Tate 猜想成立。
\end{enumerate}
\end{thm}
\begin{rmk}
(1) 中心 $Z_G$ 连通这一假设并非本质，只是使表述更简单。

(2) 对 $\hat G$ 的表示 $V$，「$V$-一般」与「强 $V$-一般」Satake 参数的定义见 \cite[定义 1.4.2]{XZ}。$\hat G\sigma_p$ 中的正则半单元素总是 $V$-一般的，反之不然。见 \cite[注 1.4.3]{XZ}。

(3) 该定理的一些特殊情形最初在 \cite{HTX,TX} 中得到证明。
\end{rmk}

该定理的证明依赖若干不同要素。借助 Shimura 簇基本链的 Rapoport–Zink 均化，第 (1) 部分可归约为某些仿射 Deligne–Lusztig 簇的不可约分支问题，后者已在 \cite[\S 4]{XZ} 中研究。
最困难的部分是 (2)，我们通过计算 $\overline\scrS_b$ 中所有 $d$ 维轮之间的相交数来证明它。由第 (1) 部分，$\overline{\scrS}_b$ 中 $d$ 维轮的相交数可编码为一个 $r\times r$-矩阵，其元素在 $p$ 处的球面 Hecke 代数中。一般来说，直接显式计算这个矩阵似乎没有希望。
不过，与例 \ref{intersection by geom Sat} 中的玩具模型类似，该矩阵可以理解为 $\Coh^{\hat G}_{fr}(\hat G\sigma_p)$ 中某些态射的复合。具体地说，先把 $G'(\bQ)\backslash G'(\bA)/K$ 实现为 Shimura 集，其 Shimura 余特征标记为 $\tau$\footnote{$\tau$ 的选取有一点微妙之处，此处从略，见 \cite[注 7.4.4]{XZ}。}。
然后由定理 \ref{introT: spectral action}，该矩阵可计算为
\[ \Hom_{\Coh^{\hat G}_{fr}(\hat G\sigma_p)}(\widetilde{V_{\tau}},\widetilde{V_{\mu}})\otimes  \Hom_{\Coh^{\hat G}_{fr}(\hat G\sigma_p)}(\widetilde{V_{\mu}},\widetilde{V_{\tau}})\to  \Hom_{\Coh^{\hat G}_{fr}(\hat G\sigma_p)}(\widetilde{V_{\tau}},\widetilde{V_{\tau}})=\mathbf{J}.\]
为证明第 (2) 部分，需要判定该配对何时非退化。这本身是表示论中有趣的问题，其解决依赖于对向量值函数的 Chevalley 限制映射的研究。该矩阵的行列式已在 \cite{XZ2} 中算出，由此得到 (2)。第 (3) 部分通过比较两个迹公式证明：$G$ 的 Lefschetz 迹公式与 $G'$ 的 Arthur–Selberg 迹公式。细节参见 \cite[\S 2]{XZ}。

\begin{ex}
设 $G=\on{GU}(1,2r)$ 是虚二次扩张 $E/\bQ$ 对应的 $(2r+1)$ 个变量的酉相似群，其无穷处签名为 $(1,2r)$。它带有标准的 Shimura 数据。固定层级 $K\subset G(\bA_f)$，记 $\Sh_K(G,X)$ 为相应的 Shimura 簇。特别地当 $r=1$ 时，它就是 Picard 模曲面。取 $p$ 为非分歧惰性素数。此时 $\overline{\scrS}_b$ 是一些 Deligne–Lusztig 簇之并，由 $G'(\bQ)\backslash G'(\bA_f)/K$ 参数化，其中 $G'=\on{GU}(0,2r+1)$ 在所有有限素位与 $G$ 同构。这些轮在 $\overline{\scrS}_b$ 中的相交模式已在 \cite{VW} 中给出。然而这些轮之间的相交数要难算得多。事实上，除把定理 \ref{introT: spectral action} 用在此情形之外，我们并不知道如何对一般的 $r$ 直接计算它们（$r=1$ 的情形可以直接处理）。

此时 $\hat G=\GL_{2r+1}\times \bG_m$，$\sigma_p$ 在其上作用为 $(A,c)\mapsto (J(A^T)^{-1}J,c)$，其中 $J$ 是反对角元素全为 $1$ 的反对角矩阵。表示 $V_\mu$ 是 $\GL_{2r+1}$ 的标准表示，$\bG_m$ 通过数乘作用。可以验证 $\dim V_\mu^{\Tate}=1$（这与上面提到的 $\overline{\scrS}_b$ 的不可约分支由 $G'(\bQ)\backslash G'(\bA_f)/K$ 参数化相符）。
可取 $\tau=((1,\ldots,1),1)$，并把 $G'(\bQ)\backslash G'(\bA_f)/K$ 看作 Shimura 余特征标为 $\tau$ 的 Shimura 集。
这里按惯例把 $\hat G$ 的权格等同于 $\bZ^{2r+1}\times\bZ$。则 $\Hom_{\Coh^{\hat G}_{fr}(\hat G\sigma_p)}(\widetilde{V_{\tau}},\widetilde{V_{\mu}})$ 是 $\mathbf J$ 上的自由秩一 module。取定生成元 $\mathbf a_{\on{in}}$，它诱导 $\mathbf J\times\mH^p$-等变同态
$$S(\mathbf a_{\on{in}}): C(G'(\bQ)\backslash G'(\bA_f)/K)\to \on{H}_c^{2r}(\overline{\scrS}_{\overline{\bF}_p},\Ql(r)),$$
实现 $\overline\scrS_b$ 的轮类映射（相差一个倍数）。模 $\Hom_{\Coh^{\hat G}_{fr}(\hat G\sigma_p)}(\widetilde{V_{\mu}},\widetilde{V_{\tau}})$ 也是 $\mathbf J$ 上的自由秩一模。取定生成元 $\mathbf a_\out$，则复合
\[S(\bba_\out)\circ S(\bba_{\on{in}})=S(\bba_\out\circ\bba_{\on{in}})\]
计算这些轮（由 $\overline\scrS_b$ 的不可约分支给出）的相交矩阵。

元素 $h:=\bba_\out\circ\bba_{\on{in}} \in \mathbf J$ 已在 \cite[例 6.4.2]{XZ2} 中显式计算（差一个显然的修正）。在 Satake 同构 $\mathbf J\cong H_{G_{\bQ_p}}\otimes \Ql$ 下，它也可以显式写成 Hecke 算子的组合（\cite{XZ3}）
\begin{equation}
\label{E:Hecke operator in Tpj}
h=p^{r(r+1)}\sum_{i=0}^{r} (-1)^i(2i+1)p^{(i-r)(r+i+1)}\sum_{j=0}^{r-i} \begin{bmatrix} 2r+1-2j \\ r-i- j\end{bmatrix}_{v=-p}T_{p,j}.
\end{equation}
这里 $T_{p,j}=1_{K_p\la_j(p)K_p}$，其中 $\la_i=(1^i,0^{2r-2i+1},(-1)^i,0)$；$\begin{bmatrix} n \\ m\end{bmatrix}_v$ 是二项式系数的 $v$-类比
\[
[0]_v=1,\quad [n]_v=\frac{v^n-1}{v-1},\quad [n]_v!=[n]_v[n-1]_v\cdots[1]_v,\quad \begin{bmatrix} n \\ m\end{bmatrix}_v=\frac{[n]_v!}{[n-m]_v![m]_v!}.
\]
换言之，此情形下 $\overline\scrS_b$ 中轮的相交矩阵由 Hecke 算子 \eqref{E:Hecke operator in Tpj} 计算。细节与更多讨论参见 \cite{XZ3}。
\end{ex}

\section{附录：由对应构造的一个范畴}
设 $C\rightrightarrows X$ 是空间的群胚。
本附录（\S \ref{SS: cat of corr}）的主要目标是引入范畴 $\on{D}^C(X)$：当 $C=X\times_YX$ 且 $X\to Y$ 紧合时，它抓住了 $\on{D}(Y)$ 的大部分信息。

固定完全域 $k$，记 $\mathrm{pt}$ 为 $\Spec k$。
所谓叠（相应地，完全叠）$\mX$，指 $\Aff_k$（相应地，$\Aff_k^\pf$）上带 fpqc 拓扑的叠，其对角线可由代数空间（相应地，完全代数空间）表示，且存在从概形（相应地，完全概形）$X\to \mX$ 的平滑（相应地，完全平滑）满射。
后文中，$k$ 上（完全）有限表现的（完全）叠也称空间。关于完全代数几何基础的一些讨论，参看 \cite[附录 A]{Zh1}、\cite{BS} 与 \cite[\S A.1]{XZ}。

\subsection{上同调对应的回顾}
\label{Sec:cohomological correspondence}
我们大量用到上同调对应的形式化，故按 \cite[\S A.2]{XZ} 简要回顾若干事实。更详尽的处理见文献同前处。

设 $E$ 是系数环，$\cha k$ 在 $E$ 中可逆。
若 $X$ 是 $k$ 上（完全）有限表现概形且 $E$ 有限，记 $\on{D}(X)=\on{D}(X,E)$ 为 $E$-系数 étale 层（无界）导出范畴。对一般空间 $X$ 与可能更一般的系数环 $E$（例如 $E=\Ql$），可通过某种极限定义 $\on{D}(X)=\on{D}(X,E)$（精确表述见 \cite[\S A.1.15]{XZ}）。
记 $\on{D}^*_c(X)\subset \on{D}(X,E)$（$*=+,-,b,\emptyset$）为其可构造导出范畴，$\on{P}(X)\subset\on{D}_c^b(X)$ 为反常层范畴。记 $\omega_X\in \on{D}_c^b(X)$ 为 $X$ 上的对偶层，$\bD\mF=R\underline{\Hom}(\mF,\omega_X)$ 为 Verdier 对偶。按定义有典范赋值映射
\begin{equation}
\label{E: VD counit}
 \ev_\mF: \Delta^*(\bD\mF\boxtimes\mF)=\bD\mF\otimes \mF\to \omega_X,
\end{equation}
且若 $\mF\in\on{D}_c^b(X)$，取 Verdier 对偶即得典范余赋值映射
\begin{equation}
\label{E: VD unit}
\coev_\mF: E\to \Delta^!(\mF\boxtimes\bD\mF),\quad
\end{equation}

下面引入的概念很有用。
\begin{dfn}
\label{D: basechangeable}
空间交换方块
\begin{equation}
\label{E: basechangeable}
\begin{CD}
D@>p>>C\\
@VbVV@VVaV\\
Y@>f>>X
\end{CD}
\end{equation}
称为\emph{可基变换}\zy{[译者按（上下文推断，无外部依据）：base changeable 是原作者自定的术语，取「可基变换」，指该方块在任意基变换下仍给出同构。]}，若诱导态射 $h:D\to C\times_XY$ 可由（完全）紧合代数空间表示。
\end{dfn}

下面的引理是直接的。
\begin{lem}
\label{L: criterion basechangeable}
\begin{enumerate}
\item 如 \eqref{E: basechangeable} 的交换图可基变换，若
\begin{itemize}
\item $p$ 可由（完全）紧合代数空间表示且 $f$ 分离；
\item 或 $b$ 可由（完全）紧合代数空间表示且 $a$ 分离。
\end{itemize}
\item 在交换图
\begin{equation}
\label{E: two base changeable}
\begin{CD}
E@>q>>D@>p>>C\\
@VcVV@VbVV@VVaV\\
Z@>g>>Y@>f>>X
\end{CD}
\end{equation}
中，若两个内方块均可基变换，则外矩形也可基变换。
\end{enumerate}
\end{lem}

这一命名由下述事实支持。

\begin{lem}
\label{L:proper base change}
假设有如上 \eqref{E: basechangeable} 的可基变换空间交换图。
则存在自然的基变换同态
\begin{equation}
\label{E: base change isom}
\on{BC}^*_!: a^*f_!\to p_!b^*,
\end{equation}
若 \eqref{E: basechangeable} 是笛卡儿方块，则它是同构。
若 $p$ 与 $f$ 分离，它进入交换图
\[\begin{CD}
a^*f_!@>>> p_!b^*\\
@VVV@VVV\\
a^*f_*@>>>p_*b^*,
\end{CD}\]
其中下方箭头是自然伴随。
此外，给定 \eqref{E: two base changeable} 且两方块均可基变换，基变换同态 $a^*(fg)_!\to (pq)_!c^*$ 等于复合
$$a^*(fg)_!=a^*f_!g_!\xrightarrow{\on{BC}_!^*} f_!b^*g_!\xrightarrow{\on{BC}_!^*} f_!g_!c^*=(fg)_!c^*.$$
同样，$(fg)^*a_!\to c_!(pq)^*$ 等于复合
$$(fg)^*a_!=g^*f^*a_!\xrightarrow{\on{BC}_!^*} g^*b_!f^*\xrightarrow{\on{BC}_!^*} c_!g^*f^*=c_!(fg)^*.$$
\end{lem}

\begin{definition}
\label{AD:correspondences}
设 $(X_i,\mF_i)$（$i=1,2$）是两对数据，其中 $X_i$ 是空间，$\mF_i\in \on{D}(X_i,E)$。\emph{上同调对应} $(C,u):(X_1,\mF_1)\to (X_2,\mF_2)$ 指空间 $C\xrightarrow{c_1\times c_2} X_1\times  X_2$ 连同一个态射 $u:c_1^*\mF_1\to c_2^!\mF_2$。称 $C$ 为 $(C,u)$ 的\emph{支撑}。为简单起见，$(C,u)$ 有时也记作 $C$ 或 $u$。两个上同调对应 $(C,u)$ 与 $(C',u')$（从 $(X_1,\mF_1)$ 到 $(X_2,\mF_2)$）之间的（同）构有显然的概念。从 $(X_1,\mF_1)$ 到 $(X_2,\mF_2)$、支撑在 $C$ 上的上同调对应的同构类之集记作 $\on{Corr}_C((X_1,\mF_1),(X_2,\mF_2))$，按定义
\[\on{Corr}_C((X_1,\mF_1),(X_2,\mF_2))\cong \Hom_C(c_1^*\mF_1,c_2^!\mF_2).\]

若 $(C,u):(X_1,\mF_1)\to (X_2,\mF_2)$ 与 $(D,v):(X_2,\mF_2)\to (X_3,\mF_3)$ 是两个上同调对应，定义其复合为 $(C\times_{X_2}D,v\circ u):(X_1,\mF_1)\to (X_3,\mF_3)$，其中 $v\circ u$ 是如下复合
\[p^*c_1^*\mF_1\xrightarrow{p^*u} p^*c_2^!\mF_2\xrightarrow{\on{BC}^{*!}} q^!d_1^*\mF_2\xrightarrow{q^!v} q^!d_2^!\mF_3,\]
这里 $\on{BC}^{*!}$ 是 \eqref{E: base change isom} 中的基变换同态，$p,q$ 分别是 $C\times_{X_2}D$ 到 $C$ 与 $D$ 的投影。
\end{definition}

下面是一些例子。

\begin{ex}
\label{Ex:examples of correspondences}
\begin{enumerate}
\item 对态射 $f: X\to Y$ 与 $\mF \in \on{D}(X)$，有上同调对应 $(X\xleftarrow{\id}X\xrightarrow{f} Y, u:\mF\to f^!f_!\mF)$，从 $(X,\mF)$ 到 $(Y,f_!\mF)$，称为\emph{前推对应}，简记作 $(\Ga_f)_!$。

\item 对态射 $f: X\to Y$ 与 $\mF \in \on{D}(Y)$，有自然的上同调对应 $(Y\xleftarrow{f} X\xrightarrow{\id} X, u:= \mathrm{id}:f^*\mF\to f^*\mF)$，从 $(Y,\mF)$ 到 $(X,f^*\mF)$，称为\emph{拉回对应}，简记作 $\Ga_f^*$。

\item 若 $(C,u): (X_1,\mF_1)\to (X_2,\mF_2)$ 是上同调对应，且 $\mF_i\in \on{D}_c^b(X_i)$，则 $(C,\bD u): (X_2, \bD \mF_2)\to (X_1,\bD \mF_1)$ 也是上同调对应，称为\emph{对偶对应}。

\item 设 $(X,\mF)$ 是一对数据，$\mF\in \on{D}_c^b(X)$。则 \eqref{E: VD counit} 与 \eqref{E: VD unit} 给出上同调对应及其对偶
\[\begin{aligned}
&\ev_\mF\in \on{Corr}_{X}((X\times X, \bD\mF\boxtimes \mF),(\on{pt}, E)),\\
&\coev_\mF\in \on{Corr}_{X}((\on{pt}, E), (X\times X, \mF\boxtimes\bD\mF)).
\end{aligned}\]

\item 假设 $k=\bar k$。设 $X_1\leftarrow C\to X_2$ 是对应。则
\begin{equation*}
\label{E:corr and cycle}
\begin{split}
\on{Corr}_C((X_1,E[d_1]), (X_2, \omega_{X_2}[d_2]))&=\Hom_{\on{D}^b_c(C)}(E[d_1],\omega_C[d_2])\\
&=\on{H}^{\on{BM}}_{d_1-d_2}(C).
\end{split}
\end{equation*}
故若 $2\dim C=d_1-d_2$，则 $\on{Corr}_C((X_1,E[d_1]), (X_2, \omega_{X_2}[d_2]))$ 可等同于 $C$ 的最高次 Borel–Moore 同调。特别地，若 $X_1,X_2,C$ 是有限集，看作 $k$ 上的零维空间，则 $\on{Corr}_C((X_1,E),(X_2,E))$ 可等同于 $C$ 上 $E$ 值函数之空间。
\end{enumerate}
\end{ex}

现假设有交换图
\[\begin{CD}
X_1@<c_1<< C@>c_2>> X_2\\
@Vf_1VV@VVfV@VVf_2V\\
Y_1@<d_1<<D@>d_2>>Y_2
\end{CD}\]
且设 $(C, u:c_1^*\mF_1\to c_2^!\mF_2)$ 是从 $(X_1,\mF_1)$ 到 $(X_2,\mF_2)$ 的上同调对应。
设左交换方块可基变换（定义 \ref{D: basechangeable}），则有
\begin{equation}
\label{ASS:pushforward correspondence}
d_1^*(f_1)_!\mF_1\xrightarrow{\on{BC}^*_!} f_!c_1^*\mF_1\xrightarrow{f_!u} f_!c_2^!\mF_2\to d_2^!(f_2)_!\mF_2,
\end{equation}
其中最后一个映射是自然伴随。
为记号简便，上述映射仍记作 $f_!(u)$。则 $(D,f_!(u))$ 是从 $(Y_1,(f_1)_!\mF_1)$ 到 $(Y_2,(f_2)_!\mF_2)$ 的上同调对应。称它为\emph{上同调对应的前推}。

特别地，若 $(Y_1\xleftarrow{d_1} D\xrightarrow{d_2} Y_2)$ 为 $(S=S=S)$（$S$ 是某基概形），且 $c_1$ 紧合（故左方块可基变换），则得到层的同态
\[f_!(u): (f_1)_!\mF_1\to (f_2)_!\mF_2.\]
更具体地，若 $S=\Spec k$ 代数闭且 $c_1$ 紧合，则得到上同调上的诱导映射 $\on{H}_c(u): \on{H}_c^*(X_1,\mF_1)\to \on{H}_c^*(X_2,\mF_2)$。

\begin{rmk}
若 $f_1$ 分离且 $f$ 可由完全紧合代数空间表示（故左方块可基变换），则 $f_!(u)$ 分解为自然映射 $(f_1)_!\mF_1\to (f_1)_*\mF_1$ 与
\[d_1^*(f_1)_*\mF_1\to f_*c_1^*\mF_1\to f_!c_2^!\mF_2\to d_2^!(f_2)_!\mF_2\]
的复合。后一映射即 SGA 5 中考虑的上同调对应的前推。不过在我们处理的大多数情形中，$d_1$ 分离而 $c_1$ 可由完全紧合代数空间表示。
\end{rmk}

上同调对应的前推与上同调对应的复合相容。
\begin{lem}
\label{AL:pushforward compatible with composition}
假设有交换图
\[
\begin{CD}
X_1@<c_1<< C@>c_2>> X_2@<c'_1<< C'@>c'_2>> X_3\\
@Vf_1VV@VVfV@VVf_2V@VVf'V@VVf_3V\\
Y_1@<d_1<<D@>d_2>>Y_2@<d'_1<<D'@>d'_2>>Y_3
\end{CD}
\]
其中第一与第三个方块可基变换，$c_1,  c'_1$ 紧合，$d_1,d'_1$ 分离；设 $u: c_1^*\calF_1 \to c_2^!\calF_2$ 与 $v: c'^*_1\calF_2 \to c'^!_2 \calF_3$ 分别是从 $(X_1, \calF_1)$ 到 $(X_2, \calF_2)$ 与从 $(X_2, \calF_2)$ 到 $(X_3, \calF_3)$ 的上同调对应。
记 $\widetilde C := C \times_{X_2} C'$，$\widetilde D: = D \times_{Y_2} D'$，并设 $\tilde f: \widetilde C \to \widetilde D$ 是自然诱导的态射。
则以下命题成立。
\begin{enumerate}
\item 下图中的右方块
\[\begin{CD}
X_1@<c_1<<C@<p_C<<\tilde C \\
@Vf_1VV@VVfV@VV\tilde fV\\
Y_1@<d_1<<D@<P_D<<\tilde D
\end{CD}\]
可基变换。注意这保证了 $\tilde f_!(v \circ u)$ 按引理 \ref{L: criterion basechangeable} (1) 有定义。
\item 从 $(Y_1, (f_1)_!\calF_1)$ 到 $(Y_3, (f_3)_!\calF_3)$、支撑在 $\widetilde D$ 上的上同调对应满足等式
\[
\tilde f_!(v\circ u) = f'_!(v) \circ f_!(u).
\]
特别地，若 $c_1$ 与 $c'_1$ 可由（完全）紧合代数空间表示，则
$$\on{H}_c(v\circ u) = \on{H}_c(v) \circ \on{H}_c(u):\on{H}^*_c(X_1\otimes\bar k, \calF_1) \to \on{H}^*_c(X_3\otimes \bar k, \calF_3).$$
\end{enumerate}
\end{lem}

下面假设有交换图
\[\begin{CD}
Y_1@<d_1<< D@>d_2>> Y_2\\
@Vf_1VV@VVfV@VVf_2V\\
X_1@<c_1<<C@>c_2>>X_2
\end{CD}\]
给定 $u:c_1^*\mF_1\to c_2^!\mF_2$，使 $(C,u)$ 是从 $(X_1,\mF_1)$ 到 $(X_2,\mF_2)$ 的上同调对应，我们希望定义 $(C,u)$ 的「拉回」。
至少有三种情形可以做到。
\begin{enumerate}
\item[(1)] 假设 $f$ 是（上同调）光滑、等维的，相对维数为 $d_f$，并选定迹映射 $\Tr_f: R^{2d_f}f_!E(d)\to E$（见 \cite[定义 A.1.18]{XZ}）。
则得到从 $(Y_1, f_1^*\calF_1)$ 到 $(Y_2, f_2^!\calF_2\langle-2d_f\rangle)$ 的上同调对应 $(D, f^*(u))$，定义如下：
\begin{equation}
\label{ASS:smooth pullback correspondence} 
f^*(u)\colon d_1^*f_1^*\mF_1\cong f^*c_1^*\mF_1\xrightarrow{f^*u} f^*c_2^!\mF_2\cong f^!c_2^!\mF_2\langle-2d_f\rangle \cong d_2^!f_2^!\mF_2\langle-2d_f\rangle.
\end{equation}
这里及下文，$\langle d\rangle$ 按惯例表示平移与扭转 $[d](d/2)$。
称此为上同调对应的\emph{光滑拉回}。注意它依赖于 $\Tr_f$ 的选取。

\item[(2)] 若上图的左内方块是笛卡儿的，则有
\begin{equation}
\label{ASS:smooth pullback correspondence2} 
f^!(u)\colon d_1^*f_1^!\mF_1\xrightarrow{\on{BC}^{*!}} f^!c_1^*\mF_1\xrightarrow{f^!u}f^!c_2^!\mF_2= d_2^!f_2^!\mF_2.
\end{equation}
若进一步 $f_1$ 上同调光滑，可选取迹映射 $\Tr_{f_1}$ 把 $f^*_1\simeq f^!_1\langle-2d_{f_1}\rangle$ 等同起来，并取 $\Tr_f$ 为 $\Tr_{f_1}$ 的基变换。此时 \eqref{ASS:smooth pullback correspondence2} 与 \eqref{ASS:smooth pullback correspondence} 一致。

\item[(3)] 同理，若上图的右内方块是笛卡儿的，则可定义
\begin{equation}
\label{ASS:smooth pullback correspondence3} 
f^*(u)\colon d_1^*f_1^*\mF_1= f^*c_1^*\mF_1\xrightarrow{f^*u} f^*c_2^!\mF_2\xrightarrow{\on{BC}^{*!}} d_2^! f_2^*\mF_2.
\end{equation}
若 $f$ 上同调光滑，按上面选取 $\Tr_f$，并取 $\Tr_{f_2}$ 为 $\Tr_f$ 的基变换以等同 $f_2^*\cong f_2^!\langle-2d_{f}\rangle$。此时 \eqref{ASS:smooth pullback correspondence3} 与 \eqref{ASS:smooth pullback correspondence} 一致。
\end{enumerate}

与前推的情形类似，在适当条件下，上同调对应的拉回也与上同调对应的复合相容。
\begin{lem}
\label{AL:pullback compatible with composition}
假设有交换图
\[
\begin{CD}
Y_1@<d_1<<D@>d_2>>Y_2@<d'_1<<D'@>d'_2>>Y_3\\
@Vf_1VV@VVfV@VVf_2V@VVf'V@VVf_3V\\
X_1@<c_1<< C@>c_2>> X_2@<c'_1<< C'@>c'_2>> X_3
\end{CD}
\]
其中 $f$ 上同调光滑、等维，相对维数为 $d_f$。假设第三个内方块是笛卡儿的。
记 $\widetilde C := C \times_{X_2} C'$，$\widetilde D: = D \times_{Y_2} D'$，并设 $\tilde f:\tilde D\to\tilde C$ 是自然映射。则 $\tilde f$ 是 $f$ 沿 $\tilde C\to C$ 的基变换，故上同调光滑。选取 $\Tr_f$，并取 $\Tr_{\tilde f}$ 为 $\Tr_f$ 的基变换。
设 $u: c_1^*\calF_1 \to c_2^!\calF_2$ 与 $v: c'^*_1\calF_2 \to c'^!_2 \calF_3$ 分别是从 $(X_1, \calF_1)$ 到 $(X_2, \calF_2)$ 与从 $(X_2, \calF_2)$ 到 $(X_3, \calF_3)$ 的上同调对应。则上同调对应 $\tilde f^*(v\circ u): (Y_1, f_1^*\calF_1) \to (Y_3, f_3^!\calF_3\langle 2d_f\rangle )$ 等于
\[
(Y_1, f_1^*\calF_1) \xrightarrow{f^*(u)} \big(Y_2, f_2^!\calF_2\langle 2d_f\rangle \big) \xrightarrow{f'^!(v)} \big(Y_3, f_3^!\calF_3\langle 2d_f\rangle \big).
\]
\end{lem}

\subsection{一个由对应构成的范畴}\label{SS: cat of corr}
设 $(\overleftarrow{c},\overrightarrow{c}): C\rightrightarrows X$ 是空间中的群胚。特别地，有下图
\begin{equation*}
\label{E: diagram}
\xymatrix{
&& C\times_XC\ar_{\overleftarrow{p}}[ld]\ar^{\overrightarrow{p}}[dr]\ar^m[d]& &\\
&C\ar_{\overleftarrow{c}}[ld]\ar_{\overrightarrow{c}}[rd]& C\ar_{\overleftarrow{c}}[lld]\ar^{\overrightarrow{c}}[rrd]&C\ar^{\overleftarrow{c}}[ld]\ar^{\overrightarrow{c}}[rd]&\\
X&&X&&X
}
\end{equation*}
其中 $\overleftarrow{p},\overrightarrow{p}: C\times_XC\to C$ 是第一、第二因子的投影，$m: C\times_XC\to C$ 是乘法映射。另外有映射 $\delta: X\to C$，使 $(\overleftarrow{c},\overrightarrow{c})\circ \delta:X\to X\times X$ 是对角映射。

现假设 $m$ 紧合且可由代数空间表示，$\delta$ 是闭浸入。
如下定义 $E$-线性范畴 $\on{D}^C(X)$：对象是 $X$ 上的层；态射为
\begin{equation}
\label{E:HomDC}
\Hom_{\on{D}^C(X)}(\mF_1,\mF_2)=\Hom_{\on{D}(C)}(\overleftarrow{c}^*\mF_1, \overrightarrow{c}^!\mF_2),
\end{equation}
即支撑在 $C$ 上、从 $\mF_1$ 到 $\mF_2$ 的上同调对应。

上一小节回顾的上同调对应形式化允许定义态射的合成律
\begin{eqnarray*}
&&\Hom_{\on{D}^C(X)}(\mF_1,\mF_2)\otimes\Hom_{\on{D}^C(X)}(\mF_2,\mF_3)\\
&=&\on{Corr}_C((X,\mF_1),(X,\mF_2))\otimes \on{Corr}_C((X,\mF_2),(X,\mF_3))\\
&\to& \on{Corr}_{C\times_XC}((X,\mF_1),(X,\mF_3)) \\
&\xrightarrow{m_!}& \on{Corr}_C((X,\mF_1),(X,\mF_3))\\
&=&\Hom_{\on{D}^C(X)}(\mF_1,\mF_3).
\end{eqnarray*}
它含一个全子范畴 $\on{P}^C(X)$，由那些在 $X$ 上反常的 $\mF\in\on{D}^C(X)$ 组成。

考虑上述范畴的动机如下。
设 $f:X\to Y$ 是可表示的紧合态射，$C=X\times_YX$，$\overleftarrow{c}$ 与 $\overrightarrow{c}$ 是两投影。
按定义，
\begin{lem}
\label{L: fully faithful}
函子 $f_!: \on{D}(X)\to\on{D}(Y)$ 分解为
\[\on{D}(X)\to \on{D}^C(X)\to \on{D}(Y)\]
且第二个箭头全忠实。
\end{lem}
\begin{proof}
第一个断言是显然的。
第二个断言来自典范同构（由紧合基变换）
\[\Hom(\overleftarrow{c}^*\mF_1,\overrightarrow{c}^!\mF_2)=\Hom(f_!\mF_1,f_!\mF_2),\]
它与合成相容。
\end{proof}

\begin{rmk}
\label{R: DC enh}
(1) 作上述定义的一个理由在于：有些情形下即使叠 $Y$（从而 $\on{D}(Y)$）性态很差，范畴 $\on{D}^C(X)$ 仍可有恰当的定义。

(2) 上面作为 $E$-线性范畴引入的 $\on{D}^C(X)$ 已足够本文的应用。不过我们说明如何把它定义为 $\infty$-范畴，至少在 $\overleftarrow{c}$ 与 $\overrightarrow{c}$ 都是有限或代数闭域 $k$ 上分离（完全）有限型概形的（完全）紧合态射的情形。记 $\Shv(-)$ 为 $\on{D}(-)$ 的自然 $\infty$-范畴加强。若 $X$ 在 $\ell$-上同调维数有限的域上分离且（完全）有限型，则当 $E=\bF_\ell,\bZ_\ell, \bQ_\ell$ 等时，$\Shv(X,E)$ 由紧对象生成，紧对象即可构造复形。

注意 $\on{D}^C(X)$ 是 $E$-线性范畴 $2$-范畴中下图之余极限的单子式描述
\[\on{D}(C\times_XC)\ \substack{\longrightarrow\\[-1em] \longrightarrow \\[-1em] \longrightarrow}\ \on{D}(C)\ \substack{\longrightarrow\\[-1em] \longrightarrow} \ \on{D}(X),\]
其中函子由紧合前推给出。
于是显然可以把 $\infty$-范畴 $\Shv^C(X)$ 定义为
\[\cdots \substack{\longrightarrow\\[-1em] \longrightarrow \\[-1em] \longrightarrow\\[-1em] \longrightarrow}\ \Shv(C\times_XC)\ \substack{\longrightarrow\\[-1em] \longrightarrow \\[-1em] \longrightarrow}\ \Shv(C)\ \substack{\longrightarrow\\[-1em] \longrightarrow} \ \Shv(X)\]
关于紧合前推的余极限。
由 Barr–Beck–Lurie 定理，$\Shv^C(X)$ 有关于 $\Shv(X)$ 的单子式描述。特别地，其同伦范畴的 Hom 空间可用上同调对应描述：设 $\mF_1,\mF_2\in \Shv(X)$，经自然函子 $\Shv(X)\to \Shv^C(X)$ 看作 $\Shv^C(X)$ 的对象。则
$$\Hom_{\on{h}\Shv^C(X)}(\mF_1,\mF_2)=R\Hom_{\on{D}(C)}(\overleftarrow{c}^*\mF_1,\overrightarrow{c}^!\mF_2).$$
若 $C=X\times_YX$，$f:X\to Y$ 是某紧合满射，则 $f_!:\Shv(X)\to \Shv(Y)$ 诱导 $\infty$-范畴的等价 $\Shv^C(X)\simeq \Shv(Y)$。
\end{rmk}

\begin{thebibliography}{9999}

\bibitem[BL]{BL}A. Beauville and Y. Laszlo, \emph{Conformal blocks and generalized theta functions}, Comm. Math. Phys. 164 (1994), no. 2, 385-419.

\bibitem[BD]{BD}A. Beilinson, V. Drinfeld, \emph{Quantization of Hitchin's integrable system and Hecke eigensheaves}, Preprint, available at \url{www.math.uchicago.edu/~mitya/langlands}.

\bibitem[BN]{BN}D. Ben-Zvi, D. Nadler, \emph{The character theory of a complex group}, arXiv:0904.1247.

\bibitem[BFO]{BFO}R. Bezrukavnikov, M. Finkelberg, V. Ostrik, \emph{Character $D$-modules via Drinfeld center of Harish-Chandra bimodules}, Invent. Math. 188 (2012), no. 3, 589-620.

\bibitem[BKV]{BKV}R. Bezrukavnikov, D. Kazhdan, Y. Varshavsky, \emph{A categorical approach to the stable center conjecture}, Ast\'erisque 369 (2015), 27-97.

\bibitem[BS]{BS}B. Bhatt, P. Scholze, \emph{Projectivity of the Witt vector affine Grassmannian}, Invent. Math. 209 (2017), no. 2, 329-423.

\bibitem[BG]{BG}K. Buzzard, T. Gee, \emph{The conjectural connections between automorphic representations and Galois representations}, Proceedings of the LMS Durham Symposium 2011.

\bibitem[DM]{DM}P. Deligne, J. S. Milne, \emph{Tannakian categories}, Hodge cycles, motives, and Shimura varieties, LNM 900 (1982), 101?228.

\bibitem[Dr]{Dr}V. G. Drinfeld, \emph{Two-dimensional $\ell$-adic representations of the fundamental group of a curve over a finite field and automorphic forms on $\GL(2)$}, Amer. J. Math. 105 (1983), 85-114.

\bibitem[EG]{EG}D. Edidin, W. Graham, \emph{Equivariant Intersection Theory}, Invent. Math., 131 (1998) 595-634.

\bibitem[Fa]{Fa}G. Faltings, \emph{Algebraic loop groups and moduli spaces of bundles}, J. Eur. Math. Soc. 5 (2003), no. 1, 41-68.

\bibitem[FKK]{FKK}B. Fontaine, J. Kamnitzer, and G. Kuperberg, \emph{Buildings, spiders, and geometric Satake}, Compos. Math. 149 (2013), no. 11, 1871-1912.

\bibitem[Ga]{Ga}D. Gaitsgory, \emph{From geometric to function-theoretic Langlands (or how to invent shtukas)}, arXiv:1606.09608.

\bibitem[GL17]{GL}A. Genestier, V. Lafforgue, \emph{Chtoucas restreints pour les groupes r\'eductifs et param\'etrisation de Langlands locale}, arXiv:1709.00978.

\bibitem[Gi95]{Gi1}V. Ginzburg, \emph{Perverse sheaves on a loop group and Langlands' duality}, arXiv:math/9511007.

\bibitem[Gros]{Gr}B. Gross, \emph{On the Satake isomorphism}, Galois representations in arithmetic algebraic geometry (Durham, 1996), London Math. Soc. Lecture Note Ser., vol. 254, Cambridge Univ. Press, Cambridge, 1998, pp. 223-237.

\bibitem[GI]{Illusie}
A. Grothendieck, r\'edig\'e par L. Illusie, \emph{Formule de Lefschetz}, expos\'e III, \textbf{SGA} 5, Springer LNM 589 (1977), 73-137.

\bibitem[He]{He} D. Helm, \emph{Towards a geometric Jacquet-Langlands correspondence for unitary Shimura varieties}, Duke Math. J. 155 (2010), 483-518.

\bibitem[HTX]{HTX} D. Helm, Y. Tian, and L. Xiao, \emph{Tate cycles on some unitary Shimura varieties mod $p$}, Alegbra and Number Theory 11 (2017), 2213-2288.

\bibitem[Ja]{Ja}U. Jannsen, \emph{Motives, numerical equivalence, and semi-simplicity}, Invent. math. 107 (1992), no. 3, 447-452.

\bibitem[Ki]{Ki} M. Kisin, \emph{Integral models for Shimura varieties of abelian type}, J. Amer. Math. Soc. 23 (2010), no. 4, 967-1012.

\bibitem[Kl]{Kl}S.L. Kleiman, \emph{Motives}, algebraic geometry, Oslo 1970. Oort, F. (ed.), 53-82. Groningen, 1972.

\bibitem[La]{La} V. Lafforgue, \emph{Chtoucas pour les groupes r\'eductifs et param\'etrisation de Langlands globale}, arXiv:1209.5352.

\bibitem[Lu1]{Lu}G. Lusztig, \emph{Singularities, character formulas, and a $q$-analog of weight multiplicities}, In Analysis and topology on singular spaces, II, III (Luminy, 1981), 208--229, Ast\'erisque, 101-102, Soc. Math. France, Paris, 1983.

\bibitem[Lu2]{Lu2}G. Lusztig, \emph{A bar operator for involutions in a Coxeter group}, Bull. Inst. Math. Acad. Sin. (N.S.) 7 (2012), 355-404.

\bibitem[Lu3]{Lu3}G .Lusztig, \emph{Truncated convolution of character sheaves}, Bull. Inst. Math. Acad. Sin. (N.S.) 10 (2015), 1-72.

\bibitem[Lu4]{Lu4}G. Lusztig, \emph{Unipotent representations as a categorical centre}, Represent. Th. 19 (2015), 211-235.

\bibitem[Lu5]{Lu5}G. Lusztig, \emph{Non-unipotent character sheaves as a categorical centre}, Bull. Inst. Math. Acad. Sin. (N.S.) 11 (2016), 603-731.

\bibitem[Lu6]{Lu6}G. Lusztig, \emph{Non-unipotent representations and categorical centres}, Bull. Inst. Math. Acad. Sin. (N.S.) 12 (2017), 205-296.

\bibitem[LV]{LV}G. Lusztig, D. Vogan, \emph{Hecke algebras and involutions in Weyl groups}, Bull. Inst. Math. Acad. Sinica (N.S.) 7(2012), 323-354.

\bibitem[LY]{LY}G. Lusztig, Z. Yun, \emph{A $(-q)$-analog of weight multiplicities}, Journal of the Ramanujan Math. Soc., 28A (Special Issue-2013) 311-340.

\bibitem[MV]{MV}
I. Mirkovi\'c, K. Vilonen,
\emph{Geometric Langlands duality and representations of algebraic groups over commutative rings},
Ann. Math. 166 (2007), no. 1, 95-143.

\bibitem[Ng]{Ng}B.C. Ng\^o, \emph{Le lemme fondamental pour les alg\`ebres de Lie}, Publ. Math. Inst. Hautes tudes Sci. No. 111 (2010), 1-169.

\bibitem[Os]{O}V. Ostrik, \emph{Pivotal fusion categories of rank $3$}, Mosc. Math. J.15 (2015) 373-396.

\bibitem[PWZ]{PWZ}R. Pink, T. Wedhorn, and P. Ziegler, \emph{$F$-zips with additional structure}, Pacific J. Math. 274 (2015), 183-236.

\bibitem[Ri]{Ri}T. Richarz, \emph{Affine Grassmannians and geometric Satake equivalences}, Int Math Res Notices, Volume 2016, Issue 12, pp. 3678-3716.

\bibitem[Sc]{S}P. Scholze, \emph{$p$-adic geometry}, arXiv:1712.03708.

\bibitem[SW]{SW}P. Scholze, J. Weinstein, \emph{$p$-adic geometry}, Berkeley lecture notes,
available at \url{http://www.math.uni-bonn.de/people/scholze/Berkeley.pdf}.

\bibitem[TX]{TX}
Y. Tian and L. Xiao, \emph{Tate cycles on quaternionic Shimura varieties over finite fields}, arXiv:1410.2321.

\bibitem[Va]{Vasiu}A. Vasiu, \emph{Good Reductions of Shimura Varieties of Hodge Type in Arbitrary Unramified Mixed Characteristic, Part I, II}, arXiv:0707.1668, 0712.1572.

\bibitem[VW]{VW}
I. Vollaard and T. Wedhorn,
The supersingular locus of the Shimura variety of $GU(1,n-1)$, II, \emph{
Invent. Math.} \textbf{184} (2011), 591--627.

\bibitem[XZ1]{XZ}L. Xiao, X. Zhu, \emph{Cycles on Shimura varieties via geometric Satake}, arXiv:1707.05700.

\bibitem[XZ2]{XZ2}L. Xiao, X. Zhu, \emph{On vector-valued twisted conjugate invariant functions on a group}, arXiv:1802.05299.

\bibitem[XZ3]{XZ3}L. Xiao, X. Zhu, \emph{Cycles on Shimura varieties via geometric Satake: Examples}, preprint.

\bibitem[Zha]{Zha} C. Zhang, \emph{Ekedahl-Oort strata for good reductions of Shimura varieties of Hodge type}, arXiv:1312.4869.

\bibitem[Zh1]{Zh0}X. Zhu, \emph{The Geometric Satake Correspondence for Ramified Groups}, Ann. Sci. \'Ec. Norm. Sup\'er., 48 (2015), 409-451.

\bibitem[Zh2]{Zh1}
X. Zhu, \emph{Affine Grassmannians and the geometric Satake in mixed characteristic}, Ann. Math. 185 (2017), No. 2, 403-492.

\bibitem[Zh3]{Zh2}
X. Zhu, \emph{An introduction to affine Grassmannians and the geometric Satake equivalence}, IAS/Park City Mathematics Series, 24. (2017) 59-154.

\end{thebibliography}


\end{document}
